%% GEB LaTeX v2 — Part 05 / 7
%% 自动生成  split_v2.py
%% 独立编译（从 GEB_LaTeX_v2/ 目录）：
%%   xelatex -output-directory=split part05.tex

%% GEB 精排模板 ── XeLaTeX
\PassOptionsToPackage{unicode}{hyperref}
\PassOptionsToPackage{hyphens}{url}

%% ─────────────────────────────────────────────
%%  文档类：memoir（精排书籍首选）
%% ─────────────────────────────────────────────
\documentclass[
  12pt,
  twoside,
  openright,
  oldfontcommands,
]{memoir}

%% memoir 纸张尺寸（必须在字体等 usepackage 之前）
\setstocksize{260mm}{185mm}   % 大16开
\settrimmedsize{\stockheight}{\stockwidth}{*}
\settrims{0pt}{0pt}

%% ─────────────────────────────────────────────
%%  字体
%% ─────────────────────────────────────────────
\usepackage{fontspec}
\usepackage{xeCJK}

% ※（U+203B）Georgia 无此字符，强制用 CJK 字体渲染
\xeCJKDeclareCharClass{CJK}{"203B}

% geb.ttf：EPUB 自定义字体，含 Tai Viet 占位字形（U+AA9E–U+AAAF）
% 以及 U+31CF、U+39DF 等特殊字形
\newfontfamily\gebfont{geb.ttf}[Path=./fonts/]

% 中文正文：思源宋体（Noto Serif CJK SC）
\setCJKmainfont[
  BoldFont    = {Noto Serif CJK SC SemiBold},
  ItalicFont  = {Kaiti SC},          % 楷体作斜体
]{Noto Serif CJK SC}

% 中文无衬线（标题等）
\setCJKsansfont[
  BoldFont = {Heiti SC},
]{Heiti SC}

% 中文等宽
\setCJKmonofont{PingFang SC}

% 西文正文：Palatino（等高数字，基线与思源宋体对齐更佳）
\setmainfont[
  BoldFont       = {Palatino Bold},
  ItalicFont     = {Palatino Italic},
  BoldItalicFont = {Palatino Bold Italic},
]{Palatino}

% 西文无衬线
\setsansfont{Helvetica Neue}

% 数学字体
\usepackage{amsmath,amssymb}
\usepackage{unicode-math}

%% ─────────────────────────────────────────────
%%  版面：大16开（185×260mm）页边距（memoir 原生布局）
%% ─────────────────────────────────────────────
\setlrmarginsandblock{25mm}{20mm}{*}
\setulmarginsandblock{22mm}{25mm}{*}
\setheaderspaces{*}{7mm}{*}
\setheadfoot{\baselineskip}{12mm}
\checkandfixthelayout

%% ─────────────────────────────────────────────
%%  颜色
%% ─────────────────────────────────────────────
\usepackage{xcolor}
\definecolor{accent}{HTML}{1A3A5C}     % 深蓝（章号、页眉线）
\definecolor{accentlight}{HTML}{4A7DAA}% 中蓝（节标题）
\definecolor{rulegray}{HTML}{BBBBBB}   % 分隔线

%% ─────────────────────────────────────────────
%%  章节样式
%% ─────────────────────────────────────────────
\usepackage{titlesec}

% 章：大号宋体居中，上方装饰线，下方细实线
\titleformat{\chapter}[display]
  {\normalfont\centering}
  {\color{accent}\small\CJKfamily{}\sffamily 第\ \thechapter\ 章}
  {0pt}
  {\color{accent}\LARGE\bfseries\CJKfamily{}\rmfamily}
  [\vspace{4pt}\textcolor{rulegray}{\titlerule[0.4pt]}]

\titlespacing*{\chapter}{0pt}{30pt}{40pt}

% 节：宋体加粗左对齐，左侧蓝色竖线装饰
\titleformat{\section}
  {\normalfont\large\bfseries\color{accentlight}}
  {}
  {0pt}
  {\colorbox{white}{\raisebox{0pt}{\rule[-1pt]{3pt}{14pt}}\hspace{6pt}}}
  []

\titlespacing*{\section}{0pt}{18pt}{8pt}

% 小节
\titleformat{\subsection}
  {\normalfont\normalsize\bfseries}
  {}
  {0pt}
  {}

%% ─────────────────────────────────────────────
%%  页眉页脚（memoir 内建）
%% ─────────────────────────────────────────────
\makepagestyle{geb}
\makeevenhead{geb}{\small\thepage}{}{\small\itshape\leftmark}
\makeoddhead {geb}{\small\itshape\rightmark}{}{\small\thepage}
\makeheadrule {geb}{\textwidth}{0.4pt}
\pagestyle{geb}
\renewcommand{\chaptermark}[1]{\markboth{第\ \thechapter\ 章\enspace #1}{第\ \thechapter\ 章\enspace #1}}
\renewcommand{\sectionmark}[1]{\markright{\thesection\enspace #1}}

%% ─────────────────────────────────────────────
%%  中文标准名称
%% ─────────────────────────────────────────────
\renewcommand{\contentsname}{目录}
\renewcommand{\figurename}{图}
\renewcommand{\tablename}{表}
\renewcommand{\appendixname}{附录}
\renewcommand{\bibname}{参考文献}
\renewcommand{\indexname}{索引}

%% ─────────────────────────────────────────────
%%  段落间距与行距（memoir 内建命令）
%% ─────────────────────────────────────────────
\setSpacing{1.3}
\setlength{\parindent}{2em}
\setlength{\parskip}{3pt plus 1pt minus 1pt}
\raggedbottom

%% ─────────────────────────────────────────────
%%  图片与浮动体
%% ─────────────────────────────────────────────
\usepackage{graphicx}
\usepackage{float}
\usepackage{caption}
\captionsetup{
  font=small,
  labelfont={bf,color=accent},
  labelsep=quad,
  width=0.88\linewidth,
  justification=centering,
  singlelinecheck=false,
}

%% pandocbounded：memoir 兼容版，限制图片最大宽度和高度
\newcommand*\pandocbounded[1]{%
  \begingroup
  \setkeys{Gin}{width=\linewidth,height=0.8\textheight,keepaspectratio}%
  #1%
  \endgroup
}
\def\fps@figure{htbp}

%% ─────────────────────────────────────────────
%%  引用块样式
%% ─────────────────────────────────────────────
\usepackage{mdframed}
\mdfdefinestyle{gebquote}{
  leftline=true,
  rightline=false,
  topline=false,
  bottomline=false,
  linewidth=2pt,
  linecolor=accentlight,
  backgroundcolor=black!3,
  innerleftmargin=12pt,
  innerrightmargin=8pt,
  innertopmargin=6pt,
  innerbottommargin=6pt,
}
\renewenvironment{quote}
  {\begin{mdframed}[style=gebquote]\small}
  {\end{mdframed}}

%% ─────────────────────────────────────────────
%%  表格
%% ─────────────────────────────────────────────
\usepackage{longtable,booktabs,array,multirow,calc}
\usepackage{etoolbox}
\makeatletter
\patchcmd\longtable{\par}{\if@noskipsec\mbox{}\fi\par}{}{}
\makeatother
\IfFileExists{footnotehyper.sty}{\usepackage{footnotehyper}}{\usepackage{footnote}}
\makesavenoteenv{longtable}

% 美化 booktabs 规则
\setlength{\aboverulesep}{0pt}
\setlength{\belowrulesep}{0pt}
\setlength{\heavyrulewidth}{0.8pt}
\setlength{\lightrulewidth}{0.4pt}

%% ─────────────────────────────────────────────
%%  脚注
%% ─────────────────────────────────────────────
\usepackage[
  bottom,
  hang,
  flushmargin,
]{footmisc}
\setlength{\footnotemargin}{1.5em}
\renewcommand{\footnoterule}{%
  \vspace{4pt}\noindent\textcolor{rulegray}{\rule{0.3\linewidth}{0.4pt}}\vspace{4pt}%
}

%% ─────────────────────────────────────────────
%%  目录样式
%% ─────────────────────────────────────────────
\usepackage{titletoc}

%% memoir 的 chapternumberline/partnumberline 需映射到 \numberline
%% 才能被 titletoc 识别为「有编号」条目
\makeatletter
\renewcommand{\chapternumberline}[1]{\numberline{#1}}
\renewcommand{\partnumberline}[1]{\numberline{#1}}
\makeatother

\titlecontents{chapter}[0em]
  {\addvspace{6pt}\bfseries\color{accent}}
  {\makebox[3em][l]{\thecontentslabel}}
  {}
  {\hfill\contentspage}

\titlecontents{section}[3em]
  {\addvspace{2pt}\small}
  {}
  {}
  {\titlerule*[0.5em]{.}\contentspage}

%% ─────────────────────────────────────────────
%%  超链接
%% ─────────────────────────────────────────────
\usepackage{bookmark}
\IfFileExists{xurl.sty}{\usepackage{xurl}}{}
\urlstyle{same}
\usepackage{hyperref}
\hypersetup{
  pdftitle={哥德尔、艾舍尔、巴赫——集异璧之大成},
  pdfauthor={〔美〕侯世达},
  pdflang={zh-CN},
  colorlinks=true,
  linkcolor={accent},
  urlcolor={accentlight},
  citecolor={accent},
  pdfcreator={XeLaTeX + GEB Template},
}

%% ─────────────────────────────────────────────
%%  杂项
%% ─────────────────────────────────────────────
\usepackage{microtype}
\setlength{\emergencystretch}{3em}
\providecommand{\tightlist}{%
  \setlength{\itemsep}{0pt}\setlength{\parskip}{0pt}}

% 对话角色加粗
\newcommand{\speaker}[1]{\textbf{#1}}

%% ─────────────────────────────────────────────
%%  GEB v2 字体命令
%% ─────────────────────────────────────────────
\newCJKfontfamily\kaitifamily[BoldFont={Kaiti SC Bold}]{Kaiti SC}
\newcommand{\kaititext}[1]{{\kaitifamily #1}}
\newcommand{\kaitismall}[1]{{\kaitifamily\small #1}}
\newcommand{\kaitibold}[1]{{\kaitifamily\bfseries #1}}
\newcommand{\heititext}[1]{{\sffamily\bfseries #1}}
\newcommand{\heitismall}[1]{{\sffamily\bfseries\small #1}}
\newcommand{\dialogname}[1]{{\sffamily\bfseries #1}}
\newcommand{\westerntext}[1]{{\rmfamily #1}}
\newcommand{\westernital}[1]{{\rmfamily\itshape #1}}

%% ─────────────────────────────────────────────
%%  GEB v2 排版环境（memoir 内建 hangparas，不需要 hanging 宏包）
%% ─────────────────────────────────────────────

% 对话悬挂缩进行
\newenvironment{hangdialog}{%
  \begin{hangparas}{2em}{1}%
}{%
  \par\end{hangparas}%
}

% 楷体舞台指示（居中小楷）
\newenvironment{stagedir}{%
  \begin{center}\kaitifamily\small%
}{%
  \end{center}%
}

% 对话引导语（楷体斜体缩进）
\newenvironment{dialogguide}{%
  \begin{adjustwidth}{1em}{0em}\kaitifamily\itshape\small%
}{%
  \end{adjustwidth}%
}

% 仿宋/楷体引用块（带左竖线）
\newenvironment{fsquote}{%
  \begin{mdframed}[style=gebquote]\kaitifamily\small%
}{%
  \end{mdframed}%
}

% 楷体加粗引用框
\newenvironment{ktquote}{%
  \begin{mdframed}[style=gebquote]\kaitifamily\small%
}{%
  \end{mdframed}%
}

% 禅宗公案块
\newenvironment{zenkoan}{%
  \begin{mdframed}[style=gebquote]\kaitifamily\small%
}{%
  \end{mdframed}%
}

% 禅宗公案诗
\newenvironment{zenkoanpoem}{%
  \begin{center}\kaitifamily\small%
}{%
  \end{center}%
}

% 图片副说明（小楷居中，灰色）
\newenvironment{imgsub}{%
  \begin{center}\kaitifamily\small\color{rulegray}%
}{%
  \end{center}%
}

% 对话诗歌
\newenvironment{dialogpoem}{%
  \begin{adjustwidth}{3em}{0em}\kaitifamily%
}{%
  \end{adjustwidth}%
}

% 章节前置数字大标题（如"第一章"）
\newenvironment{chappreheader}{%
  \begin{center}\large\sffamily\bfseries\color{accent}%
}{%
  \end{center}\vspace{-10pt}%
}

% 部分前置大标题
\newenvironment{partpreheader}{%
  \begin{center}\Large\sffamily\bfseries\color{accent}%
}{%
  \end{center}\vspace{-10pt}%
}

%% ─────────────────────────────────────────────
%%  扉页（自定义，替换 maketitle）
%% ─────────────────────────────────────────────
\newcommand{\gebmaketitle}{%
  \begin{titlingpage}
    \vspace*{6\baselineskip}
    \begin{center}
      {\fontsize{28}{34}\selectfont\bfseries\color{accent}
        哥德尔、艾舍尔、巴赫}\\[12pt]
      {\Large\color{rulegray}——集异璧之大成}\\[40pt]
      \textcolor{rulegray}{\rule{0.4\linewidth}{0.4pt}}\\[20pt]
      {\large \textit{Gödel, Escher, Bach: An Eternal Golden Braid}}\\[6pt]
      {\normalsize \color{accentlight} Douglas R.\ Hofstadter}\\[40pt]
      {\normalsize 〔美〕侯世达　著\\[4pt]
        本书翻译组　译}\\[60pt]
      {\small\color{rulegray} 商务印书馆}
    \end{center}
    \vfill
  \end{titlingpage}
}

%% pandoc 在 body 中会注入 \babel@toc，提供空定义避免报错
\makeatletter
\providecommand{\babel@toc}[2]{}
\makeatother

%% memoir + XeLaTeX 下需手动告知 xdvipdfmx 纸张尺寸
\AtBeginDocument{%
  \special{papersize=\the\paperwidth,\the\paperheight}%
}

%% ═══════════════════════════════════════════
%%  正文开始
%% ═══════════════════════════════════════════

\begin{document}
\mainmatter
\setcounter{chapter}{32}
\markboth{}{}

\chapter{咏叹调及其种种变奏}

这几天阿基里斯夜里一直无法入睡。今晚他的朋友乌龟来了，要陪他度过这烦人的时光。

\dialogname{乌龟：}亲爱的阿基，听说失眠弄得你很苦恼，我真难过。我希望我来陪伴你能让你放松下来，摆脱那些使你没法入睡的讨厌刺激。也许我会让你烦得不得了，烦得你最后终于睡着了。要是这样，我也算对你有所帮助了。

\dialogname{阿基里斯：}噢，不，恐怕到现在我已经见过世界上最出色的烦人能手了，他们都想要让我烦得睡过去------可是真不幸，结果都没什么用。你不会比他们更有办法的。别这么想了，龟兄，我邀你来是希望能扯点儿数论里的玩艺儿消遣消遣，这样我至少可以比较愉快地捱过这漫长的时光。你知道吗，我发现谈点儿数论对我这骚乱的精神大有好处。

\dialogname{乌龟：}真是个有趣的怪念头！这让我想起------嗨，是有点像------那个可怜的凯瑟林侯爵。

\dialogname{阿基里斯：}他是谁呀？

\dialogname{乌龟：}哦，他是十八世纪时萨克森的一位侯爵------老实说，是一位后知后觉的侯爵------可就是因为他------哎，要我讲下去吗？这是个挺有趣的故事呢。

\dialogname{阿基里斯：}那还用说，讲下去吧。

\dialogname{乌龟：}有一段时期这位好侯爵患了失眠症，很是痛苦。恰巧有个不错的音乐家住在同一个镇上。于是他委托这位音乐家谱写了一系列变奏曲，由侯爵的宫廷羽管键琴师在他的不眠之夜里为他演奏，以便愉快地消磨这些时间。

\dialogname{阿基里斯：}那个当地的音乐家能胜任吗？

\dialogname{乌龟：}我想他胜任了，因为曲子谱好以后，侯爵非常慷慨地回报了他------侯爵赠给他一只金质高脚杯，里面有一百枚金币。

\dialogname{阿基里斯：}真的吗？我首先得知道他哪儿来的那么个高脚杯，还有那些金币。

\dialogname{乌龟：}可能是他在博物馆里看见的，然后做了些手脚。

\dialogname{阿基里斯：}你是不是暗示说他把它们弄到手后溜走了？

\dialogname{乌龟：}好了好了，我不想那么说，不过\ldots\ldots 那种时候，侯爵是能弄走好多东西的，无论如何，有一点很清楚：那些音乐给侯爵带来了巨大的快乐。因为他不断请求他的羽管键琴师------那时还只是个小伙子，叫哥德堡------为他演奏那三十首变奏曲中的曲子。结果（这多少带点讽刺）那些变奏曲倒和那位年轻人的名字``哥德堡''联系在一起了，而不是跟那位著名的侯爵。

\dialogname{阿斯里斯：}你是说，那位作曲家是巴赫，而那音乐就是所谓的《哥德堡变奏曲》？

\dialogname{乌龟：}正是！实际上，这部作品取名为《咏叹调及其种种变奏》，一共有三十首。你知道巴赫是怎么把这三十首辉煌的变奏曲组织在一起的吗？

\dialogname{阿基里斯：}你告诉我吧。

\dialogname{乌龟：}所有曲子------除了最后一首------都基于一个单一的主题，他称之为``咏叹调''。实际上，把曲子结合在一起的不是一支共同的旋律，而是一个共同的和声基础。旋律可以变化，但背后那个主题是稳定的。只是在最后一支变奏曲里巴赫变得随便起来。那是种``终了后的终了''，包含有与原主题几乎无关的多余的音乐------事实上，那是两支德国民歌。那支变奏曲称为``拼凑曲''。

\dialogname{阿基里斯：}《哥德堡变奏曲》还有什么不同寻常的地方吗？

\dialogname{乌龟：}有。每隔两支变奏曲有一首卡农。第一首卡农中两个轮唱着的声部在同一个音符上加入。在第二首中，轮唱着的声部中的一个加入时比另一个高一个音符。在第三首中，一个声部比另一个声部高两个音符加入。如此类推，直到最后一首卡农加入，恰好相差九个音程。一共十首卡农。而------

\dialogname{阿基里斯：}等一下。我好像记起来在什么地方读到过最近新发现了十四首哥德堡卡农\ldots\ldots{}

\dialogname{乌龟：}是不是那份刊物上还说最近发现十一月一共有四十四天，其中新添的那十四天以前不为人们所知？

\dialogname{阿基里斯：}没错儿。一个名叫沃尔夫的人------一位音乐学家------听说在斯特拉斯堡有一份特殊的《哥德堡变奏曲》副本，就到那里看了看。令人吃惊的是，在纸的背面。作为一种``终了后的终了''，他发现了这十四首新卡农。它们全都基于《哥德堡变奏曲》主题的前八个音符。所以现在大家都知道《哥德堡变奏曲》中的变奏曲一共有四十四首，而不是三十首。

\dialogname{乌龟：}也就是说，在另外的什么音乐学家在什么其他鬼地方再发现另一些变奏曲之前，共有四十四首。尽管这看起来不大可能，但也说不定，即便是几乎肯定不会了，也还是有可能会有另一些曲子被发现，然后又发现一些，然后又发现又发现又发现\ldots\ldots 喔，恐怕这是没个头的！我们大概永远不会知道是否------或何时------我们能得到全部的《哥德堡变奏曲》。

\dialogname{阿基里斯：}这倒是个奇怪的想法。大概，每个人都认为这一最新发现仅只是个偶然现象，而我们现在确实真地得到了所有的《哥德堡变奏曲》。不过还是假定你对吧，什么时候也许又会出现一些新的变奏曲。我们可以抱有这种指望。这样的话，``《哥德堡变奏曲》''这个名字的含义就会稍稍有所变化，不仅要包括已经知道的曲子，还要包括任何其他最终会出现的曲子。总数------称之为``g''吧------肯定是有限的，你同意吧？------可仅仅知道g是有限的并不等于知道g有多大。因此，这一信息无法使我们得知什么时候我们能确定最后一支《哥德堡变奏曲》。

\dialogname{乌龟：}说的一点不错。

\dialogname{阿基里斯：}告诉我------巴赫是什么时候写下这些著名变奏曲的？

\dialogname{乌龟：}这是1742年的事，那时他在莱比锡当``康托尔''{[}Cantor{]}------就是教堂圣咏班的指挥。

\dialogname{阿基里斯：}1742？嗯\ldots\ldots 这个数让人想起什么。

\dialogname{乌龟：}理当如此，因为这恰巧是个很有意思的数，它是两个奇素数的和：1723与19。

\dialogname{阿基里斯：}天哪！多么奇特的事情！那么，碰见具有这种性质的偶数的频率会是多少呢？让我们来看看\ldots\ldots{}

6=3+3

8=3+5

10=3+7=5+5

12=5+7

14=3+11=7+7

16=3+13=5+11

18=5+13=7+11

20=3+17=7+13

22=3+19=5+17=11+11

24=5+19=7+17=11+13

26=3+23=7+19=13+13

28=5+23=11+17

30=7+23=11+19=13+17

哦，瞧------按我的这一小张表，似乎这是个很常见的现象。可我到现在也还没看出这张表里有什么显见的规律。

\dialogname{乌龟：}也许就没有什么能找得出的规律。

\dialogname{阿基里斯：}当然得有！我只不过是不够聪明，现在还看不出来。

\dialogname{乌龟：}你好像非常肯定。

\dialogname{阿基里斯：}在我看来这是毫无疑问的。我想\ldots\ldots 会不会所有的偶数（除去4）都能写成两个奇素数之和？

\dialogname{乌龟：}嗯\ldots\ldots 这问题让人想起什么\ldots\ldots 啊，我知道了！你不是第一个问这个问题的人。嗯，事实上，1742年就有个业余数学家提出了这个问题，是在------

\dialogname{阿基里斯：}你说的是1742年吗？请原谅我打断你，不过我注意到1742恰巧是个很有意思的数，它是两个奇素数的差：1747与5。

\dialogname{乌龟：}天哪！多么奇特的事情！那么，碰见具有这种性质的偶数的频率会是多少呢？

\dialogname{阿基里斯：}我还是别让你分心吧，请接着讲你的故事。

\dialogname{乌龟：}哦，对------我刚才说到，1742年有个业余数学家------他的名字我一下子记不起来了------寄了封信给欧拉。欧拉那时正在波茨坦腓德烈大帝的宫廷里，于是------哎，要我讲下去吗？这可不是个没劲的故事。

\dialogname{阿基里斯：}那还用说，讲下去吧。

\dialogname{乌龟：}很好。在信里，这位业余的数论爱好者向伟大的欧拉提出了一个未证明的猜想：``每个偶数都能表示成两个奇素数之和。''这家伙的名字是什么来着？

\dialogname{阿基里斯：}我有点想起这个故事了，是在一本什么数论书里。他的名字是不是``席勒艾舍尔''？

\dialogname{乌龟：}嗯\ldots\ldots 不太像，看来不是。那名字就在我的舌尖上------是------是------啊对了！是``哥德巴赫''！这家伙是哥德巴赫。

\dialogname{阿基里斯：}我就知道是这类名字。

\dialogname{乌龟：}是的，你的猜测帮我想起了它。真怪，人怎么会偶尔搜寻他的记忆，就像在图书馆里查一本不知书号的书一样\ldots\ldots 我们还是回到1742吧。

\dialogname{阿基里斯：}好的，咱们回到1742。我想问问：欧拉是否证明了哥德巴赫的这个猜想是对的？

\dialogname{乌龟：}也真怪极了，他甚至从未想过这值得花功夫。不过，并非所有的数学家都像他这么不屑。事实上，它引起了很多人的注意，并被称为``哥德巴赫猜想''。

\dialogname{阿基里斯：}有人证明了它吗？

\dialogname{乌龟：}还没有。可已经有了一些值得注意的很接近的结果。比如，1931年俄国数论专家施尼勒曼证明了：任何一个数------偶的或奇的------都能表示成不多于300000个素数之和。

\dialogname{阿基里斯：}这结果真怪。它有什么用处呢？

\dialogname{乌龟：}这就把问题限制在了有穷领域，在施尼勒曼给出证明之前，有可能当你取越来越大的偶数时，就需要越来越多的素数来表示。没准有什么偶数需要一万亿个素数来表示！现在我们知道了不会有这种事------300000个素数（或更少些）的和就足够了。

\dialogname{阿基里斯：}我明白了。

\dialogname{乌龟：}后来，1937年，一个名叫维诺格拉多夫的狡猾家伙------也是个俄国人------设法建立了一个与期望目标接近得多的结果，就是：每个足够大的奇数都可以表示成不多于三个奇素数之和。例如，1937=641+643+653。我们可以把能表示成三个奇素数之和的奇数称为具有``维诺格拉多夫性质''的数。这样，所有足够大的奇数都具有维诺格拉多夫性质。

\dialogname{阿基里斯：}很好------可``足够大''是什么意思？

\dialogname{乌龟：}意思是说，有数量有限的某些奇数也许不具有维诺格拉多夫性质，但会有一个数------称之为``v''吧------大于它的所有奇数都是具有维诺格拉多夫性质的。不过，维诺格拉多夫无法确定v有多大。所以某种意义上说，v就像g，那个有限的、但不知多大的哥德堡变奏曲数目。仅仅知道v是有限的并不等于知道v有多大。因此，这一信息无法使我们得知什么时候我们能确定最后那个需要由多于三个奇素数来表示的奇数。

\dialogname{阿基里斯：}我明白了。于是任何足够大的偶数2N都能表示成四个素数之和，因为可以先把2N-3表示成三个素数之和，然后再添上素数3。

\dialogname{乌龟：}正是如此。另一个角度的研究得出了一个很接近的结果，它是这样一个\heititext{定理}：``任何一个偶数都可以表示成一个素数和另外一个数的和，后者是至多两个素数的积。''

\dialogname{阿基里斯：}这个两素数之和的问题肯定把你引入了一个陌生境地。我想知道的是：如果你考虑两个素数之差，你会被引向哪里。要是我来做一小张表，列一些偶数以及它们如何表示成两个奇素数之差，就像我对和所做的那样，我敢打赌我会获得对这个难题的某种洞见。我们来看看\ldots\ldots{}

2=5-3, 7-5, 13-11, 19-17, 等等。

4=7-3, 11-7, 17-13, 23-19, 等等。

6=11-5, 13-7, 17-11, 19-13, 等等。

8=11-3, 13-5, 19-11, 31-23, 等等。

10=13-3, 17-7, 23-13, 29-19, 等等。

我的天！看上去对这些偶数来说，所能找到的表示成奇素数之差的方式是没有头的。可我到现在也还没看出这张表里有什么显见的规律。

\dialogname{乌龟：}也许本来就没有什么可以找得出的规律。

\dialogname{阿基里斯：}哼，你这家伙总是没完没了地鼓吹``紊乱''、``无序''。谢谢了，我可不想再听这种话。

\dialogname{乌龟：}你是不是认为每个偶数都能以某种方式表示成两个奇素数之差？

\dialogname{阿基里斯：}从我的表来看，回答当然是肯定的。不过我仍然可以假定回答也可能是否定的。这并没有真的使我们走的太远，是吧？

\dialogname{乌龟：}这里有很多值得注意的方面，我说的是在这个问题上会有更深入的洞见。

\dialogname{阿基里斯：}真奇怪，这个问题和哥德巴赫原来的那个问题这么相像。也许该称之为``哥德巴赫变奏''。

\dialogname{乌龟：}真该这样。不过，哥德巴赫猜想与这个哥德巴赫变奏之间有一个显著的不同之处。我很愿意跟你谈谈这一点。我们假设：如果一个偶数是两个奇素数之和，就是具有``哥德巴赫性质''；而如果它是两个奇素数之差，则具有``乌龟性质''。

\dialogname{阿基里斯：}我认为你该把它称作``阿基里斯性质''。说到底，是我提出的这个问题。

\dialogname{乌龟：}我正要说明，如果一个数不具有乌龟性质，则我们就说它具有``阿基里斯性质''。

\dialogname{阿基里斯：}好，那就这样吧\ldots\ldots{}

\dialogname{乌龟：}现在想一想，比如说，一万亿是否具有哥德巴赫性质，或具有乌龟性质？当然，它也可能两者兼备。

\dialogname{阿基里斯：}我可以想下去，可我怀疑我是否能就这两个问题中的任何一个给你一个回答。

\dialogname{乌龟：}别这么快就放弃。假如我是问你其中一个问题，前一个或者后一个，你会选择哪个问题去考虑？

\dialogname{阿基里斯：}我想我会扔一枚硬币。我没看出两者有什么差别。

\dialogname{乌龟：}啊哈！这可有天壤之别啊！如果你选了哥德巴赫性质，说的是素数之和，那么就仅限于使用那些介于2和一万亿之间的素数。对吧？

\dialogname{阿基里斯：}当然。

\dialogname{乌龟：}因此，你为一万亿找一个两素数之和的表示这一过程是注定会终止的。

\dialogname{阿基里斯：}啊！我知道关键所在了。要是我选择把一万亿表示成两个素数之差，涉及到的素数就会没有个头，有可能大到要花一万亿年才能找到它们。

\dialogname{乌龟：}或者还有可能它们根本就不存在！毕竟，问题就是这么提的------这样的素数是否存在？它们到底会有多大倒无关紧要。

\dialogname{阿基里斯：}对。如果它们不存在，那么搜寻过程就会永远进行下去，永远不作肯定答复，也不作否定答复。然而尽管如此，其答案是否定的。

\dialogname{乌龟：}所以，如果你有了某个数，想测试一下它是否具有哥德巴赫性质或者乌龟性质，两种测试的差别就在于，前者涉及的搜寻过程是注定会终止的，而后者则有可能是无穷尽的------没有任何保证。它会无休无止、没完没了，永远给不出一个回答。然而另一方面，在某种情形下，它或许会在第一步时就停住了。

\dialogname{阿基里斯：}我清楚了，哥德巴赫性质与乌龟性质之间是有着巨大差别的。

\dialogname{乌龟：}是的，而那两个相似的问题就是关于这两个极为不同的性质的。哥德巴赫猜想如成立，则所有偶数具有哥德巴赫性质；哥德巴赫变奏如成立，则所有偶数具有乌龟性质。两个问题都还没有解决，但有趣的是，虽然它们看起来很相像，它们涉及到的却是有关全体整数的极为不同的性质。

\dialogname{阿基里斯：}我明白你的意思。哥德巴赫性质对于任何一个偶数来说是可检测的，或者说是可识别的性质，因为我知道怎样测试它是否存在------进行搜索就是了。这个过程必定会到达一个终点，并给出肯定的或否定的回答。而乌龟性质则难对付多了，因为盲目的搜索可能会永远给不出回答。

\dialogname{乌龟：}啊，也许对于乌龟性质会有一些更聪明的搜索办法，按照其中的一个办法做下去总能到达一个终点，并给出回答。

\dialogname{阿基里斯：}难道这一搜索不是仅仅在答案是肯定时才会终止吗？

\dialogname{乌龟：}那倒不一定。没准儿会有什么办法证明：一旦搜索的时间超过某一限度，回答就必定是否定的。甚至也可能会有某种其他的搜索素数的办法------而不是这种盲目的方式------能在该性质存在时发现它，而若是不存在，也能给出回答。任何一种情形下否定的回答都只需有限的搜索。可我不知道这种事情是否能被证明。在无穷的空间里进行搜索总是很难驾驭的，这你也知道。

\dialogname{阿斯里斯：}所以就目前来看，你知道对乌龟性质来说还没有注定会终止的测试------然而没准儿还是会有这样一种搜索。

\dialogname{乌龟：}对。我可以设想某人从事对这种搜索的搜索，但我同样不能保证这个``元搜索''会终止。

\dialogname{阿基里斯：}这种怪异的情况给我的印象太深了，某个偶数------比如说一万亿------没能具有乌龟性质，竟是由无穷多支离破碎的信息所导致的。想想也挺有趣，把这些信息都卷在一起打成捆，称之为------就如你仗义的龟兄所说的------一万亿的``阿基里斯性质''。它事实上是作为整体的数论系统的一个性质，而不只是属于一万亿这个数的。

\dialogname{乌龟：}这是个很有意思的见解，阿基，不过我还要补充一点，把这一事实与一万亿这个数联系在一起是很有意义的。为了能说明这一点，我建议你考虑一个简单的陈述：``29是一个素数''。好，事实上这个陈述确实就意味着2乘以2不等于29，5乘以6不等于29，如此下去，对吧？

\dialogname{阿基里斯：}我认为必然是这样。

\dialogname{乌龟：}但是你把这些事实搜集起来，打成捆，并与29这个数联系起来，仅仅说上一句，``29是素数''？

\dialogname{阿基里斯：}是的\ldots\ldots{}

\dialogname{乌龟：}而涉及到的事实其数量实际上是无限的，对吧？毕竟，诸如``4444乘以3333不等于29''这一事实也是其中一部分，不是吗？

\dialogname{阿基里斯：}严格说来，我是这么认为。可你我都知道你不可能把两个比29都大的数乘起来得到29。所以实际上，说``29是素数''仅概括了有限多个关于乘法的事实。

\dialogname{乌龟：}你要是愿意，当然你可以这么看问题。但是你想想：两个大于29的数的乘积不可能等于29这一事实涉及了全体整数系统结构。在这种意义下，那事实本身就概括了无穷多的事实。你不能逃避这一事实，阿基，当你说``29是素数''时，你实际上表述了无穷多的事情。

\dialogname{阿基里斯：}也许是吧，可我觉得它像是单单一个事实。

\dialogname{乌龟：}那是因为有无穷多的事实已经包括在你预先拥有的知识里了------它们隐含地嵌在你看待事物的方式中了。你没有看到明显的无穷，是因为它已在你把握的图像中被隐含地捕获了。

\dialogname{阿基里斯：}我想你多半是对的。这似乎还是有点怪：把全体整数的性质集成一块，在这个块上贴一个标签，``29的素数性''。

\dialogname{乌龟：}也许看起来怪点儿，可这也是看待事物一种相当便利的办法。现在我们还回到你那个假设性的想法。如果像你所说的，一万亿这个数具有阿基里斯性质，那么不论你把它加上什么样的一个素数，你都不会得到另一个素数。这种状况会由无穷多个别的数学``事件''所导致。那么所有这些``事件''是否必须都是出自同一来源？它们是否必须有共同的缘起？因为若不是这样，那些事实就是由某种``无穷巧合''造成的，而非背后的规律性。

\pandocbounded{\includegraphics[keepaspectratio]{./media/OEBPS/Images/img_071.jpg}}

图71．秩序与紊乱

艾舍尔作（蚀版画，1950）

\dialogname{阿基里斯：}``无穷巧合''？自然数中没有巧合的事------没有任何事情不是出自某种背后的规律性。不说一万亿，就说7吧。处理它更容易些，它比较小。7具有阿基里斯性质。

\dialogname{乌龟：}你肯定？

\dialogname{阿基里斯：}是的。我告诉你为什么。如果你加上2，就得到9，不是素数。如果你加上随便什么其他素数，你就是把两个奇数加在一起，结果是个偶数------这样你还是得不到一个素数。因此在这里，7的``阿基里斯性质''不过是------杜馔一个术语------上两条理由的``逻辑后承''，而远不是什么``无穷巧合''。这恰好支持了我的论断：说明某个算术真理时决不需要无穷步骤的推理。假若有什么算术上的事实是由无穷多的无关巧合所导致，那么你将永远无法给出一个有穷的证明。但这是荒唐的。

\dialogname{乌龟：}你这看法有些道理，而且也不止你一个人这么想。不过------

\dialogname{阿基里斯：}难道还有人不同意这个观点吗？那么这种人就必然要相信``无穷巧合''了，他们必然相信在人类创造物中最完美、最和谐、最漂亮的部分------自然数系统------当中，也有紊乱。

\dialogname{乌龟：}也许他们就是这样，可你是否想过，这种紊乱或许是美与和谐的不可分割的组成部分？

\dialogname{阿基里斯：}紊乱，完美的一部分？秩序与紊乱构成了悦人的整体？真是异端邪说！

\dialogname{乌龟：}据说你最喜欢的艺术家艾舍尔就在一幅画里暗示了这种异端观点\ldots\ldots 说起紊乱这事，我相信你会有兴趣听听两类不同的搜索，两种搜索都保证会终止。

\dialogname{阿基里斯：}当然。

\dialogname{乌龟：}第一类搜索------非紊乱型的搜索------可以用一个关于哥德巴赫性质的测试来说明。你只需看那些小于2N的素数，如果某一对素数加起来得到2N，那么2N就具有哥德巴赫性质，否则就不具有。这种测试不仅注定会终止，而且你还能预先知道什么时候终止。

\dialogname{阿基里斯：}那么这是个可预知其终止的测试喽。你是不是要告诉我检验某些数论性质时，要用到那种虽然保证会终止，但预先无法知道会过多久才终止的测试？

\dialogname{乌龟：}你预见得太对了，阿基，这种测试的存在就说明自然数系统从某种意义上说带有固有的紊乱。

\dialogname{阿基里斯：}这种情况下，我只好说人们还不够了解那个测试。如果他们多作些研究，他们会弄清在终止前那个测试至多要花多少时间。无论如何，整数中的模式必定是有来由的，不可能只是一些无法预知的紊乱模式！

\dialogname{乌龟：}我可以理解你这种直觉信念，阿基。不过，这一点并非总能得到证实。当然，很多情况下你是绝对正确的------仅仅因为某人不知道某事，不能得出结论说那件事不可知。可是有那么一些整数的性质，可以证明检验它们的有终测试是存在的，同时还能证明没有办法预先知道那些测试要花多少时间。

\dialogname{阿基里斯：}真难以相信。听起来就像是可恶的魔鬼阴险诡诈地溜进了自然数这属于上帝的美丽乐园！

\dialogname{乌龟：}有一点大概能安慰安慰你：想要定义出一个性质，使得能有一个有终止测试来检验它，但其终止是不可预知的，这绝不是件容易的事情，也绝不那么自然。大多数``自然的''整数性质的确都能接受可预知其终止的测试。例如：素数性、可开方性、是否是10的方幂等等。

\dialogname{阿基里斯：}是的，这些性质的测试全都是直截了当的。你告诉我一个那样的性质吧：检验它的唯一方法是用一个有终止测试，但不可预知其终止。

\dialogname{乌龟：}说起来太复杂了，我现在直犯困，换一个吧。我来跟你说一个很容易定义的性质，然而却还没有一个检验它的有终止测试。我没说永远不会发现这样一个测试，这一点请你注意------只是现在还没有。我们从一个数开始吧------你来选一个数。

\dialogname{阿基里斯：}15怎么样？

\dialogname{乌龟：}很好。我们从你这个数开始。如果它是奇数，我把它乘以3，再加1。如果它是偶数，就取其一半。然后重复这个过程。如果一个数经过这样的一些操作得到1，就称之为``妙极''的数，否则就称之为``非妙极''的。

\dialogname{阿基里斯：}那么15是妙极的还是非妙极的呢？我们看看吧。

\begin{longtable}[]{@{}lll@{}}
\toprule\noalign{}
\endhead
\bottomrule\noalign{}
\endlastfoot
15 & 　是奇数，所以做3n+1： & 46 \\
46 & 　是偶数，所以取一半： & 23 \\
23 & 　是奇数，所以做3n+1： & 70 \\
70 & 　是偶数，所以取一半： & 35 \\
35 & 　是奇数，所以做3n+1： & 106 \\
106 & 　是偶数，所以取一半： & 53 \\
53 & 　是奇数，所以做3n+1： & 160 \\
160 & 　是偶数，所以取一半： & 80 \\
80 & 　是偶数，所以取一半： & 40 \\
40 & 　是偶数，所以取一半： & 20 \\
20 & 　是偶数，所以取一半： & 10 \\
10 & 　是偶数，所以取一半： & 5 \\
5 & 　是奇数，所以做3n+1： & 16 \\
16 & 　是偶数，所以取一半： & 8 \\
8 & 　是偶数，所以取一半： & 4 \\
4 & 　是偶数，所以取一半： & 2 \\
2 & 　是偶数，所以取一半： & 1 \\
\end{longtable}

嚯！从15到1这么转弯抹角的一大圈！不过我终于还是得到1了。这说明15是具有妙极这一性质的。那么什么数会是非妙极的呢\ldots\ldots{}

\dialogname{乌龟：}你注意了没有，在这么简单定义的运算下，那些数是怎样上下摆动的？

\dialogname{阿基里斯：}注意了。特别让我吃惊的是，经过十三次运算后，我得到16，仅比我起步时的15大1。在某种意义上说我几乎回到了出发点------而在另一种意义上，我已远离了出发点。另外还有一点使我诧异的是为了解决问题，我不得不动用大到160那样的一个数。我真想知道这都是为什么。

\dialogname{乌龟：}是的，有个无穷的``天空''供你飞，而且事先很难知道你会被抛到这天空里多高的地方。而且很可能你会只是不断地往高飞，越来越高，再也不下来了。

\dialogname{阿基里斯：}是吗？我猜想这倒是可能的------可这得要多奇特的巧合啊！你得正好碰上一个接一个的奇数，其中只掺有很少的几个偶数。我怀疑这到底是否会发生------可我不能肯定。

\dialogname{乌龟：}你从27开始怎么样？不过，我可是什么都没保证啊。有时间的话你试试，就当是玩吧。我劝你准备好一大张纸。

\dialogname{阿基里斯：}嗯\ldots\ldots 听起来挺有意思。说实在的，把妙极（或非妙极）与那个起始的数联系在一起还是让我觉得有点不安，即使很明显那是全体整数系统的性质。

\dialogname{乌龟：}我明白你的意思，``29是素数''，或``金子是宝贵的''这两句话没什么不同------它们都是说某一特定实体被嵌入特定语境，从而具有某一性质。

\dialogname{阿基里斯：}我想你是对的。这个``妙极''问题真是奇妙至极了，就看那些数的摆动方式吧------一会儿增大，一会儿减少。模式应该是有规则的，可是表面上却显得相当紊乱。所以我完全能想象为什么至今还没有一个检验妙极性质的有终止测试。

\dialogname{乌龟：}说起有终止过程和无终止过程，还有那些介于两者之间玩艺儿，我想起了我的一个朋友。他是个作家，正在写一本书。

\dialogname{阿基里斯：}哦？这太好了！是本什么书？

\dialogname{乌龟：}《金、银、铜------聚宝藏之精华》。挺有意思吧？

\dialogname{阿基里斯：}老实说，这名字弄得我有点糊涂了。说到底，金、银、铜之间有什么关系？

\dialogname{乌龟：}对我来说似乎是很清楚的。

\dialogname{阿基里斯：}假若书名是，比如说，``金、马、牛''，或者``风、银、牛''，嗨，我就可以把它看成\ldots\ldots{}

\dialogname{乌龟：}也许你更喜欢``风、马、铜''？

\dialogname{阿基里斯：}噢，太对了！但原来那个书名太糟糕。没人能懂。

\dialogname{乌龟：}我会告诉我的朋友。他会很高兴能有一个生动易记的书名（他的出版商也会的）。

\dialogname{阿基里斯：}我真高兴。可我们的讨论怎么会让你想起他的那本书了？

\dialogname{乌龟：}啊，是这样。那本书里有篇对话，在其中他想让读者去搜索终了，从而甩掉读者。

\dialogname{阿基里斯：}竟会有这种愿望，够奇怪的。这怎么做到呢？

\dialogname{乌龟：}你无疑会注意到，一些作者在他的故事结束前的几页里如何费尽心机地建立起巨大的紧张状态------可是那个把书捧在手上的读者能感觉出故事就要结束了。所以，有某种额外的信息以一定的方式预先提醒了读者。书的物理性质多少有点破坏了那种紧张。如果，比如说，在小说的终了处有许多废话，那就会好多了。

\dialogname{阿基里斯：}废话？

\dialogname{乌龟：}对。我的意思是，有许多多余的书页，不是故事的真正组成部分，但能用来藏住终了，使得草率浏览或只凭对书的物理感觉都无法确切地得知终了在哪里。

\dialogname{阿基里斯：}我明白了。这样的话，故事的真正终了可能会出现在，比如说，书的物理终了前五十页或前一百页了？

\dialogname{乌龟：}是的。这就有了能让人兴奋的东西，因为读者事先不知道有多少页是废话，多少页是故事。

\dialogname{阿基里斯：}如果这能广为实践，会是很有效的。可是有个问题：假如这些废话非常明显------比如有许多空白，或者有充满句号或满篇杂乱无章的字词的书页，那还不如没有。

\dialogname{乌龟：}你说得对。得弄得像是正常的一页书。

\dialogname{阿基里斯：}可是一般说来即使草率地浏览正常的一页书也足以把一个故事同另一个故事区分开来。所以你必须把那些废话弄得与真正的故事非常相似。

\dialogname{乌龟：}没错。我的设想是这样：把故事讲完后不中断，继续讲下去，讲一些看上去是该故事的继续，而实际上只是些废话，与真正的主题毫无关系的事。某种意义上说，这些废话是种``终了后的终了''，带有与原主题几乎无关的多余的文学意象。

\dialogname{阿基里斯：}真够狡猾的！可这就有个问题：你无法弄清真正的终了到底在哪儿，都混在废话里了。

\dialogname{乌龟：}这也就是我的作家朋友和我所得出的结论。这是个问题，可我还是觉得那个设想相当吸引人。

\dialogname{阿基里斯：}哎，我有个建议。可以把从真正的故事到废话部分的转折做成这样：一个有智能的阅读者对正文充分专注仔细地审查之后，能查明什么地方是前者的结束、后者的开始。也许这会使他花费相当的时间。也许没有办法预先知道到底要花多长时间\ldots\ldots 但出版者能够保证对真正终了的一个充分专注仔细的搜索一定会终止，即使他无法在其终止前说出到底要花多少时间。

\dialogname{乌龟：}很好------但``充分专注仔细''是什么意思？

\dialogname{阿基里斯：}意思是：读者必须专注地寻找出现在正文某处的某种微小、然而却说明问题的特征。那特征标志着终了。他还必须有足够的聪明去思考并搜寻许多这类特征，直到发现了真正标志终了的那个特征。

\dialogname{乌龟：}比如说某些字的出现频率的突然变化，或一些成语的一部分被莫名其妙地丢掉了，或者突然涌现出一连串语法错误？

\dialogname{阿基里斯：}可能是吧。或者是某种隐藏着的信息，对于一个充分专注仔细的读者，它会揭示出真正的终了。你甚至还可以加进一些额外的角色或事件，它们与前面故事的精神很不协调。一个头脑简单的读者会囫囵吞枣地一直把书读完，而有头脑的读者就能够准确地找到分界线在哪里。

\dialogname{乌龟：}这真是个有独创性的想法，阿基。我会电告我的朋友，说不定他在他的对话里就用上你这方法了。

\dialogname{阿基里斯：}我不胜荣幸。

\dialogname{乌龟：}好啦，阿基，我得回去了。我恐怕都有点站不稳了，趁现在还能找着路，我最好现在就走。

\dialogname{阿基里斯：}你为了我而在这儿待了这么长时间，而且是在夜里的这种时候，真让我受宠若惊。我敢保证你那数论娱乐对于我的辗转不眠是绝好的良方。真的------没准我今晚能睡个好觉了。龟兄，我愿送给你一件特别的礼物，以表我的感激之情。

\dialogname{乌龟：}噢，别这样，阿基。

\dialogname{阿基里斯：}这是我的一点心意，龟兄。到柜子那儿去，你会在上面找一个金色的小包。

（乌龟慢吞吞地踱到阿基里斯的柜子前。）

\dialogname{乌龟：}你不是指这个人造革的包吧？

\dialogname{阿基里斯：}就是那个。不过那不是人造革的，它是变造革的。

\dialogname{乌龟：}什么是变奏歌？

\dialogname{阿基里斯：}哦，你听岔了，不是什么``变奏歌''，我是说``变造革''，这是种金、银、铜的聚合物，看上去像皮子。请收下吧，龟兄，作为我最衷心的祝愿。

\dialogname{乌龟：}的确太谢谢你啦，阿基，嗯\ldots\ldots 哎，这包上面印的是什么呀？哦，是英文呐，多古怪的一份名单啊：

\heititext{D}e Morgan

A\heititext{b}e1

Bo\heititext{o}1e

Bro\heititext{u}wer

Sier\heititext{p}inski

Weier\heititext{s}trass

\dialogname{阿基里斯：}我相信这一定是一份完整的所有伟大数学家的名单。

\dialogname{乌龟：}真的吗？我看看------德·摩根、阿贝尔、布尔、布劳维、谢尔平斯基、维尔斯特拉斯------哦，没错，他们确实全都是世界著名的数学家。可这名单上怎么没有康托尔的名字？让我瞧瞧，我想他的名字的拼法是``Cantor''。

\dialogname{阿基里斯：}我真想象不出他的名字怎么能不包括在内。我也弄不明白为什么对角线上的那些字母用了黑体。

\dialogname{乌龟：}下面写着：``从该对角线中减1，就能找到巴赫在莱比锡的职位''。

\dialogname{阿基里斯：}我也看到了，可我完全不知所云，哎，喝上两口上好的威士忌怎么样？我架子上的那只瓶子里正好有点儿。

\dialogname{乌龟：}不必啦。我已经累坏了。我得赶紧回去了。\kaitismall{（他漫不经心地打开了小包。）}哎，等等，阿基------这里面装的是一百枚金币！

\dialogname{阿基里斯：}你要是能收下，我将非常非常高兴，龟兄。

\dialogname{乌龟：}可是------可是------

\dialogname{阿基里斯：}别再客气了。那小包，那金子------都是你的了。为这个无与伦比的晚上，我衷心地感谢你。好了，再见吧，龟兄。

\dialogname{乌龟：}我没能做到善始善终，我本想帮你克服失眠症，可现在只好无功受禄、不了了之。好吧，谢谢你的礼物，再见。

（这时，传来什么人的敲门声。）

都这么晚了，谁还会来敲门？

\dialogname{阿基里斯：}我一点儿也想不出来会是谁。我觉得多少有点问题。你还是躲到柜子后面去吧，免得碰上什么麻烦事。

\dialogname{乌龟：}好主意。\kaitismall{（费力地挤到柜子后面。）}

\dialogname{阿基里斯：}谁啊？

\dialogname{声音：}开门！我们是警察！

\dialogname{阿基里斯：}请进，门开着呢。

（两个粗壮的警察推门进来，身上带着闪亮的徽章。）

\dialogname{警察：}\kaitismall{（指着自己的铜徽章）}我姓银。他姓金。这里是不是住着一个叫阿基里斯的？

\dialogname{阿基里斯：}我就是。

\dialogname{警察：}那好，阿基里斯，我们有理由相信有一个变造革的包在这儿，里面有一百枚金币。今天下午有人在博物馆里把它弄到手后溜走了。

\dialogname{阿基里斯：}哎哟我的老天！

\dialogname{警察：}如果是这儿，阿基里斯，由于你是唯一可能的嫌疑，很遗憾我必须拘留你。我带来了搜查证------

\dialogname{阿基里斯：}噢。先生们，你们来我真高兴！我被乌龟先生还有他那个变造革的包吓得一晚上都不自在。现在你们终于来了，这下我可得救了！快，先生们，就在那柜子后面，罪犯就在那里！\kaitismall{（警察往柜子后面看去，看见乌龟在里面缩成一团，抱着他那变造革的包，浑身打颤。）}

\dialogname{警察：}就在这儿！那么乌龟先生就是那个混蛋喽，嗯？我可没想到是他！可是他被当场抓住了！

\dialogname{阿基里斯：}把这恶棍带走，好心的先生们！谢天谢地，这可将是我最后一次见他和那个变造革的包了！

\phantomsection\label{Chapter13.xhtml}{}

第十三章

\chapter{BlooP和FlooP和GlooP}

自我意识和无序

BlooP、FlooP和GlooP不是神话中的巨人，不是唐老鸭的小侄子们，也不是船沉时发出的冒泡声------它们是三种计算机语言，其中每种都有特殊的用途。这些语言是专门为本书的这一章发明的。它们将被用来解释``递归''这个词的某些新意义------特别是``原始递归''和``一般递归''这两个概念。事实将证明，这些语言有助于阐明TNT中自我相关的机制。

我们似乎很突然地从大脑和心智跳到了数学和计算机科学中的技术。虽然这个跳跃从某些方面来看有点突然，但它还是有意义的。我们已经看到，某种自我意识似乎是意识的关键所在。现在我们要在更形式化的背景下进一步分析``自我意识''，例如在TNT的背景下。在TNT和心智之间还有很大一段距离，但一些想法将会是富有启发性的，或许会以隐喻的方式被传回到我们关于意识的思考之中。

关于TNT的自我意识，令人惊奇的一点就是它密切地联系于自然数中的有序与无序问题。特别是我们将会看到，一个充分复杂以至能反映自身的有序系统不可能是完全有序的------它必定包括某种奇怪的无序特征。对于那些具有某种阿基里斯式想法的读者来说，这将是很难接受的。但是，存在一种``魔术式的''补偿------存在某种关于无序的有序，这本身已经形成了一个研究领域，称为``递归函数论''。遗憾的是，我们所能做到的，只是略微展示一下这个课题的魅力。

可体现性与冰箱

像``足够复杂''或``足够强有力''及类似的说法，在前面已经多次出现了。但这些说法是什么意思呢？让我们回到龟蟹之战，并问这样一个问题：``是什么使某个东西有资格成为一台唱机的？''螃蟹可能会声称它的冰箱是一台``完备的''唱机。然后为了证明这一点，它可能会随便拿个唱片放在冰箱顶上，说：``你看------它在播放这张唱片了。''而乌龟，如果想要反驳这个禅宗式的行动，就必须说：``不------你的冰箱保真度太低，已经不能算是唱机了：它根本不能重现唱片上记录的声音（更不必说那种让它自我破坏的声音了）''。只有当``唱机X''真是一台唱机时，乌龟才能造出一张叫做``我不能在唱机X上播放''的唱片。乌龟的方法很狡猾，因为它抓的是系统的强项，而非弱项。因此它要求``保真度足够高''的唱机。

对形式化的数论系统来讲也是一样。说TNT是一个形式化的N，其原因是TNT的符号以正确的方式活动：这也就是指它的定理不像冰箱似的一声不响------它们确实说出了N中的真理。当然，pq系统中的定理也是这样。pq系统也能作为一个``形式化的数论系统''吗？还是说它更像个冰箱？它比冰箱是强一点，但仍然很弱。被pq系统所包含的N中的核心真理太少了，因此它还不能称作是一个``数论系统''。

那么，这些N中的``核心真理''是什么呢？它们是``原始递归真理''，这就是说，它们仅仅涉及到``可预测其终止''的计算。这些核心真理在N中的作用就像欧里几得的前四个公设在几何学中的作用一样：它们使你可以在比赛开始前就能淘汰掉某些``力量不够强''的选手。从此以后，``全部原始递归真理的可体现性''将作为我们判别一个系统是否``足够强有力''的标准。

元数学中的岩头之斧

上述观念的重要性表现在下列关键性的事实之中：如果你有数论的一个足够强有力的形式化体现，那么哥德尔定理就是可应用的，结果你的系统是不完备的。反过来，如果你的系统不足够强有力（即不是所有原始递归真理都是定理），那么因为这个缺陷，它也不完备。这里我们在元数学中碰到了``岩头之斧''的另一种形式化体现：不管系统怎么办，哥德尔之斧总要砍掉它的脑袋！另一点需要注意的是，这完全平行于《对位藏头诗》中的``高低保真度之别''。

实际上，人们发现很弱的系统依然会受到哥德尔方法的攻击。``所有原始递归真理都需要体现成定理''，这个标准过于严格了。这有点像这样一个贼：他只偷``足够阔气''的人，其标准是盗窃对象应当至少带有一百万元现金。幸运的是，对TNT系统来说，我们将能够像这个贼那样动手，因为那里确实有一百万元现金------这就是说，TNT确实包含了所有原始递归真理作为其定理。

现在，在我们进入对原始递归函数的谓词的详细讨论之前，我想把这一章的主题和前几章的主题联系起来，以提供一个更好的背景。

选择适当的过滤器来发现有序

我们很早就已经看到，形式系统可能是难以驾驭的，因为它们有加长和缩短符号串的规则，这可能会导致在大量符号串中进行无终止的搜索。哥德尔配数法的发现，显示了任何对一个具有特殊的符号性质的串的搜索都有一个算术中的表兄弟：对一个具有相应的特殊算术性质的整数的同构搜索。结果，对于形式系统判定过程的要求，涉及了对整数中不可预测其长度的搜索之谜（一种无序）的解决。在本章前面的对话中，我可能过分强调了和整数有关的问题中显示出的无序现象。事实上，和``妙极性''问题相比，人们已经驯服了一些更复杂的无序现象，发现它们不过是些很温顺的家伙。因此，阿基里斯对数的规则性和可预测性的坚定信仰，应当得到一定的尊敬------特别是由于这反映了到30年代为止几乎所有数学家的信念。为了说明为什么有序与无序的对比是这样一个微妙而有意义的问题，也为了把这一问题和对意义的定位及揭示等问题联系起来，我想引用《量子是实在的吗？》一文中优美而又令人难忘的一段------这是已故的杰·尧奇写的一篇伽利略式的对话：

萨尔维亚蒂：假设我给你两个数列，如

7 8 5 3 9 8 1 6 3 3 9 7 4 4 8 3 0 9 6 1 5 6 6 0 8 4\ldots\ldots{}

和

1, -1/3, +1/5, -1/7, +1/9, -1/11, +1/13, -1/15\ldots\ldots{}

如果我问你，辛普利奇奥，第一个序列的下一个数是什么，你会怎样说？

辛普利奇奥：我无法告诉你。我觉得它是个随机序列，其中没有规律。

萨尔维亚蒂：那么第二个序列呢？

辛普利奇奥：这个容易，下一个一定是+1/17。

萨尔维亚蒂：正确。但如果我告诉你第一个序列也是按某个规律构造的，而且这个规律事实上和你刚才在第二个序列中发现的一样，你会说什么？

辛普利奇奥：我看这不大可能。

萨尔维亚蒂：但事实的确如此。因为第一个序列只不过是第二个序列的和的十进制小数部分（展开式）的开始几位。这个和是π/4。

辛普利奇奥：你用的数学技巧太多，但我看不出这与抽象和实在有什么关系。

萨尔维亚蒂：和抽象的关系显而易见。除非你已经通过抽象过程开发了一种过滤器，用它可以看出在表面的随机性背后的简单结构，否则第一个序列看上去就是随机的。

自然规律恰恰是以这种方式被发现的。自然只是提供给我们大量的现象，它们表面上杂乱无章，直到我们选择了某些有意义的事件，而且把它们从特定的，关系不大的环境中抽象出来，使它们成为理想化的时候为止。只有那时它们才会展现出光彩夺目的真实结构。

萨哲杜：这个想法太妙了！由此看来当我们设法理解自然时，我们应当把现象看作待理解的消息。只不过每个消息在我们为它设立一种编码之前好像是随机的。这种编码表现为一种抽象，就是说，我们有选择地把某些无关的东西忽略掉，这样就通过一种自由选择选取了消息的部分内容。这些无关信号构成了``背景噪声''，它们将限制我们的消息的精确性。

但由于这种编码不是绝对的，在同一组作为原材料的数据中就可能存在着多个消息，因此改变编码就会在某些以前认为只是噪声的东西中发现同样深刻的意义。相反地，在新的编码中以前的消息可能会变得没有意义。

这样一种编码就预先设定了在不同的、互补的方面之间的一种自由选择。每个方面都同样自称是``真实的''，如果我可以用这个可疑的词的话。

在这些方面之中，有些可能目前还完全不为我们所知，但它们可能把自己暴露给一个具有完全不同的抽象系统的观察者。

但请告诉我，萨尔维亚蒂，这样，我们怎么能仍然声称我们在客观的实在世界中发现了某些事物呢？这难道不是意味着我们只不过是根据自己的想象在构造事物，而实在仅仅存在于我们自身之中吗？

萨尔维亚蒂：我不认为一定是这样的，但这是一个需要进一步反思的问题。\footnote{}

在这里，尧奇所涉及到的消息不是来自一个``有感知能力的物体''，而是来自自然界本身。我们在第六章提出的关于意义和消息的关系的问题同样适用于来自自然的消息。自然是无序的还是模式化的？智能在确定这个问题的答案的过程中能起什么作用？

从哲学中退出来，我们仍能考虑表面上的随机序列中深藏的规律性。第五章中的函数Q(n)是否也具有一个简单的非递归解释？是不是每个问题都像《一首无的奉献》中提到的那片果树林一样，从某个特定的角度看进去，其中的秘密就一览无余了？还是说在数论中存在着一些不论从哪个角度看都是神秘的问题？

有了这段开场白，我觉得现在应该继续前进，去定义所谓``长度可预测的搜索''的精确意义了。这将用BlooP语言来完成。

BlooP语言的基本步骤

我们的议题是搜寻具有各种性质的自然数。为了讨论任意搜索的``长度''，我们必须定义一些基本``步骤''，任何搜索都由它们组成。这样，搜索长度就可以根据其中的步骤数来度量。被我们当作基本步骤的有：

\begin{itemize}
\tightlist
\item
  两个自然数相加；
\item
  两个自然数相乘；
\item
  确定两个数是否相等；
\item
  确定两个数的相对大小。
\end{itemize}

循环和上界

如果我们想用这些步骤严格地构造一个测试，例如测试一个数是否是素数，我们很快就会发现必须在其中包含一个``控制结构''------即对操作次序的描述：何时需要回过头来重新尝试某些东西，何时跳过一些步骤，何时停止，以及诸如此类的事情。

在典型的情况下，任何``算法''------即对任务完成过程的明确描述------都由下列成分组成：（1）需完成的特定运算，（2）控制语句。因此，当我们为表示长度可预测的计算开发我们的语言时，我们必须同时包括基本控制结构。事实上，BlooP的特色也恰在于其有限的控制结构集。它不允许你转移到任意的步骤或无限制地循环某一步骤。在BlooP中，基本上唯一的控制结构就是``有界循环''（即``bounded loop''，这也正是``BlooP''的来历）：重复执行一组指令，重复次数不能大于某个预先确定的最大值。这个最大值称为该循环的``上界''，或``顶''。如果顶是300，那么这个循环可以进行0次、7次或300次------但不能是301次。

在一个程序中，并不要求程序员准确地给出所有上界的数值------事实上它们是无法预知的。但是，每个上界都可以在进入该循环之前通过计算来确定。例如，如果你要计算\gebfont{ꪢ}的值，就该设置两个循环。首先，你求出3\textsuperscript{n}的值，这包括n次乘法。然后，你求2的3\textsuperscript{n}次方，这包括3\textsuperscript{n}次乘法。这样，第二个循环的上界是第一个循环的计算结果。

下面说明你怎样才能把这个过程表示在一个BlooP程序之中：

\begin{longtable}[]{@{}ll@{}}
\toprule\noalign{}
\endhead
\bottomrule\noalign{}
\endlastfoot
DEFINE PROCEDURE ``TWO-TO-THE-THREE-TO-THE''{[}N{]}： & 定义过程名和参数N \\
BLOCK 0: BEGIN & 第0块开始 \\
　CELL(0)\gebfont{⇐}1 & 把``1''送入0号单元 \\
　　LOOP N TIMES & 循环执行N次 \\
　BLOCK 1: BEGIN & 第1块开始 \\
　　CELL(0)\gebfont{⇐}3*CELL(0) & 0号单元内容乘``3''结果仍送入0号单元 \\
　BLOCK 1: END & 第1块结束 \\
　　CELL(1)\gebfont{⇐}1 & 把``1''送入1号单元 \\
　　LOOP CELL(0) TIMES： & 循环次数在0号单元中 \\
　　BLOCK 2: BEGIN & 第2块开始 \\
　　CELL(1)\gebfont{⇐}2*CELL(1) & 1号单元内容乘``2''结果仍送入1号单元 \\
　　BLOCK 2: END & 第2块结束 \\
　　OUTPUT\gebfont{⇐}CELL(1) & 1号单元内容送入输出单元 \\
BLOCK 0: END & 第0块结束 \\
\end{longtable}

BlooP的约定

能够看着一个用计算机语言书写的算法，指出它是做什么的，这需要学习。但是，我希望这个算法足够简单，不必仔细研究就能理解。在上述程序中，定义了一个``过程''，它带有一个``输入参数''N，它的``输出''就是所要得到的值。

这个过程定义采用所谓``块结构''，这就是说其中某些部分可以被看作一个单元，或``块''。在一块中的所有语句可以被当作一个单元来执行。每一块都有个编号（最外面一层是BLOCK 0），而且用BEGIN和END标记开始和结束。在我们这个例子中，BLOCK 1和BLOCK 2每个仅含一个语句------但不久你就会见到更长的块。一个LOOP语句总是意味着重复地执行紧随其后的块。正像在上例中所见到的，块可以互相嵌套。

上述算法的策略前面已经描述过了。你首先取一个辅助变量，叫CELL(0)，把它的初值赋为1，然后，在一个循环中反复地把其值乘以3，直到这样做完N次为止。下一步，你对CELL(1)做类似处理------把它置成1，反复乘2，做CELL(0)次，然后停止。最后，你把CELL(1)中的值赋给OUTPUT。这个值就是要传回给外部世界的------这是该过程唯一的外部可见行为。

关于所用的记号，在这里还有几点需要说明。首先，左箭头``\gebfont{⇐}''的意思是这样：

求出其右边公式的值，然后把结果赋给其左边的CELL（或OUTPUT）。

因此CELL(1)\gebfont{⇐}3*CELL(1)这条命令的意义就是把存在CELL(1)中的值乘3倍。你可以把每个CELL想象成某个计算机的存贮器中的一个字。在一个CELL和一个真实的字之间的唯一区别就是：后者只能存放某个确定限度之内的整数，而我们允许一个CELL存放任意自然数，不管有多大。

在BlooP中，每个过程在被调用时都能得出一个值------即称为OUTPUT的变量的值。在开始执行任何过程时，OUTPUT都把0值作为其缺席选择值。这样，即使过程根本没给OUTPUT重新赋值，OUTPUT也能始终有一个明确定义的值。

IF条件语句和分支

现在让我们看另一个过程，它将展示BlooP的其它特点，这些特点使该语言更加具有一般性。如果只会做加法，你怎样才能求出M-N的值呢？窍门在于把各种各样的数和N相加，直至你发现一个与N的和等于M的数为止。但是如果M小于N会发生什么情况呢？如果我们想从2中取走5会怎样？这在自然数域中是没有答案的。但我们希望我们的BlooP过程不管怎样应当给出一个答案------比如说0。下面就是一个做减法的BlooP过程：

\begin{longtable}[]{@{}ll@{}}
\toprule\noalign{}
\endhead
\bottomrule\noalign{}
\endlastfoot
DEFINE PROCEDURE ``MINUS''{[}M, N{]}: & 定义过程名和变量M \\
　BLOCK 0: BEGIN & 第0块开始 \\
　IF M＜N, THEN & 若M＜N，则： \\
　QUIT BLOCK 0 & 退出第0块 \\
　LOOP AT MOST M+1 TIMES: & 循环至多M+1次； \\
　BLOCK 1: BEGIN & 第1块开始 \\
　IF OUTPUT+N=M, THEN： & 若OUTPUT+N=M，则： \\
　ABORT LOOP 1 & 中止循环1 \\
　OUTPUT\gebfont{⇐}OUTPUT+1 & OUTPUT的值加1，送回 \\
　BLOCK 1: END & 第1块结束 \\
　BLOCK 0：END & 第0块结束 \\
\end{longtable}

这里我们利用了OUTPUT开始时为0这一隐含性质。如果M小于N，那么减法就是不可能的，我们只需直接跳到BLOCK 0的末尾，同时答案是0。这就是QUIT BLOCK 0这一行的意思。但如果M不小于N，那么我们就越过这条QUIT语句，按顺序执行下一条命令（在这里是个LOOP语句）。这就是IF语句在BlooP中的工作方式。

这样，我们就进入了LOOP 1。这样称呼它，是因为它要重复执行的块是BLOCK 1。我们试着把0加到N上，然后是1、2、3\ldots\ldots 直到我们找到那个能导出M的数。这时我们对当前循环做ABORT，即跳到紧跟着标识循环块底的END后面的一条语句。在上述例子中，我们跳到了BLOCK 1: END的后面，即到达了算法的最后一条语句，而且任务已经完成了。现在OUTPUT中放的就是正确答案。

注意，往下跳有两种不同的指令：QUIT和ABORT。前者从属于块，后者从属于循环。QUIT BLOCK n意味着跳到BLOCK n的最后一行，而ABORT LOOP n意味着跳到BLOCK n的最后一行之后。只有在一种情况下二者的区别是重要的：当你在一个循环之中，想要继续循环，但又要退出这一轮所在的块。这时你用QUIT就能达到预期的效果。

还应当注意的是在循环上界前面注有AT MOST（至多），这告诫你循环可能在达到上界之前就被中止了。

自动组块

下面是BlooP的最后两个需要加以说明的特点，它们都是非常重要的。其一是：一旦一个过程已经被``定义''了，它就可以在后面的过程定义中被``调用''。其效果是一旦某个操作在一个过程中被定义了，它就会被认为和基本步骤一样简单。这样，BlooP就具有了自动组块的特点。你可以把这和一个溜冰好手学习新花样的过程相比：他不是把新动作看作长长的基本肌肉活动序列，而是看作以前学过的一些动作的组合，而那些动作自身又是用更早所学的动作所组成的，如此等等------这种嵌套或组块可以上溯许多层次，直至遇到基本肌肉活动为止。这样，BlooP程序的能力就和溜冰者的技能一样，可以突飞猛进地增长了。

BlooP测试

BlooP的另一个特点是：某些特定过程可以用YES或NO作为其输出，而不是输出一个整数的值。这种过程是``测试''，而非``函数''。为了说明这个区别，测试的名字必须以一个问号结束。另外，在一个测试中，OUTPUT的缺省值当然不能是0，而是NO。

让我们看一个例子，其中把BlooP的最后这两个特点表现在一个算法里。这个算法要测试一个函数的自变量是否是素数。

\begin{longtable}[]{@{}ll@{}}
\toprule\noalign{}
\endhead
\bottomrule\noalign{}
\endlastfoot
DEFINE PROCEDURE ``PRIME?''{[}N{]}: & 定义过程名和变量N \\
BLOCK 0: BEGIN & 第0块开始 \\
　IF N\textless0, THEN: & 若N\textless0，则： \\
　QUIT BLOCK 0； & 退出第0块 \\
　CELL(O) \gebfont{⇐}2 & 把CELL(0)置成2 \\
　LOOP AT MOST & 循环至多N-2次 \\
　MINUS{[}N, 2{]}TIMES & \\
　BLOCK 1 & BEGIN第1块开始 \\
　　IF REMAINDER {[}N, CELL(0){]}=0, THEN: & 若N除以CELL(0)余0，则： \\
　　QUIT BLOCK 0 & 退出第0块 \\
　　CELL(0) \gebfont{⇐}CELL(0)+1 & CELL(0)的值加1 \\
　BLOCK 1: END & 第1块结束 \\
　OUTPUT \gebfont{⇐}YES & 把OUTPUT置成YES \\
BLOCK 0: END & 第0块结束 \\
\end{longtable}

注意，我在这个算法中调用了两个过程：MINUS和REMAINDER（假设后者已经事先定义了，你可以自己完成这个定义，REMAINDER{[}N, R{]}的作用是求N除以R所得的余数）。这个素数测试过程的工作方式是逐个检查N可能有的因子，从2开始直到N-1为止。如果其中某一个能整除N（即余数为0），则跳到过程尾部，由于此时OUTPUT仍然保有其缺省值，故答案为NO。仅当N没有整数因子时，它才能通过整个LOOP 1，那时我们会遇到语句OUTPUT \gebfont{⇐}YES，执行了这个语句后，过程就结束了。

BlooP程序包含了过程组成的链

我们已经看到了在BlooP中怎样定义过程，但过程定义只是一个程序的一部分。一个``程序''包括一条``过程定义链''（其中每个过程仅仅调用前面定义的过程），还可能跟着一个或多个对所定义过程的``调用''。例如，过程TWO-TO-THE-THREE-TO-THE的定义和后面跟着的调用

TWO-TO-THE-THREE-TO-THE{[}2{]}

合起来构成了一个完整的BlooP程序，它的运行结果是512。

如果你只有一条过程定义链，那什么东西都不会被执行。它们就在那里等着调用，等着用特定的数值启动它们。这就像一台绞肉机等着肉来绞------或者说更像一串连在一起的绞肉机，前面一台绞完了就送进后一台\ldots\ldots 对于绞肉机来说，这种安排可能没什么必要，但对于BlooP程序来说，这种结构是非常重要的。我们可以可以称它为一个``待调用程序''。这个概念显示在图72中

\pandocbounded{\includegraphics[keepaspectratio]{./media/OEBPS/Images/img_072.png}}

图72．一个待调用BlooP程序的结构

由于这个程序是自足的，每个过程定义中只能调用定义在它前面的过程。

现在BlooP是我们用来定义可预知终止的计算的语言。可以用BlooP来计算的函数的标准名称是``原始递归函数''；而可以用BlooP测试来验证的性质的标准名称是``原始递归谓词''。这样函数\gebfont{ꪢ}就是个原始递归函数，而命题``n是个素数''则是个原始递归谓词。

凭直觉就能看出哥德巴赫性质是原始递归的。为了使其明确化，这里有一个用BlooP写的过程定义，来显示如何测试这一性质是否存在：

\begin{longtable}[]{@{}ll@{}}
\toprule\noalign{}
\endhead
\bottomrule\noalign{}
\endlastfoot
DEFINE PROCEDURE ``GOLDBACH?''{[}N{]}: & 定义过程名和变量N \\
BLOCK 0: BEGIN & 第0块开始 \\
　CELL(0)\gebfont{⇐}2 & 把CELL(0)置成2 \\
　LOOP AT MOST N TIMES & 循环至多N次 \\
　BLOCK 1: BEGIN & 第1块开始 \\
　　IF\{PRIME?{[}CELL(0){]}AND RPIME?{[}MINUS{[}N, CELL(0){]}{]}\}THEN: & 若CELL(0)和N-CELL(0)均为素数，则： \\
　　BLOCK 2: BEGIN & 第2块开始 \\
　　　OUTPUT \gebfont{⇐}YES & 把OUTPUT置成YES \\
　　　QUIT BLOCK 0 & 退出第0块 \\
　　BLOCK 2: END & 第2块结束 \\
　　CELL(0)\gebfont{⇐}CELL(0)+1; & CELL(0)加1 \\
　BLOCK 1: END & 第1块结束 \\
BLOCK 0: END & 第0块结束 \\
\end{longtable}

和通常一样，我们先假设NO，直到证明YES为止。而且我们是靠一股蛮劲在其和为N的数对之中进行搜索。如果两个数均为素数，我们就退出最外层的块，否则我们就返回去再试，直至穷尽全部可能性为止。

（警告：哥德巴赫性质是原始递归的，但这一事实并没有使``是否所有的数都具有哥德巴赫性质？''成为一个简单的问题------远非如此！）

推荐的练习题

你能写一个类似的过程来测试乌龟性质（或阿基里斯性质）是否存在吗？如果能，就请写出来。如果不能，是仅仅因为你不知道上界，还是因为在BlooP中刻画这样的算法时遇到了根本性的障碍？如果对在对话中提出的``妙极性''性质提同样的问题，会得出什么答案？

下面我将列出一些函数和性质，你最好花一些时间来确定你是否相信它们是原始递归的（即可用BlooP编程序）。这意味着你必须认真考虑它们所要求的计算中包含哪些操作，以及是否能给出其中涉及的全部循环的项。

FACTORIAL{[}N{]}=N!（N的阶乘），例如：

\begin{itemize}
\tightlist
\item
  FACTORIAL{[}4{]}=24
\end{itemize}

REMAINDER{[}M, N{]}=用N除M所得的余数，例如：

\begin{itemize}
\tightlist
\item
  REMAINDER{[}24, 7{]}=3
\end{itemize}

PI-DIGIT{[}N{]}=π在小数点后的第N位数字，例如：

\begin{itemize}
\tightlist
\item
  PI-DIGIT{[}1{]}=1
\item
  PI-DIGIT{[}2{]}=4
\item
  PI-DIGIT{[}1000000{]}=1
\end{itemize}

FIBO{[}N{]}=斐波那契数列的第N项，例如：

\begin{itemize}
\tightlist
\item
  FIBO{[}9{]}=34
\end{itemize}

PRIME-BEYOND{[}N{]}=大于N的最小素数，例如：

\begin{itemize}
\tightlist
\item
  PRIME-BEYOND{[}33{]}=37
\end{itemize}

PERFECT{[}N{]}=第N个``完全''数（一个数的所有约数之和等于它本身，如28=1+2+4+7+14），例如：

\begin{itemize}
\tightlist
\item
  PERFECT{[}2{]}=28
\end{itemize}

PRIME?{[}N{]}=YES若N是素数，否则为NO。

PERFECT?{[}N{]}=YES若N是完全数，否则为NO。

TRIVIAL?{[}A, B, C, N{]}=YES若A\textsuperscript{N}+B\textsuperscript{N}=C\textsuperscript{N} 成立，否则为NO。例如：

\begin{itemize}
\tightlist
\item
  TRIVIAL?{[}3, 4, 5, 2{]}=YES
\item
  TRIVIAL?{[}3, 4, 5, 3{]}=NO
\end{itemize}

PIERRE?{[}A, B, C{]}=YES若A\textsuperscript{N}+B\textsuperscript{N}=C\textsuperscript{N} 可以被某个大于1的N值所满足，否则为NO。例如：

\begin{itemize}
\tightlist
\item
  PIERRE?{[}3, 4, 5{]}=YES
\item
  PIERRE?{[}1, 2, 3{]}=NO
\end{itemize}

FERMAT?{[}N{]}=YES若A\textsuperscript{N}+B\textsuperscript{N}=C\textsuperscript{N} 可以被某一组正数A、B、C 所满足，否则为NO。例如：

\begin{itemize}
\tightlist
\item
  FERMAT?{[}2{]}=YES
\end{itemize}

TORTOISE-PAIR?{[}M, N{]}=YES若M和M+N均为素数，否则为NO。例如：

\begin{itemize}
\tightlist
\item
  TORTOISE-PAIR?{[}5, 1742{]}=YES
\item
  TORTOISE-PAIR?{[}5, 100{]}=NO
\end{itemize}

TORTOISE?{[}N{]}=YES若N为两个素数之差，否则为NO。例如：

\begin{itemize}
\tightlist
\item
  TORTOISE?{[}1742{]}=YES
\item
  TORTOISE?{[}7{]}=NO
\end{itemize}

WJU-WELL-FORMED?{[}N{]}=YES若N可以看作WJU系统中的一个良构串，否则为NO。例如：

\begin{itemize}
\tightlist
\item
  WJU-WELL-FORMED?{[}310{]}=YES
\item
  WJU-WELL-FORMED?{[}415{]}=NO
\end{itemize}

WJU-PROOF-PAIR?{[}M, N{]}=YES若在WJU系统中可以把串序列M看作串N的一个推导，否则为NO。例如：

\begin{itemize}
\tightlist
\item
  WJU-PROOF-PAIR?{[}3131131111301, 301{]}=YES
\item
  WJU-PROOF-PAIR?{[}311130, 30{]}=NO
\end{itemize}

WJU-THEOREM?{[}N{]}=YES若N对应的WJU系统中的串是一个定理，否则为NO。例如：

\begin{itemize}
\tightlist
\item
  WJU-THEOREM?{[}311{]}=YES
\item
  WJU-THEOREM?{[}30{]}=NO
\item
  WJU-THEOREM?{[}701{]}=NO
\end{itemize}

TNT-THEOREM?{[}N{]}=YES若N对应的TNT串是个定理，例如：

\begin{itemize}
\tightlist
\item
  TNT-THEOREM?{[}666111666{]}=YES
\item
  TNT-THEOREM?{[}223666111666{]}=NO
\item
  TNT-THEOREM?{[}7014{]}=NO
\end{itemize}

FALSE?{[}N{]}=YES若N对应的TNT串是数论中的一个假命题，否则为NO。例如：

\begin{itemize}
\tightlist
\item
  FALSE?{[}666111666{]}=NO
\item
  FALSE?{[}223666111666{]}=YES
\item
  FALSE?{[}7014{]}=NO
\end{itemize}

后面七个题目与我们将要进行的元数学探索关系密切，因此很值得你认真分析。

可表示性和可体现性

在继续讨论一些关于BlooP及其``亲戚''FlooP的问题之前，让我们先回顾一下当初引进BlooP的原因，并把它和TNT联系起来。我在前面说过，一旦全部原始递归的概念对一个形式系统来说都是可体现的，那么这个系统就达到了使用哥德尔方法所需的``临界质量''。这到底意味着什么呢？首先，我们必须区别这样两个概念：可表示性和可体现性。``表示''一个谓词只是从自然语言到严格的形式化表述的翻译问题。这和该谓词是不是定理没有关系。而``体现''一个谓词就是强得多的概念了。这意味着：

⑴．该谓词的全部为真的例均为定理；

⑵．全部为假的例均为非定理。

``例''在这里是指用数值取代谓词中所有自由变量后所得到的串。例如，谓词m+n=k在pq系统中是可体现的，因为该谓词的每个为真的例都是定理，每个为假的例都是非定理。这样，任何具体的加法算式，无论真假，都能被译成pq系统中的``可判定串''。但是，pq系统不能表示------更不要说体现------自然数的任何别的性质。因此在对数论进行形式化的竞争之中，它的确不是一个强有力的候选。

而TNT的优点是能够表示数论中的任何谓词。例如，很容易写一个TNT串来表示谓词``b具有乌龟性质''。因此，根据表示能力来看，TNT满足我们的全部要求。

但是，问``在TNT中哪些性质是可体现的？''，这恰恰就是问``TNT作为一个公理系统有多强？''。是不是所有可能的性质在TNT中都是可体现的？若是如此，那么TNT就能回答数论中的任何问题，它就是完全的。

原始递归谓词在TNT中是可体现的

虽然完全性将被发现是个幻梦，TNT至少相对于原始递归谓词来说还是完全的。换句话说，如果数论中的某个陈述的真假能用计算机在一个长度可预知的时间段中判定，那它在TNT中一定也是可判定的。也可以把这个问题表述成如下形式

如果能为自然数的某个性质写出一个BlooP测试，那么这个性质在TNT中是可体现的。

存在非原始递归的函数吗？

用BlooP测试可以检查许多种不同的性质，包括一个数是否为素数或完全数、是否具有哥德巴赫性质、是否是2的幂，如此等等。这自然地会使人考虑是不是数的每种性质都能被某个适当的BlooP程序所检查。虽然我们现在还无法测试一个数是否是妙极的，但我们不必过于因此心烦意乱，因为这也许仅仅是由于我们对妙极性还缺乏认识，随着进一步的挖掘，我们也许会为其中包含的循环上界发现一个通用的公式。那时就可以马上为妙极性写出一个BlooP测试了。对乌龟性质可能也是同样的。

因此，这个问题实质上是``是否总能给出运算长度的上界------还是说在自然数系统中存在一种内在的混乱，致使有时无法事先预测运算的长度？''令人震惊的是，恰好是后一种情况出现了，其原因我们不久就会看到。这种事情准会把毕达哥拉斯------他第一个证明了二的平方根是无理数------气昏了头。在我们的论证中，将使用享有盛名的``对角线法''，这是由集合论的奠基人康托尔发现的。

B库、索引编号和蓝程序

让我们首先设想一个新奇的概念：装着所有可能的BlooP程序的库。不用说，这个库------``B库''是个无穷大的库。我们要考虑由B库通过三个连续的过滤操作所得到的一个子库。第一个过滤器为我们选出其中的``待调用''程序。然后我们再去掉这个子库中的全部``测试''，只留下``函数''（顺便说明一下，在待调用程序中，链中的最后一个过程决定了整个程序应当被看成是测试还是函数）。第三个过滤器将选出那些恰有一个输入参数的函数（同样只是指链中最后一个过程）。剩下来的是一个包含所有这样的待调用BlooP程序的完整库：这些程序计算那些恰有一个输入参数的函数。

让我们称这些特殊的BlooP程序为``蓝程序''（Bloo的发音在英语里和``蓝''{[}blue{]}的发音很相近，而P是``程序''{[}program{]}的首字母）。现在我们要做的事是为每个蓝程序指定一个互不混淆的``索引编号''。怎样做呢？最简单的办法------我们正是要这样做------就是根据长度把它们排列起来：可能的蓝程序中最短的为1号，次短的为2号，等等。当然，会有许多长度相同的程序。为解决这个问题，我们使用字典次序。在这里，``字典次序''是在拓广的意义下使用的，其中字母表里包括BlooP中的所有附加字符，以某一次序排列起来，例如下面这样：

ABCDEFGHIJKLMN

OPQRSTUVWXYZ+*

0123456789\gebfont{⇐}=＜＞

(){[}{]}\{\}-\textquotesingle?:;,.

------注意：其中最后一个是个空格！这样，一共有五十六个字符。为方便起见，我们可以把所有长度为1的蓝程序放在第一列，两个字符组成的程序放在第2列，等等。显然前几列完全是空的，而后面一些列中有许多许多项（虽然每一列中的项数是有限的）。排在第一的蓝程序大概是这样一个：

\begin{longtable}[]{@{}ll@{}}
\toprule\noalign{}
\endhead
\bottomrule\noalign{}
\endlastfoot
DEFINE PROCEDURE ``A''{[}B{]}: & 定义过程名和参数B \\
BLOCK 0: BEGIN & 第0块开始 \\
BLOCK 0: END. & 第0块结束 \\
\end{longtable}

这个程序实在有点傻：不管输入是什么，输出值总是0。它出现在第56列中，因为它有56个字符（请数一下必要的空格，包括那些隔开相邻两行的空格）。

在遇到第56列之后不久，后面的列变得越来越满，因为把符号组合起来形成蓝程序的方法实在太多了。不过不要紧------我们不想设法把这个无穷无尽的目录表列出来。抽象地说，我们所关心的是，这个目录表应当是良定义的，即每个蓝程序能从中得到一个唯一确定的索引编号。这才是关键所在。

让我们这样指示被第k个蓝程序所计算的函数：

蓝程序\{\#k\}{[}N{]}

在这里，k是该程序的索引编号，N是唯一的输入参数。例如，蓝程序\#12可以是送回其输入值的二倍的那个函数：

蓝程序\{\#12\}{[}N{]}=2×N

上述等式的意义是：左面指定的程序将给出一个返回值，这个值和人从右面的常规代数公式中算出的相等。作为另一个例子，或许第5000个蓝程序计算其输入参数的立方：

蓝程序\{\#5000\}{[}N{]}=N\textsuperscript{3}

对角线法

好，现在我们就要耍那个``花招''了，即使用康托尔的对角线法。我们将用这个蓝程序的目录表定义一个单变量函数------``蓝对角''{[}N{]}------而且将会发现这个函数根本不在表中（这就是把它的名字加上引号的原因）。而``蓝对角''很显然是个良定义的单变量可计算函数，这样我们就不得不下结论说，存在一些不能在BlooP中编程序的函数。

下面就是``蓝对角''{[}N{]}的定义：

等式（1）：``蓝对角''{[}N{]}=1+蓝程序\{\#N\}{[}N{]}

我们的策略是：把每台``绞肉机''的索引编号送给它自己去加工，然后把它的输出值加1。为了说清楚这一点，让我们找出``蓝对角''{[}12{]}。我们已经看到蓝程序\{\#12\}就是函数2N，因此，``蓝对角''{[}12{]}的值一定就是1+2×12，也就是25。类似地，``蓝对角''{[}5000{]}应该等于125000000001，因为这个数是1加上5000的立方。同理，你可以为任意选定的自变量找出``蓝对角''的值。

``蓝对角''的独特之处在于：它不出现在全部蓝程序所构成的目录表中。它不可能在那里。理由是这样的：如果它是个蓝程序，它就一定得有个索引编号------比如说它是蓝程序\#X。这个假设可以被表示成

等式（2）：``蓝对角''{[}N{]}=蓝程序\{\#X\}{[}N{]}

但等式（1）和（2）是互不相容的。当我们要去计算``蓝对角''{[}X{]}的值时，这就会变得很显然。因为我们在计算时要令二式中的N均取值为X。如果在公式（1）中做此代换，我们得到：

``蓝对角''{[}X{]}=1+蓝程序\{\#X\}{[}X{]}

但如果对公式（2）做此代换，我们得到：

``蓝对角''{[}X{]}=蓝程序\{\#X\}{[}X{]}

而``蓝对角''{[}X{]}不可能既等于某个数又等于该数加一。但这正是上述两个公式的意思。因此我们只得返回去丢掉某个造成这个矛盾的假设。唯一可能的选择是丢掉公式（2）所表示的假设：函数``蓝对角''{[}N{]}可以写成一个蓝BlooP程序。而这就证明了``蓝对角处于原始递归函数的范围之外''。这样，我们就达到了目的，一举摧毁了阿基里斯所珍爱的那个朴素观念：每个数论中的函数都必定能在可预测的步数内被算出来。

这里出现了一些很微妙的事情。比如说你可能会考虑这个问题：对每个特定的N值来说，在计算``蓝对角''{[}N{]}时所包括的步数总是可预测的------但这些不同的预测方式不能全都结合成一个普适的方法，以预测``蓝对角''{[}N{]}的运算长度。这是一个``无穷阴谋''，它联系于乌龟那个``无穷巧合''的概念，以及ω不完全性。这里我们不再仔细追寻这些联系。

康托尔本人给出的对角线论证

为什么这种方法叫作``对角线''论证呢？这个术语来自康托尔本人的对角线论证，而许多其它论证都是后来从其中派生的。为了讲解康托尔当时所创造的这一论证，我们要稍微离开正题，不过这样做还是值得的。康托尔也是致力于表明某个项不存在于一个特定的表之中。具体点说，康托尔要表明的是：如果建立了一个实数``登记表''，它就不可避免地要把某些实数拒之门外------因此，实际上``一个完全的实数登记表''这个观念本身就是自相矛盾的。

必须看到，这不仅适用于有限的登记表，而且也适用于无穷长的登记表。这个结果比下述命题深刻得多：``实数有无穷多个，因此显然不能列在一个有限长的登记表中''。康托尔的结果的实质在于说明了存在（至少）两种不同类型的无穷：一种无穷描述了在一个无穷长的登记表中能有多少表目，另一种无穷描述了实数的数量（即一条线或线段上点的数量）------因此后者``更大''，其理由是实数不能被塞进一个具有前一种无穷的长度的表之中。让我们看看为什么康托尔的论证中涉及了对角线这个词。

我们只考虑0和1之间的实数就够了。为了进行论证，先假设能够给出一个无穷长的表，其中每个正整数N对应于0和1之间的一个实数r\textsubscript{(N)}，而且0和1之间的每个实数都将出现在表中的某个位置上。由于实数都能写成无穷小数，我们可以设想这张表的开始部分是这样的：

r\textsubscript{⑴}：.\heititext{1} 4 1 5 9 2 6 5 3\ldots\ldots{}

r\textsubscript{⑵}：.3 \heititext{3} 3 3 3 3 3 3 3\ldots\ldots{}

r\textsubscript{⑶}：.7 1 \heititext{8} 2 8 1 8 2 8\ldots\ldots{}

r\textsubscript{⑷}：.4 1 4 \heititext{2} 1 3 5 6 2\ldots\ldots{}

r\textsubscript{⑸}：.5 0 0 0 \heititext{0} 0 0 0 0\ldots\ldots{}

其中对角线上的数字用黑体印出了：1、3、8、2、0\ldots\ldots 现在要用这些数字来构造一个特殊的实数d，d在0和1之间，但我们将发现它不在表内。为了构造d，我们依次取出对角线上的数字，然后把其中的每一个都改成别的什么数字。当在这样得到的数字序列前加上小数点之后，你就得到了d。当然把一个数字改成别的什么数字会有许多种办法，相应地可以得到许多不同的d。例如，如果我们从每个对角线上的数字中减去1（并约定0减1得9），那么我们得到的d将是：

.0 2 7 1 9\ldots\ldots{}

这时，根据我们的构造方式，有：

d的第1位不同于r\textsubscript{⑴}的第1位；

d的第2位不同于r\textsubscript{⑵}的第2位；

d的第3位不同于r\textsubscript{⑶}的第3位

\ldots\ldots 如此等等。

因此，我们有：

d不同于r\textsubscript{⑴}；

d不同于r\textsubscript{⑵}；

d不同于r\textsubscript{⑶}；

\ldots\ldots 如此等等。

换句括说，d不在表中！

对角线论证证明了什么？

现在来看康托尔的证明和我们的证明之间的根本区别------这涉及到对哪个假设进行否定。在康托尔的论证中，不牢靠的假设是可以写出上述的表。因此，d的构造就保证了下述结论的正确性：根本无法写出囊括了一切实数的表------这等于说整数的集合还没有大到足以为实数的集合作索引。而另一方面，在我们的证明中，我们知道蓝BlooP程序的登记表可以被写出来------整数集已经是够大，能够为蓝BlooP程序作索引。这样，我们只得退回去，撤消某个以前用过的更不牢靠的想法，而这个想法就是``蓝对角''{[}N{]}可以用某个用BlooP写的程序来计算。这是采用对角线法时的一个微妙差别。

\pandocbounded{\includegraphics[keepaspectratio]{./media/OEBPS/Images/img_073.jpg}}

图73．盖奥尔格·康托尔

如果把这种方法用于对话中提出的``伟大数学家的完备名单''这个更具体的例子，这一差别就会变得更为清楚。对角线本身是``Dboups''。如果我们完成所要求的对角线减法，就能从中得到``Cantor''{[}康托尔{]}。这样就有两种可能的结论。如果你坚信这张表是完全的，那你必须下结论说康托尔不是伟大的数学家，因为他的名字和表中所有的名字均不同。另一方面，如果你坚信康托尔的确是个伟大的数学家，那你必须下结论说这张``伟大数学家的完备名单''是不完全的，因为康托尔的名字没有列在表上！（那些对两者都坚信的人将多么不幸！）前一种情况相应于我们的证明：``蓝对角''{[}N{]}不是原始递归的；后一种情况相应于康托尔的证明：实数表是不完全的。

康托尔的证明按字面上的意义使用了一条对角线。而其它的``对角线''证明则基于一个更加一般的观念，这个观念是从该词的几何意义中抽象出来的。对角线法的本质在于以两种不同的方式使用同一个整数------或者可以说在两个不同层次上使用同一个整数------多亏了这样人们才可能在某个预先给定的序列外面造出一个项来。这个整数一会儿被用作一个纵向索引，另一会儿被用作一个横向索引。这在康托尔的构造中表现得很清楚。至于说函数``蓝对角''{[}N{]}，它涉及到在两个不同层次上使用同一个整数------首先作为一个蓝程序索引编号，然后再作为一个输入参数。

对角线论证阴险的重复

初看起来，康托尔的论证似乎不足以完全使人信服。难道没有办法摆脱它吗？或许把从对角线上构造出来的数d扔进去，就可能得到一张完全的表。如果你在打这个主意，那你将看到把数d扔进去根本于事无补，因为一旦你为它在表中指定了一个位置，对这张新表又可以使用对角线法，能构造出一个不在这张新表中的数d\textquotesingle。无论你重复多少次下列操作：用对角线法构造一个数，然后把它扔进表中以形成一个``更完全''的表，你仍然无法逃出康托尔方法的巨掌。你甚至可以试着去构造一张实数表，以某种方式提防着康托尔对角线法的全套明枪暗箭以及它的阴险重复，设法战胜它。这是一个有趣的练习。但如果你动手试试看，就会发现无论你怎样腾挪躲闪，还是跑不出康托尔的手心。可以说任何自称构成了``包括全部实数的表''之举，都是搬石头砸自己的脚。

康托尔对角线法的可重复性有点像乌龟那种毒辣手法的可重复性。乌龟用这个办法一个接一个地破坏了螃蟹那些``保真度''越来越高------至少螃蟹希望如此------越来越``完备''的唱机。这个办法就是为每台唱机创作出一支特定的歌，使它不能用那台唱机来播放。康托尔的计谋和乌龟的计谋都具有这种奇妙的可重复性，这可不是巧合。事实上，《对位藏头诗》本来就应当被命名为《对角藏头诗》。正像乌龟微妙地暗示给天真的阿基里斯的那样，发生在《对位藏头诗》中的事件，正是哥德尔用来证明\heititext{不完全性定理}的那个构造的一种形象化解释。由此可见，哥德尔的构造也很像一个对角线方法。这一点在下面两章中将变得很清楚。

从BlooP到FlooP

至此，我们已经用以BlooP语言写的程序为工具，定义了自然数上的原始递归函数和原始递归谓词所组成的类。我们也表明了BlooP不能囊括全部可用词语定义的自然数上的函数。我们甚至还用康托尔的对角线法构造了一个``非BlooP可编程的''函数，即``蓝对角''{[}N{]}。为什么在BlooP中不能表示``蓝对角''呢？是否能改进一下BlooP，以使得``蓝对角''成为可表示的呢？

BlooP的根本特征就是其中循环的有界性。如果我们抛掉对循环的这种要求，发明另一种语言，称其为FlooP（英文Free有``自由''的意思，我们借此表明FlooP的循环是自由的），会出现什么情况呢？FlooP除掉一点之外与BlooP完全相同：在FlooP中我们既可以用无顶的循环，也可以用有顶的循环（虽然在FlooP中写循环语句时注明其顶的唯一理由是为了好看）。这些新的循环将被称为MU-LOOP（μ循环）。这是为了遵守数理逻辑中的约定：在其中``自由''搜索（无界搜索）通常被注上一个叫``μ算子''的符号。因此，FlooP中的循环语句看上去会是这样的：

\begin{longtable}[]{@{}ll@{}}
\toprule\noalign{}
\endhead
\bottomrule\noalign{}
\endlastfoot
MU-LOOP & μ循环 \\
BLOCK n: BEGIN & 第n块开始 \\
· & · \\
· & · \\
· & · \\
BLOCK a: END; & 第n块结束 \\
\end{longtable}

这种特性使得我们可以用FlooP为``妙极性质''或``乌龟性质''这样的性质写测试程序------而这样的测试程序我们用BlooP是无法写出的，因为其中的搜索或许会无穷无尽。我将把对妙极性的测试程序留给有兴趣的读者自己去写，在这个测试中要实现下列功能：

⑴．如果其输入N是妙极的，则程序结束，给出答案YES。

⑵．若N是非妙极的，但导致了一个不同于1-4-2-1-4-2-1\ldots\ldots 的封闭循环，则程序停止，给出答案NO。

⑶．若N是非妙极的，而且导致一个``无穷上升过程''，程序将不会结束。以这种方式，FlooP的不回答就是回答。FlooP的不回答就像赵州用``无''来``废问''一样。

在第3种情况下，具有讽刺意味的是：输出变量OUTPUT的值一直是NO，但我们又一直无法得到它，因为程序仍在不停地运转。这个讨厌的第三种可能性是我们写自由循环时所必须付出的代价。在所有包含了MU-LOOP的FlooP程序中，无终止总是一种理论上的可能性。当然，有许多FlooP程序实际上相对于其所有可能的输入值都是能终止的。例如像我前面提到过的，大多数研究过妙极性的人都认为像上面提出的那样一个FlooP程序将总是能终止的，而且每次的答案都是YES。

有终止和无终止的FlooP程序

一个看起来很吸引人的想法是：最好能把FlooP过程分成两类：``有终止过程''和``无终止过程''。一个``有终止过程''无论输入什么都将终止，尽管其中的循环具有``μ性质''。而对一个``无终止过程''来说，则输入中至少有一种选择会使它永远运行下去。如果对每个FlooP程序，我们通过某种复杂的检验总能说出它属于哪一类，这将会有某种相当显著的效果（我们很快就会看到）。不必说，这种类别检测操作本身也必须是个有终止操作------否则就什么也得不到了！

图灵的妙计

我们立刻就会想到，可以让一个BlooP过程完成这个检验。但BlooP过程只接收数字输入，不能接收程序！不过，我们可以解决这问题\ldots\ldots 只要用数字对程序进行编码就行了！这个巧妙的计策只不过是哥德尔配数法的许多表现形式中的另一种罢了。令FlooP字母表中的五十六个字符分别对应于``密码子''901、902\ldots\ldots956。这样每个FlooP程序就得到了一个相当长的哥德尔数。例如，那个最短的BlooP函数（它同时也是一个有终止的FlooP程序）------

\begin{longtable}[]{@{}l@{}}
\toprule\noalign{}
\endhead
\bottomrule\noalign{}
\endlastfoot
DEFINE PROCEDURE ``A''{[}B{]}: \\
　BLOCK 0: BEGIN \\
　BLOCK 0: END \\
\end{longtable}

------将得到一个哥德尔数，它的一部分如下所示：

\begin{longtable}[]{@{}lllllllllll@{}}
\toprule\noalign{}
\endhead
\bottomrule\noalign{}
\endlastfoot
904 & 905 & 906 & 909 & 914 & 905 & \ldots{} & 905 & 914 & 904 & 955 \\
D & E & F & I & N & E & & E & N & D & • \\
\end{longtable}

现在我们的计划是写一个名为``TERMINATOR?''的BlooP测试，其功能是：如果被其输入数字所编码的那个FlooP程序是有终止的，则回答YES，否则回答NO。用这个办法就可以把这个任务交给机器去完成。如果走运，就能把有终止过程和无终止过程分开了。但是，阿兰·图灵给出的一个精巧的论证表明，任何BlooP程序都不能一贯正确地完成这种区分工作。图灵所用的技巧实际上和哥德尔的技巧基本相同，因此也密切联系于康托尔的对角线技巧。我们将不在这里把它给出来------只需说基本想法是往``终止检验程序''中送入它自己的哥德尔数。但情况并非如此简单，因为这很像设法在一个句子内部完整地引用它自己。你必须对引号也加引号，如此进行下去。这似乎会导致一个无穷回归。但是，图灵指出了一种技巧，可以把一个程序自身的哥德尔数送给它。对这同一个问题在另一个环境下的解决将在下一章中给出。本章里，我们将沿另一条途径来达到相同的目的，即证明不可能构成一个终止测试器。如果有些读者希望看到图灵方案的一个漂亮而又简单的表示，我推荐文献目录中提到的霍尔和艾利森{[}Hoare; Allison{]}的著作。

终止测试器将是法力无边的

在我们摧毁这个观念之前，让我们描述一下为什么说掌握了终止测试器将是一件值得大书特书的事：在某种意义上说，这就像有了一支法力无边的魔杖，用它可以轻而易举地解决数论中的一切问题。例如，假设我们想知道哥德巴赫变奏是否是个真实的猜想，也就是要问，是否所有自然数都具有乌龟性质？我们将首先写一个叫``TORTOISE？''的FlooP测试，它检查它的输入是否具有乌龟性质。现在这个过程的缺点------即如果乌龟性质不成立，程序将不终止------反倒成了一个优点！因为现在我们可以把过程TORTOISE？送入终止测试器。如果它回答YES，那意味着TORTOISE？对一切输入值都将要终止------换句话说，所有自然数都有乌龟性质。如果它回答NO，那我们就知道存在一个具有阿基里斯性质的数。具有讽刺意义的是，我们根本没有实际使用TORTOISE？程序------只是对它做了一次检查！

把数论中的每个问题编成程序，然后用一个终止测试器对它们来一个``一揽子解决''，这个主意有点像前面提过的那个测试公案的纯正性的想法：把公案编码在一个折叠的串中，然后去检查这个折叠的串是否具有佛性。正如阿基里斯所提示的那样，或许我们希望得到的信息在某一种表示法中比在另一种表示法中``更接近表面''。

F库、索引编号和绿程序

好，白日梦到此收场。我们怎么能证明终止测试器是不可能造出来的呢？我们关于其不可能性的论据联系于下述思路：设法对FlooP运用对角线论证，就像我们对BlooP所做的那样。我们将会看到两种情形还是有一些微妙而又关键的差别的。

像我们对BlooP所做的那样，设想包括全部FlooP程序的库。我们叫它``F库''，然后对它进行那三种过滤操作，这样最后我们将得到：

一个包含所有这样的待调用FlooP程序的完整库：这些程序计算那些恰有一个输入参数的函数。

让我们称这些特殊的FlooP程序为``绿程序''（因为它们可能``一路绿灯''------永远畅行无阻）。

现在正像我们为所有蓝程序指定了索引编号一样，我们也能为绿程序指定索引编号，即把它们排列在一个目录表中，每一列包括具有一定长度的所有绿程序，按字典顺序排列。

到此为止，对FlooP的处理完全照搬了对BlooP的处理。现在让我们看看是否还能照搬下一步：对角线技巧。如果我们试图定义下列对角线函数会怎样呢？

``绿对角''{[}N{]}=1+绿程序\{\#N\}{[}N{]}

突然，我们碰上一块暗礁：``绿对角''{[}N{]}这个函数可能无法为所有输入值N都提供明确定义的输出值。原因很简单：我们没有从F库中剔除那些无终止程序，因此就无法保证我们能为N的所有值计算出``绿对角''{[}N{]}来。有时我们会进入永无休止的计算。因此对角线论证不能在这种情况下使用，因为它要求对角线上的函数为所有可能的输入算出输出值。

终止测试器提供给我们红程序

为了弥补这个漏洞，我们将不得不使用一个终止测试器，如果它真的存在的话。因此让我们有意识地引入这样一个可疑的假设：它真的存在，并且用它作为我们的第四个过滤器。我们逐个检查表中的绿程序，去掉那些无终止的，这样最后我们将得到：

一个包含所有这样的待调用FlooP程序的完整库：这些程序计算那些恰有一个输入参数的函数，并且对所有的输入值都能终止。

让我们称这些特殊的FlooP程序为``红程序''（因为它们早晚会碰上个红灯------都必须停下来）。现在，对角线方法可以一显身手了。我们定义

``红对角''{[}N{]}=1+红程序\{\#N\}{[}N{]}

完全和对``蓝对角''的讨论平行，我们被迫断言``红对角''{[}N{]}是个良定义的单变量可计算函数，但不在红程序的目录表内，因此甚至在强有力的FlooP语言中也是不可计算的。或许现在该转到GlooP了？

GlooP\ldots{}

是的，但GlooP是什么呢？如果说FlooP是去掉限制的BlooP，那么GlooP一定是去掉限制的FlooP。但一个限制怎么可能去掉两次呢？你怎么能构造一个比FlooP更强有力的语言呢？在``红对角''中，我们已经发现了一个函数，我们人知道如何去求它的值------求值方法已经明确地用自然语言描述出来了------但似乎不能用FlooP语言编出程序来完成这项工作。这是一个尖锐的二律背反，因为还没有人曾经发现过比FlooP更强有力的计算机语言。

有人曾经详尽地研究过计算机语言的能力，我们不必自己再去做这项工作，只须报告这样一项结果：有一大批计算机语言可以被证明是与FlooP具有完全相同的描述能力的，这就是说：任何一个可以用其中某种语言编程序完成的计算过程，必然能用所有这些语言中的任何一个来编程序完成。奇怪的是，几乎任何设计计算机语言的合理企图，都以构造出这个类中的一个成员而告终------也就是说，造出了一种能力等价于FlooP的语言。当然，通过某种努力也可以构造出一种比这一类语言要弱一些的合理并有趣的计算机语言。显然，BlooP就是弱语言的一个例子，但这是一种例外，而不是规律。关键在于存在许多很自然的方式可以用来发明算法语言，而不同的人，循着不同的途径，往往最终创造出等价的语言。它们只有形式上的差别，能力是一样的。

\ldots 是个神话

事实上，绝大多数人都相信不会再有描述计算的能力强于FlooP及其等价物的语言了。这个假说在30年代被两个人互相独立地表述出来：阿兰·图灵------关于他，后面还会进一步介绍------和阿朗佐·丘奇，二十世纪杰出的逻辑学家之一。他们的结果被称为``丘奇-图灵论题''。如果我们接受这个论题，我们就必须下结论说``GlooP''是个神话------在FlooP中已经没有限制可以取消，无法通过``解放''来增强它的能力，像我们对BlooP做过的那样。

这把我们置于一种尴尬的境地，即断言人可以为N的任意值计算``红对角''，但无法编出程序让计算机来完成这项任务。因为，如果这项任务可以得到完成，那它一定是用FlooP语言来完成的。而根据任务的构造方式，它又不可能用FlooP来完成。这个结论实在太奇特了，致使我们得非常仔细地研究它赖以建立的基础。而在这当中，你会想起来，就有我们那个不牢靠的假设，即存在一个可以区别有终止和无终止的FlooP程序的判定过程。这个判定过程的想法本来就有点可疑，因为我们已经看到它的存在将导致数论中的所有问题以一种统一的方式被解决。现在我们有了双重的理由相信任何终止测试都是一种神话------根本就无法把FlooP程序装到一台``甩干机''里，把有终止程序和无终止程序分离开来。

怀疑论者可能会坚持说这不像对此类终止测试的不存在性的一个严格证明。这种反对意见是合理的。但是，图灵的方案更为严格地论证了这一点：用一个和FlooP同类的语言，不可能写出计算机程序，来对所有FlooP程序进行终止测试。

丘奇-图灵论题

让我们简略地回顾一下丘奇-图灵论题。我们将在第十七章中相当详细地讨论它------以及它的变种，而现在只需要叙述它的几种形式，把关于其价值和意义的讨论推迟到后面去完成。这里是三种相互联系的叙述此论题的方式：

⑴．人所能计算的也就是机器所能计算的。

⑵．机器所能计算的也就是FlooP所能计算的。

⑶．人所能计算的也就是FlooP所能计算的（即一般递归或部分递归）。

术语解释：一般递归和部分递归

在本章中，我们对数论中的一些概念及其与可计算函数理论的联系进行了广泛的阐述。这是一个广阔而兴旺的领域，是计算机科学和现代数学的一个吸引人的混合体。我们在结束本章之前，还要介绍一下我们曾与之打过交道的一些观念所对应的标准术语。

前面已经提到过，``BlooP可计算''是``原始递归''的同义语。现在FlooP可计算函数可以被分成两类：（1）那些能用有终止FlooP程序计算的被称为``一般递归的''，（2）那些只能用无终止FlooP程序计算的被称为``部分递归的''（对谓词有类似的定义）。人们常常只说``递归''，其意思是``一般递归''。

TNT的能力

有趣的是，TNT是如此有力，以致于不仅能体现全部原始递归谓词，而且能体现全部一般递归谓词。我们将不证明这些事实，因为这对我们的目标来说是没有必要的。我们的目标是说明TNT是不完全的。如果TNT不能体现某些原始或一般递归谓词，那这样的不完全就太``没意思''了。------因此我们应该使它能做到这一点。然后，再说明它具有某种有趣的不完全性。

\begin{enumerate}
\item
  \phantomsection\label{Chapter13.xhtml_B_1}
\end{enumerate}

\phantomsection\label{Dialog14.xhtml}{}

\chapter{G弦上的咏叹调}

乌龟和阿基里斯刚刚参观完一家麦片厂。

\dialogname{阿基里斯：}换个话题你不介意吧？

\dialogname{乌龟：}随你便。

\dialogname{阿基里斯：}那好。你知道吗，我前几天接到了一个匿名电话。

\dialogname{乌龟：}听起来挺有意思。

\dialogname{阿基里斯：}对。嗯------问题是打电话的人语无伦次，至少现在我只能这样说。他嚷了一句什么之后就把电话挂断了------噢，不，我想起来了，他是把那句话嚷了两遍然后才挂断的。

\dialogname{乌龟：}你听清是什么话了吗？

\dialogname{阿基里斯：}嗯，整个电话是这样的：

我说：喂？

\begin{adjustwidth}{2em}{0em}

电话里\kaitismall{（大声嚷嚷着）}：放在其引文形式后面得到假句子！放在其引文形式后面得到假句子！

（咔嗒。）

\end{adjustwidth}

\dialogname{乌龟：}这样的匿名电话实在古怪。

\dialogname{阿基里斯：}我也是这么想。

\dialogname{乌龟：}也许那表面上的疯颠下面是有什么意思的。

\dialogname{阿基里斯：}也许吧。

（他们走进一所宽敞的庭院，这庭院由一些古怪的石头构造的三层楼房环绕着。院子中央有一棵棕榈树，旁边是一座塔楼。紧挨着塔楼有楼梯，楼梯上坐着一个男孩，正在同窗口里的一个年轻女人讲话。）

\dialogname{乌龟：}你要带我去哪儿，阿基？

\pandocbounded{\includegraphics[keepaspectratio]{./media/OEBPS/Images/img_074.jpg}}

图74．上和下

艾舍尔作（石版画，1947）

\dialogname{阿基里斯：}我想让你从塔楼顶上观赏一下美景。

\dialogname{乌龟：}噢，那好极了。

（他们走近了那个男孩。男孩先是好奇地瞧着他们，然后对那个年轻女人说了句什么------两人``哧哧''地笑起来。阿基里斯和乌龟并没有走上男孩坐着的楼梯，而是向左拐，沿着一段通向一座小木门的短楼梯走了下去。）

\dialogname{阿基里斯：}我们可以从这里进去。跟着我。

（阿基里斯推开了门，两人走进去，开始攀登塔内很陡的螺旋形楼梯。）

\dialogname{乌龟}作这类运动我可有点不是材料。我们还得走多久？

\dialogname{阿基里斯：}没有几段楼梯了\ldots\ldots 不过我有个想法。你干嘛非要在楼梯的正面走，你何不在反面走呢？

\dialogname{乌龟：}在反面怎么走？

\dialogname{阿基里斯：}你只消手上抓紧点，然后爬转到反面去------那面有足够的空间供你活动的。你会发现，在这些台阶的上面走和下面走是一回事\ldots\ldots{}

\dialogname{乌龟}我做得对吗？

\dialogname{阿基里斯：}就这样。

\dialogname{乌龟}喂------这个小花招真把我搞懵了，我现在应该沿着楼梯向上走还是沿着楼梯向下走？

\dialogname{阿基里斯：}同刚才的方向一样。在你那面是下楼梯，在我这面就是上楼梯。

\dialogname{乌龟：}你是不是想说我下楼梯也可以到达塔楼顶？

\dialogname{阿基里斯：}我不敢说，不过这也许是行得通的\ldots\ldots{}

（于是阿基里斯在一面往上，乌龟在另一面往下，同时沿着螺旋形的楼梯绕了起来。不久，他们都走到了楼梯的尽头。）

现在用不着这个窍门了，龟兄，这面来------我拉你翻上来。

（他伸手去拉乌龟，把他拽回楼梯的另一面。）

\dialogname{乌龟：}谢谢。从反面上来倒容易些。

（他们走上屋顶，鸟瞰全城。）

这儿真美。阿基，我很高兴你把我带上来------也许我该说带下来。

\dialogname{阿基里斯：}我料到你会喜欢的。

\dialogname{乌龟：}我一直在想那个古怪的电话。我觉得现在明白些了。

\dialogname{阿基里斯：}是吗？跟我说说好吗？

\dialogname{乌龟：}行啊。你是不是跟我一样觉得``放在其引文形式后面''这个短语里有些叫人回味的东西？

\dialogname{阿基里斯：}有点儿，是的------很有一点。

\dialogname{乌龟：}你能想象什么东西被放在其引文后面吗？

\dialogname{阿基里斯：}我觉得我能想象出毛主席步入一间宴会厅的情景，那里悬挂着一幅大横幅，上面写着他著作的引文。这样就有了站在其引文后面的毛主席了。

\dialogname{乌龟：}这倒是个富有想象力的例子。不过假定我们把``放在\ldots\ldots 后面''的意思限制在仅指一些文字在纸上的先后次序，而不是这种煞费苦心地想出来的步入宴会厅的情景。

\dialogname{阿基里斯：}好吧。不过你说的``引文''到底是什么意思呢？

\dialogname{乌龟：}当你讨论一个词或一个短语时，根据惯例，要把它放在一对引号之内，比如，我们可以说：

``哲学家''这个词有三个单字。

这里，我把``哲学家''放在引号之内，以表明我们说的是``哲学家''这个词，而不是某个有血有肉的哲学家本人。这就是所谓的``使用-谈论之别''。

\dialogname{阿基里斯：}噢？

\dialogname{乌龟：}我来解释一下，假如我对你说：

哲学家挣大钱。

那么我是在``使用''这个词，从而在你心目中制造出一个目光睿智的哲人揣着个鼓鼓囊囊的大钱包的形象。可是当我把这个词------或随便什么词------加上引号时，我就抽去了它的含义和内涵，只剩下纸上的一些符号，或者说只剩下几个音节。这就叫``谈论''。除了铅字的形状以外，这个词的其它特点都无足轻重------它可能有的任何涵义都被抽掉了。

\dialogname{阿基里斯：}这叫我想起把小提琴当苍蝇拍来使用------或者我该说``谈论''？有关小提琴的一切，除了它是固体以外，一概不重要------它所具有的任何用途、功能都被抽掉了。我看苍蝇恐怕也能如此处理。

\dialogname{乌龟：}这倒是对使用-谈论之别的一个合理推广，即使有点不太正统。不过我现在要你想想一件事物放在它自身的引文形式后面的现象。

\dialogname{阿基里斯：}那好吧。你看这个行吗？

``挺棒''挺棒。

\dialogname{乌龟：}不错，再想一个。

\dialogname{阿基里斯：}好吧。

``\,`咣当'不是我所见过的书名''\,``咣当''不是我所见过的书名。

\dialogname{乌龟：}这个例子稍加修改就能变成一个很有意思的典型。只消去掉``咣当''就行了。

\dialogname{阿基里斯：}真的吗？我来看看这会是什么。这一来就成了

``不是我所见过的书名''不是我所见过的书名。

\dialogname{乌龟：}瞧，你造了一个句子。

\dialogname{阿基里斯：}是这样。这是个有关``不是我所见过的书名''这个短语的句子，是个很糟的句子。

\dialogname{乌龟：}怎么叫糟呢？

\dialogname{阿基里斯：}因为它没什么意义。我再给你造一个：

``总是不了了之''总是不了了之。

这是什么意思？说真的，这是一种糟糕透顶的文字游戏。

\dialogname{乌龟：}我不这么想。依我看，这倒是个重要的素材。事实上，这种把一个短语放在其引文形式后面的办法极其重要，以至于我觉得该给它起个名字才好。

\dialogname{阿基里斯：}那你准备用什么名字来增加这种蠢行的尊严呢？

\dialogname{乌龟：}我想称它为``\gebfont{㧟}摁一个短语''\gebfont{㧟}摁一个短语。

\dialogname{阿基里斯：}``\gebfont{㧟}摁''？这是什么词？

\dialogname{乌龟：}要是我没数错的话，这是个双音节词。

\dialogname{阿基里斯：}我是问你为什么偏要挑这么两个字，把它们按这种顺序组合起来？

\dialogname{乌龟：}哦，这回我弄懂你问我``这是什么词''的意思了。答案是：一位名叫威拉德·范·奥尔曼·蒯恩的哲学家发明了这种办法。所以我就以他名字的谐音来命名了。但是我没法做更多的解释了。至于为什么他的名字是这么两个字------更不用说为什么要按这个顺序排列了------我无可奉告。不过我倒很愿意讲点------

\dialogname{阿基里斯：}不必麻烦了！我其实也并不想知道有关蒯恩这个名字的一切事情。不管怎么说，我现在倒是知道该如何去\gebfont{㧟}摁一个短语了。这挺好玩的。下面就是一个被\gebfont{㧟}摁了的短语：

``是个残句子''是个残句子。

它虽然挺糟，但我们仍然很喜欢它。拿一个残句子来，一\gebfont{㧟}摁，瞧，就得到一个整句子！这回还是个真句子。

\dialogname{乌龟：}\gebfont{㧟}摁一下``是个有所欠缺的国王''这一短语怎么样？

\dialogname{阿基里斯：}一个有所欠缺的国王会是------

\dialogname{乌龟：}当然是不太称职。你别往旁边岔，咱们先\gebfont{㧟}摁了再说。

\dialogname{阿基里斯：}我来\gebfont{㧟}摁这个短语，是吗？那好------

``是个有所欠缺的国王''是个有所欠缺的国王。

依我看，兴许说``国王''不如说``句子''更有意义些。得了，再给我一个。

\dialogname{乌龟：}好吧------那就再来一个吧。你试试这个：

``被\gebfont{㧟}摁时得到一首乌龟情歌''

\dialogname{阿基里斯：}这费不了什么劲儿\ldots\ldots 你听我念这次\gebfont{㧟}摁的结果：

``被\gebfont{㧟}摁时得到一首乌龟情歌''被\gebfont{㧟}摁时得到一首乌龟情歌。

嗯嗯嗯\ldots\ldots 这里有点什么蹊跷。噢，我明白了！这个句子在说它自己！你说呢？

\dialogname{乌龟：}你这是什么意思？句子可不会说话。

\dialogname{阿基里斯：}是不会说话。不过它们都要谈到点什么------而这个句子则是直截了当地------无歧义也无偏差地------谈到它自身！你就不得不转过头来回想一下\gebfont{㧟}摁到底是怎么回事了。

\dialogname{乌龟：}我看不出它在说有关它自己的什么事情。它什么时候说过``我''或``本句子''之类的话？

\dialogname{阿基里斯：}哎，你故意装傻吧？其美妙之处就在于它说了自己而又不必直接挑明！

\dialogname{乌龟：}那好吧，对我这么个笨朋友，你能不能详细讲讲？

\dialogname{阿基里斯：}嗯，你真是个满腹狐疑的乌龟\ldots\ldots 行啊，让我想想\ldots\ldots 假定我造出一个句子。就把它叫``句子J''吧，因为它里面有个空位------或者说``洞眼''。

\dialogname{乌龟：}比如说？

\dialogname{阿基里斯：}比如说

``\underline{　　}，被\gebfont{㧟}摁时，得到一支乌龟情歌''。

那么句子J的论题就依赖于我们如何填充这个空位。不过只要空位上该填的东西选好了，它也就确定了：它就是\gebfont{㧟}摁这个空位所得到的那个短语。由于它是由一个\gebfont{㧟}摁行为生成的，我们就把它叫做``句子K''。

\dialogname{乌龟：}这回懂了。如果空位里的句子是``总是被到处宣扬''，那句子K就必定是

``总是被到处宣扬''总是被到处宣扬。

\dialogname{阿基里斯：}是的，句子J宣称（至于是真是假，我并不知道）：句子K是一首乌龟情歌。无论如何，句子J在这里并没有说到它自己，而是说句子K。这些我们都是一致的吧？

\dialogname{乌龟：}不管怎么说，让我们也都同意这是一首优美的歌曲吧。

\dialogname{阿基里斯：}不过此刻我要另造一个东西来填空，那就是：

``被\gebfont{㧟}摁时得到一首乌龟情歌''

\dialogname{乌龟：}啊，天哪，你开始复杂起来了。我希望这一切别太高深得让我摸不着头脑。

\dialogname{阿基里斯：}噢，别担心------你肯定能明白。由于这样选择，句子K就变成

``被\gebfont{㧟}摁时得到一首乌龟情歌''被\gebfont{㧟}摁时得到一首乌龟情歌。

\dialogname{乌龟：}噢，你这个滑头，我明白了。这样一来句子K就恰好和句子J一样了。

\dialogname{阿基里斯：}由于句子K总是句子J的论题，这就有了一个圈，所有J就反过来指向自己。不过你看到了，自指乃是一种巧合。而通常的情形是句子J与句子K彼此完全不同，但随着对句子J中空位的恰当选择，\gebfont{㧟}摁就能给你变出这种戏法来。

\dialogname{乌龟：}噢，好机巧啊。真奇怪，我自己怎么就从来没想到过这些呢？那你来说说，下面这个句子是不是个自指的？

``由六个字组成''由六个字组成。

\dialogname{阿基里斯：}嗯嗯嗯\ldots\ldots 我恐怕说不好。你刚才给出的这个句子其实不是关于它自己的，而是关于``由六个字组成''这个短语的。当然，不管怎么说这个短语是该句子的一部分\ldots\ldots{}

\dialogname{乌龟：}那么这个句子就是在谈论自己的某个部分了------这会怎么样呢？

\dialogname{阿基里斯：}这也可以叫做自指吗？

\dialogname{乌龟：}依我看，这还远不是自指。不过你也别太为这些鬼事烦恼，以后你会有充裕的时间去进一步思考它们的。

\dialogname{阿基里斯：}我会吗？

\dialogname{乌龟：}你会的。而眼下你何不试着\gebfont{㧟}摁一下``放在其引文形式后面得到假句子''这个短语呢？

\dialogname{阿基里斯：}我明白你发现什么了------就是先前那个古怪电话。\gebfont{㧟}摁它就得到：

``放在其引文形式后面得到假句子''放在其引文形式后面得到假句子。

这正是那个电话里所说的！只是在他说话时我没弄清什么地方有引号！真是句混帐话。说这种话的人都该进监狱。

\dialogname{乌龟：}为什么？

\dialogname{阿基里斯：}它太让我难受了。它和前面的那些例子不一样，我这回弄不清它到底是真句子还是假句子。我越是使劲想就越理不清楚。我的头都晕了。我真想知道编出这种东西的人得了哪种神经病，居然夜里拿它去折磨无辜的人。

\dialogname{乌龟：}我也不清楚\ldots\ldots 得了，我们该下去了吧？

\dialogname{阿基里斯：}用不着下去------我们已经在地面上了。进去吧------进去你就明白了。

（他们走进塔楼，来到一个小木门前）

从这里就能走出去，跟我来。

\dialogname{乌龟：}你能肯定吗？我可不想从三楼上掉下去把背壳摔裂。

\dialogname{阿基里斯：}我还能骗你？

（他打开门，在他们前面坐着同一个男孩，正和同一个年轻女人说话，周围的一切也都一模一样。阿基里斯和乌龟走上看上去和他们进塔楼时走下的那段楼梯完全一样的楼梯，并且发现他们自己是站在一个看上去和他们最初所进的院子完全一样的院子中。）

谢谢你澄清了那个古怪的电话，龟兄。

\dialogname{乌龟：}也要谢谢你，阿基，这次散步十分愉快。希望我们很快还能见面。

\phantomsection\label{Chapter14.xhtml}{}

第十四章

\chapter[论TNT及有关系统中形式上不可判定的命题]{\texorpdfstring{论TNT及有关系统中形式上不可判定的命题\footnote{}}{论TNT及有关系统中形式上不可判定的命题}}

``牡蛎''里的两个想法

本章的标题套用了哥德尔1931年那篇著名论文的标题------只是把``《数学原理》''换成了``TNT''。哥德尔的文章是一篇技术性的论文，致力于他那个精密严格的证明，本章则直观得多。我在此所要强调的是属于他那个证明的核心的两个关键想法。首先，是一项深刻的发现：某些TNT符号串能解释成在谈论另一些TNT符号串，简言之，作为一种语言，TNT有能力``自省''，或者说是自我审视。这是哥德尔配数法的产物。第二个关键想法是：这种自我审视的性质可以全部集中于一个单个的符号串，于是这个符号串所注视的唯一焦点就是它自己了。本质上，这种``聚焦手法''可以上溯到康托尔的对角线方法。

依我看，谁要是有意于深刻地了解哥德尔的证明，就必须得认识到，从根本上说，这个证明就是由这两个主要想法融和而成的。单独来看，两个都是妙举，揉到一起就是天才的杰作了。不过，若是一定要让我作个选择，说哪一个想法更为深刻的话，我会毫不犹豫地选择第一个------符号演算系统的哥德尔配数想法，因为这一想法关系到什么是意义、什么是指称的整个观念。这是个远远超出了数理逻辑范围的想法。在数理逻辑中，康托尔的手法虽然有十分丰富的数学结果，却与实际生活很少关联。

第一个想法：证明对

我们现在就来着手分析证明本身。第九章曾给出过一个相当精细的概念，说明哥德尔同构是关于什么的。现在则要描述一个数学概念，它使我们能把诸如``符号串0=0是TNT的定理''这类陈述翻译成数论陈述。这得从``证明对''谈起。一个证明对是一对以某种特殊方式结合在一起的自然数。这里的想法是：

{两个自然数（分别用m和n表示）形成一个TNT证明对，当且仅当m是某个TNT推导的哥德尔数，而该推导的最末一行是哥德尔数为n的符号串。}

对于WJU系统也有相应的概念，而且，直观上先考虑这种情形要更容易一些。所以，我们离开TNT证明对一会儿，先看看WJU证明对，它们的定义是类似的：

{两个自然数m和n形成一个WJU证明对，当且仅当m是WJU系统的一个推导的哥德尔数，而这个推导的最末一行是哥德尔数为n的符号串。}

让我们看两个WJU证明对的例子。第一个，令m=3131131111301, n=301。这样一对m和n的值确实构成WJU证明对，因为m是下面的WJU推导的哥德尔数：

WJ

WJJ

WJJJJ

WUJ

其中最后一行WUJ的哥德尔数为301，也就是n。与此对照，令m=31311311130及n=30，它们为什么就不构成WJU证明对呢？为了看清答案，我们写出以这个m为码的所谓推导：

WJ

WJJ

WJJJ

WU

这个所谓的推导中有一步无效！那就是从第二行到第三行，也就是从WJJ到WJJJ。WJU系统中没有一条推导规则能产生这样一个形式步骤。相应地------这是至关重要的------也就没有一条算术的推导规则能使我们从311到3111。有了第九章的讨论，这大概是一项不足称道的考察，但却是哥德尔同构的核心所在：我们在任何形式系统中做的事情都有一个与之平行的算术处理。

在任何时候m=31311311130和n=30当然都不构成WJU证明对。不过，仅此并不意味着30不是WJU数。也许会有另一个m值能与30构成WJU证明对（早先的推理使我们知道WU不是WJU定理，所以实际上是不会有一个数能与30形成WJU证明对的）。

那么TNT证明对又怎么样呢？下面有两个类似的例子，一个仅仅是自称为TNT证明对，另一个则是真正有效的TNT证明对。你能指出谁是谁吗？（顺带说一下，有些地方编进了密码子``611''，目的是把TNT推导中相邻的两行隔开。按这个意义，``611''起了标点符号的作用。而在WJU系统中所有的行都以``3''起始，这已足够------无需另加标点了。）

\begin{longtable}[]{@{}lll@{}}
\toprule\noalign{}
\endhead
\bottomrule\noalign{}
\endlastfoot
⑴ & m= & 626, 262, 636, 223, 123, 262, 111, 666, 611, 223, 123, 666, 111, 666 \\
& n= & 123, 666, 111, 666 \\
⑵ & m= & 626, 262, 636, 223, 123, 262, 111, 666, 611, 223, 333, 262, 636, 123, 262, 111, 666 \\
& n= & 223, 333, 262, 636, 123, 262, 111, 666 \\
\end{longtable}

说清哪个是哪个是相当简单的，只要把它们翻译成老记号，再作些例行公事的测试，看看

⑴．这个以m为码的所谓推导是否真是一个合法的推导；

⑵．如果是，它的最后一行是否与那个以n为码的符号串相符？

第二步是不足道的，而第一步在下述的意义上也完全是直截了当的：不会碰到没完没了的搜索，也没有暗藏着的无休止循环。想一想上面有关WJU系统的例子，同时在心里把WJU系统的规则换成TNT的规则，把WJU系统的公理换成TNT的公理就行了。两种情形的算法是相同的，把它写明白，就是：

\begin{itemize}
\tightlist
\item
  在该推导中一行一行地看下去。
\item
  标出哪些是公理。
\item
  对不是公理的那些行，逐个检查它们是否能根据某条推理规则从前面的某些行中得到。
\item
  如果全部非公理的行都是根据推理规则从出现在前面的行中得到的，那就是一个合法的推导；否则就是一个冒牌的推导。
\item
  每一阶段要执行的任务都是明白清楚的，而这些任务的数量是很容易事先确定的。
\end{itemize}

``是证明对''是原始递归的\ldots{}

读者大约已能察觉，我之所以在此强调这些循环的有界性，是要宣布

{基本事实1：``是证明对''这一性质是原始递归的数论性质，因而可以用一个BlooP程序测试。}

这和另一个密切相关的数论性质------``是定理数''------有个值得注意的对比。要断定n是一个定理数，就是要断定有某个m值能与n形成证明对（顺带说一下，这些评注对TNT和WJU同样适用，把两者都记住也许会有好处，其中WJU作为原型）。要核对n是不是定理数，你得动手搜索所有可能作为n的证明对``搭档''的m值，从而有可能要进行一场无穷尽的追寻。没有人能说清当你想要找到一个数与n（作为第二分量）形成一个证明对时，得找多长时间。问题的关键就在于，一个系统之中既有使符号串变长的规则，也有使符号串变短的规则。两种规则的并存导致了某种程度的不可预言性。

哥德巴赫变奏的例子这时候能有所帮助。测试数对（m, n）是否形成乌龟对------即m和m+n是否都是素数------是一件不足道的事情。这之所以容易，是由于``是素数''这一性质是原始递归的，对此可进行一个事先就知道必定会有结局的检测。然而，要是想知道n是否具有乌龟性质，也就是问：``有没有什么数能与n形成乌龟对？''------那就又一次使我们陷入难以驾驭的由μ循环导致的未知之中了。

\ldots 因而可在TNT中体现

这个时候，关键的想法就是上述的基本事实1，因为我们可以从它推出

{基本事实2：``构成证明对''这一性质可以用BlooP测试，因而被一个有两个自由变元的公式在TNT中所体现。}

又一次，我们没有在意这些证明对是相对于哪个系统的，不过这实在无关紧要，因为上述两个基本事实对任何形式系统都成立。形式系统的本质就在于此：它总能用一种预先知道能有结局的方式，说清一个给定的符号串序列是否构成一个证明------而这一切都能移植到对应的算术概念之中。

证明对的威力

为了把话说具体些，不妨就假定我们要讨论WJU系统。读者可能还记得我们称为``无朋''的那个符号串，它在某个层次上的解释是：``WU是WJU系统的一个定理''。我们可以说明如何利用体现WJU证明对概念的那个公式在TNT中表示无朋。我们把体现WJU证明对的公式（其存在性由基本事实2保证）缩写为：

WJU-PROOF-PAIR\{a, a\textquotesingle\}

由于它涉及两个数的性质，所以得用含两个自由变元的公式来体现（要注意：本章中我们总是使用简朴TNT------因此得小心区分变元a，a\textquotesingle，a"）。为了断言``WU是WJU系统的定理''，我们得作一个与之同构的陈述``30是WJU系统的定理数''，然后再把它翻译成TNT记号。利用我们的缩写法，这不难做到（同时回忆一下，第八章里为了标明每个a\textquotesingle 都用某个数字去替换，我们是在该数字之后写上``/a\textquotesingle''）：

\gebfont{∃}a:WJU-PROOF-PAIR\{a, SSSSSSSSSSSSSSSSSSSSSSSSSSSSSS0/a\textquotesingle\}

数数S，共30个，应该看到这是TNT的一个闭句子，原有的自由变元一个量化了，另一个则换成了一个数字。另外，这里作了一件聪明事。基本事实2给我们指出了一种谈论证明对的办法，我们刚才又想出了该如何谈论定理数：在前面加一个存在量词就行了！上面的符号串更直接地翻译过来将是：``存在某个数a能与30（作为第二个分量）形成WJU证明对。''

假使我们要对TNT作类似的事情------比方说要表示``0=0是TNT定理''这句话，我们可以用类似的方法缩写出一个（仍是有两个自由变元的）公式（其存在性由基本事实2保证）：

TNT-PROOF-PAIR\{a, a\textquotesingle\}

（这个缩写的TNT公式的解释是：``自然数a和a\textquotesingle 构成TNT证名对''。）接下来就是按照上述无朋模式把我们的话翻译成数论陈述，就是：``存在一个自然数a，它与666, 111, 666（作为后一分量）构成TNT证明对''，表示它的TNT公式是

\begin{center}

\pandocbounded{\includegraphics[keepaspectratio]{./media/OEBPS/Images/Formula04.png}}

\end{center}

------一个TNT闭句子（我们称之为``赵州''，原因一会儿就清楚）。由此我们看到有那么一种方法，它不但可以谈论TNT证明对这样的原始递归概念，而且也能谈论与之有关而又不那么靠得住的``TNT定理数''这一概念。

为了检验你对这些想法的理解，请你想想如何把下述有关TNT的元陈述翻译成TNT记号：

⑴．0=0不是TNT定理。

⑵．\textasciitilde0=0是TNT定理。

⑶．\textasciitilde0=0不是TNT定理。

它们的解答与上述例子有什么不同？彼此之间有什么不同？这里再给你几个翻译练习：

⑷．赵州是TNT定理。（表示这句话的符号串称为``元赵州''）

⑸．元赵州是TNT定理。（表示这句话的符号串称为``元元赵州''。）

⑹．元元赵州是TNT定理。

⑺．元元元赵州是TNT定理。

（等等，等等。）

例5表明关于元元TNT的一些陈述能翻译成TNT记号；例6则表明对元元元TNT的陈述也是如此，等等。

这时候要紧的是要记住：表示一个性质和体现一个性质是不一样的，例如``是TNT定理数''由公式

\gebfont{∃}a:TNT-PROOF-PAIR\{a, a\textquotesingle\}

来表示，翻译为：``a\textquotesingle 是TNT定理数''。但我们并不担保此公式体现了这一概念，因为我们并未能担保这个性质是原始递归的------事实上我们至少会暗中怀疑它不是。（这样怀疑是合理的。``是TNT定理数''这一性质不是原始递归的，因而也就没有一个TNT公式能体现它！）相比之下，``是证明对''这个性质鉴于它的原始递归性，就既可以用上面引入的这个公式表示，也可用这个公式体现。

代入导致第二个想法

通过上面的讨论，我们已经能够看清TNT是怎样对``是否是TNT定理''这一概念进行``反省''的了。这是哥德尔证明第一部分的精华。我们现在打算着手讨论证明的第二个主要想法，方法是构想一个概念，使我们能把这种反省集中进一个单个句子。要做到这一点，我们先得看看当你用一个简单的方法从结构上修改某个公式时，该公式的哥德尔数会发生什么变化。其实，我们要考虑的只是这样一种特殊的修改：

{把全部自由变元都用具体数字替换。}

下面的对照表中，左边一列给出这种运算的一些例子，而右边一列则显示其哥德尔数的对应变化

\begin{longtable}[]{@{}ll@{}}
\toprule\noalign{}
\endhead
\bottomrule\noalign{}
\endlastfoot
公式 & 哥德尔数 \\
a=a & 262, 111, 262 \\
全部自由变元都换成2的数字表示 & ↓ \\
SS0=SS0 & 123, 123, 666, 111, 123, 123, 666 \\
\multicolumn{2}{@{}l@{}}{%
*　*　*　*　*} \\
～\gebfont{∃}a:\gebfont{∃}a\textquotesingle:a"=(SSa·SSa\textquotesingle) & 223, 333, 262, 636, 333, 262, 163, 636, 262, 163, 163, 111, 362, 123, 123, 262, 236, 123, 123, 262, 163, 323 \\
全部自由变元都换成4的数字表示 & ↓ \\
～\gebfont{∃}a:\gebfont{∃}a\textquotesingle:SSSS0=(SSa·SSa\textquotesingle) & 223, 333, 262, 636, 333, 262, 163, 636, 123, 123, 123, 123, 666, 111, 362, 123, 123, 262, 236, 123, 123, 262, 163, 323 \\
\end{longtable}

右边一列有一个与左边同构的、把一个大数变成一个更大的数的算术过程。依靠加法、乘法和10的幂运算等等，不难算术地描述从旧数到新数的函数。------不过用不着真这么做。问题的要点在于：（1）原公式的哥德尔数、（2）其数字表示被插入的那个数以及（3）结果的哥德尔数，这三者之间的关系是原始递归的。也就是说，我们可以写出一个BlooP测试，使得对任意输入的三个自然数，若它们之间有这种关系，该测试就能回答``是''，否则就会回答``否''。你可以检查下面的两组数，测验一下自己作这种测试的能力，同时弄清楚为什么这一过程不会出现隐蔽的无终止循环。

\begin{longtable}[]{@{}ll@{}}
\toprule\noalign{}
\endhead
\bottomrule\noalign{}
\endlastfoot
⑴ & 362, 262, 112, 262, 163, 323, 111, 123, 123, 123, 123, 666 \\
& 2: \\
& 362, 123, 123, 666, 112, 123, 123, 666, 323, 111, 123, 123, 123, 123, 666 \\
⑵ & 223, 362, 262, 236, 262, 323, 111, 262, 163 \\
& 1: \\
& 223, 362, 123, 666, 236, 123, 666, 323, 111, 262, 163 \\
\end{longtable}

像往常一样，一个例子是对的，另一个是错的。我们称三个数之间的这种关系为代入关系。由于这个关系是原始递归的，所以可用某个有三个自由变元的TNT公式体现。我们把这个公式缩写成

SUB\{a, a\textquotesingle, a"\}

由于这个公式体现代入关系，所以下面给出的公式就必定是TNT定理：

\begin{center}

\pandocbounded{\includegraphics[keepaspectratio]{./media/OEBPS/Images/Formula05.png}}

\end{center}

（此公式基于本节前面用并行两列的表给出的代入关系例子中的第一个。）同样也因为SUB公式体现代入关系，于是下面给出的公式自然就不是TNT定理：

SUB\{SSS0/a, SS0/a\textquotesingle, S0/a"\}

算术{㧟}摁

现在已到了决定关头：可以把拆开的各个部分组合成意义深远的整体了。我们要把TNT-PROOF-PAIR和SUB这两个公式作为部件，以某种方式把它们结合成一个单个的TNT句子，其解释为：``本TNT符号串不是TNT定理''。怎么做呢？到了现在，所需的全部零件都摆在眼前，答案却仍然不易找到。

把一个公式自身的哥德尔数代入这个公式似乎是个怪异但却意义不大的想法。这与另一个怪异但却似乎意义不大的想法颇为相似：在《G弦上的咏叹调》上作``\gebfont{㧟}摁''。\gebfont{㧟}摁显示出了一种出乎意料的价值：它表明了一种制造自指句子的新方法。第一次接触了各式各样\gebfont{㧟}摁型的自指句子之后，它就会暗暗尾随你------而一旦真的理解了这一原则，你就禁不住要赞叹它是如此的简单而又生机勃勃。\gebfont{㧟}摁的算术版本------姑且称之为``算术\gebfont{㧟}摁''------能使我们造出一个``谈论自身''的TNT句子。

我们来看一个\gebfont{㧟}摁算术的例子。我们需要一个至少有一个自由变元的公式。用下面的公式就行：

a=S0

它的哥德尔数是262, 111, 123, 666。我们将把这个数插进公式自身中------更确切地说是把这个数的数字形式插进去，结果得到

\begin{center}

\pandocbounded{\includegraphics[keepaspectratio]{./media/OEBPS/Images/Formula06.png}}

\end{center}

这个新公式断定了一句笨拙的假理------262, 111, 123, 666等于1。要是我们从符号串\textasciitilde a=S0出发来作算术\gebfont{㧟}摁，那就会做出一个真句子------这你自己就能看出。

在作算术\gebfont{㧟}摁时，你自然是在做我们早先定义的代入运算的一个特例。如果要在TNT内部谈论算术\gebfont{㧟}摁，就得使用公式

SUB\{a", a", a\textquotesingle\}

其中前两个变元是相同的。这是由于我们是在以两种不同的方式使用同一个数（康托尔对角线方法的幽灵！）。这里的数a"既是（1）原公式的哥德尔数，又是（2）插进去的数。我们给上述公式造一个缩写记号：

ARITHMOQUINE\{a", a\textquotesingle\}

用汉语讲，这个公式是说：

a\textquotesingle 是对哥德尔数为a"的公式作算术\gebfont{㧟}摁所得公式的哥德尔数。这句话又长又笨，我们引入一个简明雅致的术语来概括它，我们说：

a\textquotesingle 是a"的算术\gebfont{㧟}摁化。

这表示同样的意思。例如，262, 111, 123, 666的算术\gebfont{㧟}摁化是一个大得难以想象的数：

\begin{center}

\pandocbounded{\includegraphics[keepaspectratio]{./media/OEBPS/Images/Formula07.png}}

\end{center}

（这就是我们算术\gebfont{㧟}摁a=S0所得公式的哥德尔数。）我们能够很容易地在TNT内部谈论算术\gebfont{㧟}摁。

最后一锤

现在，要是再回顾一下《G弦上的咏叹调》，你就能看清，按蒯恩的方式达到自指的各种方法中，最复杂的一招是\gebfont{㧟}摁一个其本身就在谈论\gebfont{㧟}摁概念的句子。仅仅\gebfont{㧟}摁是不够的------你必须\gebfont{㧟}摁一个提及\gebfont{㧟}摁的句子！那好，在我们的情形，与之平行的手法必定是算术\gebfont{㧟}摁某个其本身就在谈论算术\gebfont{㧟}摁概念的公式。

我们立即来写出那个公式，并称它为G的``服''号串：

\textasciitilde{}\gebfont{∃}a:\gebfont{∃}a\textquotesingle:＜TNT-PROOF-PAIR\{a, a\textquotesingle\}∧ARITHMOQUINE\{a", a\textquotesingle\}＞

你可以明白地看出这一策略是怎样深刻地借助了算术\gebfont{㧟}摁化的。这个``服''号串当然有一个哥德尔数，我们称之为``u''。u的十进表达式的两头甚至中间的一小块可以直接读出来：

u=223, 333, 262, 636, 333, 262, 163, 636, 212, \ldots, 161, \ldots, 213

至于其余部分，就得知道``TNT-PROOF-PAIR''和``ARITHMOQUINE''写出来时究竟是什么样。这过于复杂了，而且不管怎么讲都有点离题太远了。

我们现在该做的就是\gebfont{㧟}摁这个``服''号串！这意味着把全部自由变元------这里只有一个，即a"------``踢出去''，把所有出现a"的地方换成数字u。于是得到

\begin{center}

\pandocbounded{\includegraphics[keepaspectratio]{./media/OEBPS/Images/Formula08.png}}

\end{center}

不管你信不信，这是一个哥德尔符号串，我们可以称之为``G''。于是就有两个刻不容缓必须回答的问题：

⑴．G的哥德尔数是什么？

⑵．G的解释是什么？

先说第一个。我们是怎么造的G呢？噢，是从其``服''号串入手，对之进行算术\gebfont{㧟}摁。那么根据算术\gebfont{㧟}摁化的定义，G的哥德尔数就是

u的算术\gebfont{㧟}摁化。

再看第二个。我们分成几步来把G译成汉语，一步比一步更为概括。作为粗略的第一步，我们作一个相当刻板的字面翻译：

``不存在数a和a\textquotesingle 使得：（1）它们形成TNT证明对；并且（2）a\textquotesingle 是u的算术\gebfont{㧟}摁化。''

眼下当然是有一个数，a\textquotesingle 是u的算术\gebfont{㧟}摁化------因而问题就必定出在另一个数a身上。这一考察使我们能把G的译文重新措辞而写成：

``没有一个数a能与u的算术\gebfont{㧟}摁化形成TNT证明对''。

（这一步有可能把人搞糊涂，后面还要详细解释。）你该看到发生什么事了吧？G在说：

``以u的算术\gebfont{㧟}摁化为哥德尔数的那个公式不是TNT定理。''

可是------这时已不再意外了------这个公式不是别的，就是G自己。因此，我们可以把G最终翻译成

``G不是TNT定理''。

------或者，你要是愿意，也可以说成

``我不是TNT定理''。

我们已经逐步得到了一个高层次的解释------一个TNT元语言的句子，它出自先前那个低层次的解释------一个数论句子。

TNT说``服了！''

在第九章，我们勾勒过这个惊人的构造的主要结果------TNT的不完全性。重申一下那个论证：

G是不是TNT定理呢？如果是，它说的就一定是一句真理。可实际上G在说什么呢？G就在说它自己的非定理性。于是从它是定理就能得出它不是定理------矛盾。

要是G不是定理又如何呢？这倒可以接受，此时不会导出矛盾。不过G的非定理性正是G所断定的------因而G讲的是真理，而G又不是定理，所以就存在（至少）一个不是TNT定理的真理。

现在，为了再一次解释解释那个巧妙的步骤，我要使用另一个类似的例子。取符号串

\textasciitilde{}\gebfont{∃}a:\gebfont{∃}a\textquotesingle:＜TORTOISE-PAIR\{a, a\textquotesingle\}∧TENTH-POWER\{SS0/a", a\textquotesingle\}＞

其中的两个缩写符号都代表TNT符号串，这些符号串你自己就可以写出。``TENTH-POWER\{a", a\textquotesingle\}''体现陈述``a\textquotesingle 是a"的10次幂''。逐字地译成汉语就是

``不存在数a和a\textquotesingle 满足：（1）它们形成乌龟对，（2）a\textquotesingle 是2的10次幂''。可2的10次幂显然是存在的------即1024，所以这个符号串实际上说的就是

``没有一个数a能与1024形成乌龟对''。

这又可以进一步地归结为

``1024没有乌龟性质''。

要点在于我们得到了一种代入方法：把对一个数的描述代入某个谓词，而不是用这个数的数字形式代入。这种做法靠的是另外再用一个量化的变元（a\textquotesingle）。这里，是1024这个数，其描述为``2的10次幂''；而上面所说的，是被描述为``u的算术\gebfont{㧟}摁''的那个数。

``算术\gebfont{㧟}摁时得到非定理''

我们要停下来喘口气，并回顾一下我们做了些什么。就我所知，要想比较透彻地把握住这种代入方法，最好的办法是直截了当地把它与蒯恩形式的说谎者悖论作一个比较。它们之间的对应是

\begin{longtable}[]{@{}lll@{}}
\toprule\noalign{}
\endhead
\bottomrule\noalign{}
\endlastfoot
假句子 & \gebfont{⇔} & 非定理 \\
一个短语的引文形式 & \gebfont{⇔} & 一个符号串的哥德尔数 \\
在谓语之前放上主语 & \gebfont{⇔} & 代入一开公式 \\
在谓语之前放一句引语 & \gebfont{⇔} & 把一符号串的哥德尔数代入一开公式 \\
把一个谓语自身作为引语放在该谓语之前（``\gebfont{㧟}摁化''） & \gebfont{⇔} & 把一符号串的哥德尔数代入一开公式（``算术\gebfont{㧟}摁化''） \\
\gebfont{㧟}摁时得到假句子（没有主语。仅有一个谓语） & \gebfont{⇔} & G的``服''号串（TNT的一个开公式） \\
``\gebfont{㧟}摁时得到假句子''（上面的那个谓语，加上引号） & \gebfont{⇔} & 数u（上述开公式的哥德尔数） \\
``\gebfont{㧟}摁时得到假句子''\gebfont{㧟}摁时得到假句子（\gebfont{㧟}摁上面那个谓语而形成的一个完整句子） & \gebfont{⇔} & G本身（把u代入``服''号串所形成的句子，即``服''号串的算术\gebfont{㧟}摁） \\
\end{longtable}

哥德尔第二定理

由于G的解释是真，所以其否定～G的解释就是假。可我们又知道TNT中推不出假的陈述。于是，不论是G还是其否定～G都不会是TNT定理。这样，在我们的系统中就找到了一个``漏洞''------一个不可判定命题。这会派生出不少结果。从G的不可判定性就可以得到一个古怪的事实：G和～G都不是定理，但公式却是定理，因为命题演算的规则保证全部形如＜P∨～P＞的良构公式都是定理。

这是一个简单的例子，其中，一个系统内的断言和一个谈论该系统的断言似乎很不合拍。这会使人疑惑不定：这个系统是否真的很精确地反映了自己？存在于TNT内部的那种``自我反映的元数学''与我们所用的这种元数学是否非常相符？这是哥德尔写论文时很着迷的问题之一。具体地说，他对是否可能在这种``自我反映的元数学''内证明TNT的一致性颇感兴趣。读者还该记得，如何证明系统的一致性是当时哲学上的一大难题。哥德尔找到了一个用TNT公式表示``TNT一致''这句话的简单方法，然后证明了这个公式（以及表示同一思想的任何其它公式）只在一个条件下才是TNT定理，这个条件是：TNT不一致。这个反常的结果对那些乐观地期待可以找到数学无矛盾性的严格证明的人真是当头一棒。

怎样在TNT内部表示``TNT一致''这句话呢？靠的是这样一个简单事实：不一致性是说有两个公式x和\textasciitilde x------一个是另一个的否定------都是定理。于是根据命题演算，全部良构公式就都是定理。因而，要说明TNT的一致性，只需提出一个TNT的句子，说明这个句子可以被证明为不是定理，就行了。所以，表示``TNT一致''的一个办法就是说``公式\textasciitilde0=0不是TNT定理。''。这在几页以前就作为习题提出了，其翻译为

\begin{center}

\pandocbounded{\includegraphics[keepaspectratio]{./media/OEBPS/Images/Formula09.png}}

\end{center}

通过一个冗长但很容易的推理可以证明------只要TNT一致------这个用TNT符号写出的一致性誓言不是个TNT定理。所以说，当TNT表示一些事物时，它的反省能力很强，但要让它证明这些事物，其反省能力就很差劲了。如果把这一结果借过来用于人的自我认识问题，恐怕极易引起争论。

TNT是不完全的ω

那么TNT``享有''哪种不完全性呢？我们将要看到，它的不完全性是第八章中定义的那种``欧米伽''型。这意味着有一个无穷的、金字塔形的符号串族，其每个符号串都是定理，但与这个族对应的``概述串''却不是定理。不难给出这种不是定理的概述串：

\begin{center}

\pandocbounded{\includegraphics[keepaspectratio]{./media/OEBPS/Images/Formula10.png}}

\end{center}

要想理解这个符号串何以不是定理，只需注意一下它与G本身是极为相似的------事实上只用一步就能从它推得G（根据TNT的量词互变规则），所以，如果它是定理，G也是。而G不是定理，于是这个公式也就不能是定理。

我们现在需要证明与之相应的那个金字塔中的公式都是定理。写出它们是很容易的：

\begin{center}

\pandocbounded{\includegraphics[keepaspectratio]{./media/OEBPS/Images/Formula11.png}}

\end{center}

\textasciitilde{}\gebfont{∃}a\textquotesingle:＜TNT-PROOF-PAIR\{S0/a, a\textquotesingle\}∧ARITHMOQUINE\{SSS\ldots\ldots SSS0/a", a\textquotesingle\}＞

\textasciitilde{}\gebfont{∃}a\textquotesingle:＜TNT-PROOF-PAIR\{SS0/a, a\textquotesingle\}∧ARITHMOQUINE\{SSS\ldots\ldots SSS0/a", a\textquotesingle\}＞

\textasciitilde{}\gebfont{∃}a\textquotesingle:＜TNT-PROOF-PAIR\{SSS0/a, a\textquotesingle\}∧ARITHMOQUINE\{SSS\ldots\ldots SSS0/a", a\textquotesingle\}＞

\gebfont{︙}　　\gebfont{︙}　　\gebfont{︙}

这些符号串说的都是什么？逐个翻译出来就是：

``0与u的算术\gebfont{㧟}摁化不形成TNT证明对''

``1与u的算术\gebfont{㧟}摁化不形成TNT证明对''

``2与u的算术\gebfont{㧟}摁化不形成TNT证明对''

``3与u的算术\gebfont{㧟}摁化不形成TNT证明对''

\gebfont{︙}　　\gebfont{︙}　　\gebfont{︙}

这里的每个断言都是讲两个特定的自然数是否构成证明对（与之对照，G本身则是谈一个特定的自然数是否是定理数）。由于G不是定理，当然就没有一个自然数会与G的哥德尔数构成证明对。所以这个族中的每句话都是真的。事情的关键在于``是证明对''这一性质是原始递归的，从而可体现，所以上面列出的每一句话（都是真的）必定都能翻译成TNT的定理。也就是说我们这座无穷的金字塔中的每一层都是定理。这就说明了TNT为什么ω不完全。

两个不同的补洞方法

由于G的解释为真，其否定\textasciitilde G的解释就为假。而利用TNT一致性的假定，我们又知道在TNT中推不出假陈述，于是无论是G还是其否定\textasciitilde G都不是TNT定理。在我们的系统中就找到了一个漏洞------一个不可判定命题。如果我们在哲学上足够超脱，从而能认清其原因是什么的话，那就不会因此而惊慌。这表明，就像绝对几何一样，TNT能够扩充。事实上，和绝对几何一样，TNT也能朝两个不同的方向扩充。它可以朝标准方向扩充------这与绝对几何朝欧几里得的方向扩充相应；它也能朝非标准的方向扩充------这当然就与绝对几何朝非欧方向的扩充相应了。标准型的扩充将是：

把G加进去作为新公理。

这种见解似乎相当无害甚至颇为称心，因为G毕竟是在说有关自然数系统的某件真的事情。那么非标准的扩充又如何呢？如果完全平行于平行公设的情形去做，就得到；

把G的否定加进去作为一条新公理。

我们怎么能做这么荒唐而又蹩脚的事呢？其实说到底，借用济罗拉莫·萨彻利那些令人难忘的话，～G不正是``与自然数的本质相抵触''的吗？

超自然数

我希望上面的引文的讽刺意味触动你。济罗拉莫·萨彻利处理几何时的问题恰恰出在他从``什么是真什么不是真''这样一个固定见解出发，而且只打算证明他一开始就估计为真的那些内容。尽管他的方法很聪明------包括否定第五公设，从而证明以此得到的几何中的很多``荒唐''命题------但他从未意识到有可能以其它方式来想象点和线。我们现在得提防不要重蹈覆辙。我们应该尽可能公正地考察把～G作为新公理加进TNT将意味着什么。试看下列类型的新公理：

\gebfont{∃}a:(a+a)=S0

\gebfont{∃}a:Sa=0

\gebfont{∃}a:(a•a)=SS0

\gebfont{∃}a:S(a•a)=0

如果人们从未考虑过引入它们，今天的数学该是什么样呢？这里的每一条都``与原先了解的数系的本质相抵触''，每一条也都提供了对全部数概念的一个深刻而又精彩的扩充------有理数、负数、无理数、虚数。这种可能性就是\textasciitilde G试图打开我们的眼界，使我们达到的地方。以往，数概念的每次扩充都要招来一片倒彩和嘘叫，尤其是把诸如``无理数''、``虚数''这类不受欢迎的新客叫做数时，你听到的叫声更是特别响亮。遵循这个传统，我们要把\textasciitilde G所揭示的那些数叫做超自然数，以表明我们确是深感其有悖于所有合乎常情的概念。

如果要把～G扔进TNT作为第六条公理，那我们最好了解一下它到底怎样与一个带有我们刚讨论过的那个无穷金字塔共存于一个系统之中。明摆出来，～G是讲

``存在某个数与u的算术\gebfont{㧟}摁化形成TNT证明对。''

------而该金字塔的各层却依次断言：

``0不是这个数''

``1不是这个数''

``2不是这个数''

\gebfont{︙}

这真令人费解，因为看上去这是个十足的矛盾（正因为如此才称之为``ω不一致性''）。费解的根源------与几何分叉的情形极为相似------在于尽管事实上我们明白这个系统是修改过的，但还是顽固地拒绝接受对符号解释的修改，总想不重新解释任何符号就混过去------结果当然就证明这是不行的。

和解来自于把\gebfont{∃}重新解释为``存在一个广义自然数''，而不再解释成``存在一个自然数''。这样作的时候还得用相应的方法把\gebfont{∀}也重新解释过。这意味着我们对自然数以外的某些数也打开了大门，也就是引进了``超自然数''。自然数和超自然数一起构成整个广义数。

这样一来，这个表面上的矛盾立即化为乌有，因为那个金字塔仍然讲的是老话：``没有一个自然数能与u的算术\gebfont{㧟}摁化形成TNT证明对。''这族句子丝毫没讲到超自然数，因为没有关于超自然数的数字形式。而此时\textasciitilde G却是说：``存在一个广义自然数能与u的算术\gebfont{㧟}摁化产物形成TNT证明对。''把这族句子和\textasciitilde G放在一起，显然是在讲：存在一个超自然数能与u的算术\gebfont{㧟}摁化产物形成TNT证明对，如此而已------半点矛盾也没有了，在一个包括超自然数在内的解释下，TNT+\textasciitilde G是一致的系统。

由于我们现在约定了把两个量词的解释加以扩充，这就意味着包含量词的任何一条定理也都有一个扩充了的意义，例如，交换律定理

\gebfont{∀}a:\gebfont{∀}a\textquotesingle:(a+a\textquotesingle)=(a\textquotesingle+a)

现在所讲的就是对全体广义自然数加法是可交换的------换句话说，不仅自然数的加法可交换，超自然数的加法也行。类似地，说``2不是自然数的平方''的那条TNT定理------

\textasciitilde{}\gebfont{∃}a:(a•a)=SS0

------现在告诉我们2同样也不是超自然数的平方。事实上，自然数所具有的性质，只要是由TNT定理给出的，超自然数也都具有。换言之，对于自然数能形式证明的东西，对超自然数也都成立。具体讲，这意味着超自然数不是我们熟悉的诸如分数、负数、复数等等的任何一种东西，最好还是把超自然数设想成比全部自然数都大的整数------无穷大的整数。要害是：尽管TNT的定理能排斥负数、有理数、无理数、复数等等，但没有办法排斥无穷大的整数：连表示``不存在无穷量''这句话的办法都没有。

这种说法初听起来会让人觉得很怪。与G的哥德尔数形成TNT证明对的那个数，（我们称它为``I''，不过这没什么特殊理由）到底有多大？不幸的是，我们没有什么合用的语言能描述无穷大整数的大小，所以我怕是没法表述I的量值了。可话又说回来，i-------1的平方根------到底有多大？它的大小就不能靠我们熟悉的自然数的大小来想象，你不能说：``噢，i大约和14的一半或24的9/10一样大。''你只好讲：``i的平方是-1''，而且差不多也就只能讲这些。在这里倒好像用得上亚伯拉罕·林肯的一句话。当有人问他``人的两条腿该是多长''的时候，他慢条斯理地说：``够着地就行。''这差不多就是回答I的大小的办法------I恰好就是一个指明G的证明结构的数那么大------不再大，也不再小了。

当然，TNT的任何一个定理都有很多不同的推导，所以你大概要怪我对I的刻划不唯一。的确是这么回事。不过，这仍然与i-------1的平方根------的情形类似。你可以回忆一下，还有一个数的平方也是-1，那就是-i，i和-i不是同一个数，它们恰好共同具有某个性质。麻烦仅在于这还恰好就是定义它们的性质！我们只得挑出其中一个------至于是哪个，这无关紧要------把它叫做``i''。其实根本没办法区分开它们，所以，也没准我们在这数百年间一直都说的是那个错的``i''，但这完全是无所谓的。现在，与i类似，I也不是唯一定义的。你只须把I看做是很多能与u的算术\gebfont{㧟}摁化产物形成TNT证明对的荒唐数中特定的某一个就行了。

超自然定理有无穷长推导

我们还没有直接面对把\textasciitilde G扔进去作为公理引出的后果。我们说到过，但没有强调。要害在于\textasciitilde G断言G有证明，如果一个系统的公理之一断言它自己的否定有证明，这个系统还能好得了吗？我们现在注定要焦头烂额了！幸好，事情还不那么糟糕。只要我们仅限于构造有穷的证明，就决不会证明G。所以，G与\textasciitilde G之间并不会发生不幸的冲突。超自然数I也不会引来灾难。但是，我们现在必须得习惯\textasciitilde G是断言一个真理（``G有证明''），而G是断言一个假理（``G没有证明''）这样的想法。在标准数论中正好相反------不过标准数论中没有什么超自然数，要注意，TNT的超自然定理------即G------可以讲假理，但全部自然定理讲的仍都是真理。

超自然数的加法和乘法

我要来讲一个有关超自然数的极为稀奇和意外的事实，但不加证明（我也不知道证明）。这件事会使人想到量子力学中的海森堡测不准原理。它是这样：我们可以用某种简单而又自然的方法对每个超自然数用三个一组的普通整数（包括负整数）来``编号''。这样，我们原有的那个超自然数I就可以有一个编号集（9, -8, 3）, 而它的后继I+1就可以有一个编号集（9, -8, 4）。并不是只有一种方式来给超自然数编号，不同的方法有各自的优点和缺点。根据某些编号方案，给定两个超自然数的编号之后，很容易算出其和的编号三元组。而按另一些编号方案，则很容易算出两个数的积的编号三元组。然而，没有一种编号方案能够方便地既算加法又算乘法。说严格一点，如果和的编号能用递归函数算出，那么积的编号就不是递归函数；反之，如果积的编号是递归函数，和的编号就不是。因此如果超自然的小学生在学会超自然加法之后仍不知道超自然的九九表，那只好原谅他们------而且反过来也一样！你不可能同时都知道。

超自然数是有用的\ldots{}

我们还可越出超自然数数论的范围，考虑超自然的分数（两个超自然数的比）、超自然的实数，等等。其实，利用超自然实数的概念可以在一个新的立脚点上建立微积分。诸如dx、dy等等长期使数学家头疼的无穷小量，如果看成无穷大实数的倒数，就完全是合乎情理的！高等分析的某些定理依靠``非标准分析''能够更直观地证明出来。

\ldots 但它们真实吗？

初次接触非标准数论时会被它搞得晕头转向。不过非欧几何也是把人搞得晕头转向的东西。不管是在哪种情形，人们都不禁要问：``这两种对立的理论中哪个正确？哪个讲的是真理呢？''某种意义上这类问题是无法回答的（而在另一种意义上------后面要讨论------却可以有答案）。这个问题之所以没有答案，原因在于这两种对立的理论尽管使用相同的术语，却不是在谈相同的概念。所以它们只是表面上对立，正如欧氏几何和非欧几何一样。在几何中，``点''、``线''等字眼都是未定义项，它们的意义是靠用它们写出的公理系统确定的。

对于数论，情形也类似。当我们决定要把TNT形式化时，我们预先就选好了当作解释词的那些术语------例如``数''、``加''、``乘''等等。采取形式化这一步骤使我们失去了主动权，无条件地接受这些术语可能有的各式各样的被动意义。但是，我们------像萨彻利一样------对意外的事情毫无准备。我们自以为我们完全知道有关自然数真的、实在的、唯一的理论是什么。我们并未想到会有一些关于自然数的问题是TNT解决不了的，因而朝不同的方向扩充TNT就能对这些问题随心所欲地回答。这样一来，我们就没有一个基础来谈论数论``事实上''该是这一种还是那一种。就像有个人发誓说-1的平方根``事实上''存在，或者``事实上''不存在一样。

几何学的分叉现象与物理学家

可以而且大约应该提出这样一种见解来反对前述的内容：假设在现实的物理世界中有一些实验，用某种特定的几何学来解释它们，要比用别的几何学解释省事得多，那我们说这种几何学是``真''的或许就有意义。于是在那些想要使用``正确''几何学的物理学家眼里，区别``真''几何学和其它几何学就有了某种意义。可是，不能把事情看得太简单了。物理学家总要与条件的近似化和理想化打交道。比如第五章谈过的我自己的那篇博士论文，就是以关于磁场中晶体的问题的一个极为理想的情形为出发点的。显露出的数学倒是高度的漂亮和对称。尽管------毋宁说是因为------该模型的人为性，某些基本的物理特征明显地反映在图表上了。这些特征就提示着对更现实的条件下可能发生的这类事情的一些推测。可是，要是没有使我那图表能够产生的那些简化假设，就决不会有如此的见地。可以看到，在物理学中这类事情比比皆是：物理学家用一个``非现实''的情形认识了现实中隐藏很深的一些特性。所以，说物理学家想要使用的几何的某一分支是描述了``真几何''时，要特别谨慎，因为实际上物理学家总要使用种种不同的几何学：在给定的情形下挑选看起来最简单、最方便的那一种。

此外------或许这更能说明问题------物理学家并不只研究我们生活在其中的三维空间。有整个一簇的``抽象空间''，在其中都可以进行物理演算，而这些空间与我们所生活的那个物理空间有完全不同的几何性质。那么谁会说``真几何''就是天王星和海王星在其中绕着太阳转的那个空间所确定的呢？有一种``希尔伯特''空间，量子力学的波函数就在里面晃动；有一种``动量空间''，那里住着傅立叶分量；有一种``互补空间''，其中有波向量跳跃；还有一种``相空间''，其中多粒子构型迅速变幻；等等。绝对没有理由说这些空间的几何学都得一样。实际上，它们不可能一样！因此，对物理学家来说，存在不同的和``对立''的几何学是至关重要的。

数论的分叉现象与银行家

几何说得够多了，数论又怎么样呢？不同的数论之间的相互依存也是至关重要的吗？你要是去问银行职员，我猜他准会惊诧不已而不肯相信。2加2怎么能不等于4而等于别的什么东西呢？要是2加2不得4，由此导致的令人无法忍受的不确定性不是立刻会使整个世界经济坍台吗？其实，并不是这么回事。首先，非标准数论并不危及2加2等于4这个古老观念，它与普通数论的不同仅在于它讨论无穷概念的方式。毕竟，每个TNT定理在TNT的任何扩充中都是定理！所以银行家们犯不上为采用非标准数论会引起混乱而感到绝望。

不管怎么说，对旧有事实的改变怀着恐惧心理，这反映了对数学和现实世界之间关系的一种误解。只是在你作出了要采取哪类数学这样一个决定性的选择之后，数学才告诉你有关现实世界的一些问题的答案。即使有一个与标准数论相左的数论，而且它也使用符号``2''、``3''和``=''，可其中一条定理却说``2+2=3''，那也没什么理由要银行家选用这种理论！因为它对金融工作不合用。你得使自己的数学适应世界，而不是相反。比如说，我们并不把数论用于云彩组成的系统，因为整数这个概念很难适应那里的情形。一块云彩和另一块云彩跑到一起并不是两块云彩，而仍只是一块。这并不证明1加1等于1，它只证明``一''的数论概念就其全部功用而言都不能用来数云彩。

数论中的分叉现象与元数学家

总而言之，银行家、数云彩的人、以及其余的绝大多数人，都不必为超自然数的出现而烦恼：它们丝毫不会影响世间的日常观念。实际上可能遇到点麻烦的，只是那些离不开无穷实体本质的人。我们周围这类人并不很多------不过数理逻辑学家属于这一范畴。数论中有分叉现象怎么就会影响他们呢？因为，数论在逻辑中有两重身份：（1）当把数论公理化时，它是研究对象；（2）非形式地使用数论时，它是研究形式系统的不可少的工具。这又是使用-谈论的区别。这里，（1）是谈论数论，（2）是使用数论。

这样一来，尽管数论不适用于数云彩，但数学家认定它适用于研究形式系统，这就像银行家认定实数算术适用于他们的业务一样。这是在数学之外作出的裁决。这也表明，进行数学研究时的思维过程和其它领域的很多思维过程一样，含有``缠结的层次结构''，其中一个层次可以影响别的层次。层与层之间并没有明显的界限。就像关于``什么是数学''的形式主义观点使人们相信的那样。

形式主义的哲学宣称数学家只讨论抽象的符号，而且可以几乎不关心这些符号有什么现实用处、与现实有什么关系。然而这是一幅很失真的画面。这在元数学中再清楚不过了。如果是靠数的理论本身的帮助去获得有关形式系统的真实知识，那么数学家们就是心照不宣地表明，他们相信这些被叫做``自然数''的飘渺的玩艺其实就是现实世界的组成部分------并不只是想象中的片断。正是由于这个原因，我才在前面的一个括号中说，在某种意义上，对``数论的哪种说法为`真'\,''有一个答案。关键在于：数理逻辑学家必须选出数论的一种说法来寄托他们的信念，尤其是他们不能在超自然数的存在与不存在之间保持中立。因为对于元数学中的问题，两个不同的理论可能给出不同的答案。

例如``\textasciitilde G在TNT中有没有有穷推导？''，这个问题实际上没人知道答案。尽管如此，大多数数理逻辑学家还是会毫不犹豫地回答说没有。给出这个答案的直觉基础是下述事实：如果\textasciitilde G是定理，TNT就是ω不一致的，这样的话，你要是打算给TNT一个有意义的解释，你就被迫得吞下超自然数------而这对大多数人来讲是一件极不舒服的事。毕竟，我们发明TNT时并没打算或者说并不指望超自然数是它的一部分。也就是说，我们------或者我们中的多数人------都相信有可能造出一个形式化数论，它并不会强迫你相信超自然数完全像自然数一样真实。在押宝的时候，确定数学家们要把自己的信念寄托在数论的哪一条``岔路''上，靠的就是对现实世界的直观。不过这种信念可能是错的。也许人们发明的每一种一致的形式数论都蕴涵着超自然数的存在，也就是都ω不一致。这当然是个古怪的想法，但却是可以想得到的。

如果真是这样------我抱有怀疑，但却拿不出否定的论证------则G不见得不可判定。其实，也许根本就没有不可判定的TNT公式。可能会简单得只有一个不分叉的数论------其中必定包括了超自然数。这并不是数理逻辑学家所期待的事情，可又是不能彻底排斥的事情。一般说来，数理逻辑学家相信TNT------以及与之类似的系统------是ω一致的，从而在任何这样的系统中可以构造出的哥德尔符号串都是该系统内不可判定的。这意味着逻辑学家挑选它或它的否定作为新公理都可以。

希尔伯特第十问题与乌龟

我打算介绍一下哥德尔定理的扩充以结束这一章（戴维斯{[}Davis{]}与赫什{[}Hersh{]}的文章《希尔伯特第十问题》对此有更全面的论述，参看文献索引）。为此，我先得定义丢番图方程，即以常量整数为系数和指数的多项式方程。例如

a=0

以及

5x+13y+1=0

以及

5p\textsuperscript{2}+17q\textsuperscript{17}-177=0

以及

a\textsuperscript{(123, 666, 111, 666)}+b\textsuperscript{(123, 666, 111, 666)}-c\textsuperscript{(123, 666, 111, 666)}=0

都是丢番图方程。一般说来，要知道一个给定的丢番图方程有没有整数解是一件很难的事情。事实上，希尔伯特在二十世纪初的一次讲演中，就要求数学家寻找能在有穷步内确定一个给出的丢番图方程有无整数解的一般性算法。他居然丝毫不怀疑可能没有这种算法存在！

现在可以把对G的讨论简单化了。已经证明，只要你有一个足够强有力的形式数论以及它的一个哥德尔配数法，那就有一个丢番图方程等价于G。这种等价性依赖于下述事实：当我们在一个元数学的层次上作解释时，该方程断言它自己无解。反过来说就是：如果你找到了一个解，那就可以从这个解构造出该系统内的一个关于此方程无解的证明的哥德尔数！这就是乌龟在《前奏曲》中用费马方程作为他的丢番图方程所做的事情。我们高兴地得知，一旦做到了这一点，就能从空气分子中找回老巴赫演奏出的声音！

\begin{enumerate}
\item
  \phantomsection\label{Chapter14.xhtml_B_1}
\end{enumerate}

\phantomsection\label{Dialog15.xhtml}{}

\chapter{生日大合唱哇哇哇乌阿乌阿乌阿\ldots\ldots{}}

五月里的一个晴朗日子，乌龟遇见了阿基里斯，两人便一起漫步在树林里。阿基里斯穿得漂漂亮亮，摇头晃脑地哼着一段曲子，别在他衬衣上的一枚硕大的徽章上写着：``今天是我生日！''

\dialogname{乌龟：}你好啊，阿基，什么事让你这么高兴？是不是赶上你生日了？

\dialogname{阿基里斯：}对，对！是的，今天是我的生日！

\dialogname{乌龟：}我猜也是。你那徽章上写着呢。而且，要是我没听错的话，你唱的是巴赫的《生日大合唱》。那是他1727年为萨克森国王奥古斯都五十七岁生日而写的。

\dialogname{阿基里斯：}你说的对。奥古斯都的生日和我一样，所以这首《生日大合唱》就有双重意义。不过我不想告诉你我的年龄。

\dialogname{乌龟：}喔，这完全在你。可我想知道另外一件事，从你已经告诉我的这些情况里，就能得出结论说今天是你的生日，对吗？

\dialogname{阿基里斯：}对，对！完全正确。今天是我的生日。

\dialogname{乌龟：}好极了。这正和我猜的一样。我将要推定今天是你的生日，除非------

\dialogname{阿基里斯：}------除非什么？

\dialogname{乌龟：}除非那是一个轻率仓促给出的推断。我们乌龟毕竟不喜欢跳跃推理（我们压根就不喜欢跳，尤其不喜欢跳跃推理。）我深知你也喜欢逻辑思维，那么请问，根据前面那些话，从逻辑上演绎出今天是你的生日，合理不合理？

\dialogname{阿基里斯：}我相信我已经看出你提问题的规律了，龟兄。但我这回不跳着推理了，我要按字面意思对待你的问题，直截了当地回答你：是。

\dialogname{乌龟：}好！好！那就只有一件事了，我需要知道它以便完全确定今天是------

\dialogname{阿基里斯：}对，对，对，对\ldots\ldots 我明白你要问些什么，龟兄。我得让你知道，我已不像不久前讨论欧几里得证明时那么容易上当了。

\dialogname{乌龟：}哪儿的话呀，谁说过你容易上当啦？正相反，我倒认为你是逻辑思维方面的专家、演绎科学的权威、正确推理方法的泰斗\ldots\ldots 说真的，阿基，依我看，你是推理艺术的巨匠。正因为如此，我要问你：``前面那些话是否提供了足够的证据，使我不用再费脑筋就能推出今天是你的生日？''

\dialogname{阿基里斯：}你对我真是过重了，龟兄，噢，我是说------过奖了！不过我注意到了你在翻来覆去地提类似的问题------要我说，你（就跟我一样）每次都能回答``是''。

\dialogname{乌龟：}我当然能，阿基。可要是那样做就得靠运气了，而我们乌龟------不像那些瞎猫------讨厌撞见死耗子这种事，乌龟只肯作合理猜测。啊，是的------合理猜测的威力。你想象不出有多少人在作猜测的时候，没有把全部的相关因素都考虑进去。

\dialogname{阿基里斯：}依我看，我觉得这一大堆废话里只有一个相关因素，那就是我的头一句话。

\dialogname{乌龟：}嗯，的确，我承认它至少是一个该考虑的因素------不过，你不会要我无视逻辑吧？那可是一门受人尊敬的学问，作合理猜测时，逻辑永远是一个相关因素。现在，既然有一位著名的逻辑学家和我在一起，我在思考时就只按符合逻辑的方式借助刚才那个因素，并且要通过直接问你我的直觉是否正确来证实我的预感。因此，我最后还是得直截了当地问问你：``上面那些话是否可以让我丝毫不差地推出今天是你的生日？''

\dialogname{阿基里斯：}我再说一遍：是。不过坦率地说，我明确地感到你自己已经可以给出答案了，就像前面几次一样。

\dialogname{乌龟：}你也太刻薄了！我倒真想能如你所说的那么聪明呢！可我只不过是只普普通通、极为无知的乌龟，渴望考虑到全部的相关因素，所以我需要所有这些问题的答案。

\dialogname{阿基里斯：}那好，我就来一劳永逸地解决这个问题。你先前所提的全部问题，以及你今后按这种方式将提出的全部问题，答案都一样，那就是：是。

\dialogname{乌龟：}好极了！一网打尽。你用自己独创的方法，避免了所有麻烦。我把这个巧妙的发明叫做答案模式，你不会见怪吧？答案模式把第一个``是''答案、第二个``是''答案、第三个``是''答案等等，统统缠在一起。事实上，当它缠下去一直缠到头时，它该有个名称：``ω答案模式''，``ω''是最末一个希腊字母------我不必给你讲这个了吧？

\dialogname{阿基里斯：}我不在乎你叫它什么。现在我总算轻松了，你终于承认了今天是我的生日。现在我们可以谈点别的了------比如说，你打算送我什么礼物？

\dialogname{乌龟：}等会儿------别这么急。我会承认今天是你的生日，只是要有一个条件。

\dialogname{阿基里斯：}什么条件？是我不要礼物吗？

\dialogname{乌龟：}不是不是。阿基，我其实正打算招待你一顿丰盛的生日晚餐------只要你让我确信，一下子给出全部那些``是''答案（由ω答案模式给出的），我就能不绕弯子而直接推出今天是你的生日。是这么回事吧？

\dialogname{阿基里斯：}对，就是这么回事。

\dialogname{乌龟：}好，那我现在就有第ω+1个``是''答案了。有了它以后，我就可以进一步来试着接受今天是你生日这个假设------如果上面那种做法有效的话。你是不是能在这件事上给我当当顾问，阿基？

\dialogname{阿基里斯：}都是些什么呀？我已经看穿了你那个无穷把戏。你不是不满足于第ω+1个``是''答案吗？那好，我就不光给你第ω+2个``是''答案，而且还给你第ω+3个，第ω+4个，一直给下去。

\dialogname{乌龟：}你真大方，阿基。今天是你的生日，该我送你礼物而不是反过来让你送我。直说了吧，我怀疑今天是不是你的生日。现在，有了这个新的答案模式（我要叫它2ω答案模式），我想我可以推出今天是你的生日了。不过请你告诉我，阿基，2ω答案模式确实能使我完成这一飞跃吗？我是不是又漏掉了什么东西？

\dialogname{阿基里斯：}你用不着再来套我。我已经有办法结束这场无聊的游戏了。我给你一个结束全部答案模式的答案模式！也就是说，我同时给你一系列答案模式：ω，2ω，3ω，4ω，5ω，一直下去。用这个元答案模式，我跳出整个系统，一了百了，超越了你自以为套住了我的这个无聊游戏------现在完事大吉了！

\dialogname{乌龟：}嗬！阿基，我很荣幸能接受这样一个强有力的答案模式。我觉得它一定是很少见的，人的头脑竟能发明出如此庞大的东西，我看着都害怕。我来为你这一馈赠起个名字吧，你看如何？

\dialogname{阿基里斯：}随你便。

\dialogname{乌龟：}那我就叫它``ω\textsuperscript{2}答案模式''。我们马上就可以干别的了------不过你要告诉我，有了ω\textsuperscript{2}答案模式，我是否就能推导出今天是你的生日？

\dialogname{阿基里斯：}唉，真倒霉！这一连串恼人的问题还有没有个完？下一个该什么啦？

\dialogname{乌龟：}哦，ω\textsuperscript{2}答案模式之后，还有第ω\textsuperscript{2}+1个答案。然后还有第ω\textsuperscript{2}+2个，一直下去。不过你又可以把它们全打成一捆，作为ω\textsuperscript{2}+ω答案模式，然后就会有其它的捆，诸如ω\textsuperscript{2}+2ω、ω\textsuperscript{2}+3ω\ldots\ldots 最后，能到达2ω\textsuperscript{2}答案模式，这后面，又有3ω\textsuperscript{2}答案模式和4ω\textsuperscript{2}答案模式。它们后面，还有更进一步的答案模式，诸如ω\textsuperscript{3}、ω\textsuperscript{4}、ω\textsuperscript{5}等等。就是这么进行下去，颇有些共同风格呢。

\dialogname{阿基里斯：}我能想象了。我看，过一会就该到ω\textsuperscript{ω}答案模式了。

\dialogname{乌龟：}那当然。

\dialogname{阿基里斯：}然后是\gebfont{ꪣ}和\gebfont{ꪤ}对吗？

\dialogname{乌龟：}你领会的快极了，阿基。要是你不介意，我想提个建议：你干嘛不把这些答案模式全都塞进单独一个答案模式里去呢？

\dialogname{阿基里斯：}好吧，我真怀疑这样做能有什么结果。

\dialogname{乌龟：}依我看，我们至今所提出的命名惯例显然都不适于给它命名。我们随便起个名字。叫它ε\textsubscript{0}答案模式吧。

\dialogname{阿基里斯：}真讨厌！每次都给我的答案起个名字。我本希望这个答案能使你满意，可这么一来我的希望立刻又落空了。我们刚才干嘛不留着这个答案模式不起名字呢？

\dialogname{乌龟：}哪有那么容易！不起名字就没法称呼它了。此外，这个特别的答案模式是有点不可避免的，而且还相当漂亮，要是连名字都没有就太不雅观了！你过生日时总不想做些什么观之不雅的事情，对吗？或者，今天不是你的生日吧？要说生日，今天该是我的生日！

\dialogname{阿基里斯：}什么？今天是你的生日？

\dialogname{乌龟：}对，没错。不过说老实话，今天是我叔叔的生日，但那差不多是一回事。今天晚上你打算怎么请我吃美味可口的生日晚餐啊？

\dialogname{阿基里斯：}得了吧，老龟！今天是我的生日，该你作东！

\dialogname{乌龟：}啊，可你一直没说服我相信你这个说法是真的。你热衷于用那些答案、答案模式等等东拉西扯，可我想知道的是：今天是不是你的生日。你整个把我搞晕了。唉，太不幸了。不过，无论如何，你今晚要是请我吃一顿生日晚餐，我将十分高兴。

\dialogname{阿基里斯：}那好吧，我正好知道一个地方，那里有各种各样可口的汤，而我也恰好知道我们该吃哪一种\ldots\ldots{}

\phantomsection\label{Chapter15.xhtml}{}

第十五章

\chapter{跳出系统}

一个更强有力的形式系统

如果有一个善于思考的人来批评哥德尔的证明，那么他可能采取的一项行动就是检查该证明的一般性。比方说，这样的一位批评者有可能怀疑哥德尔只是机智地抓住了一个有利条件：在TNT这个具体的形式系统中隐藏有某种缺陷。如果真是这样，那或许就能搞出一个比TNT高明的形式系统，其中不再会有什么东西落入哥德尔的圈套，从而就能使哥德尔定理大减其色。在这一章里，我们就要仔细审查TNT的一些性质，它们使TNT在上一章的那些讨论面前显得是个极其脆弱的系统。

一个很自然的想法是：如果TNT的根本麻烦在于它有一个``漏洞''------换句话说，是有一个不可判定的句子，也就是G------那何不直接把这个漏洞补上呢？何不将G添加到TNT中作为第六条公理呢？当然，与其它公理相比，G是个庞然大物，于是所得到的系统------TNT+G------也就因其各公理之间的不协调而显得颇为滑稽。但不管怎么说，添加G仍不失为一条合理的建议。我们就这样做做看，希望新的系统真是个高明的系统------它不但没有超自然的东西，而且还是完全的。自然，至少在这样一个方面TNT+G要比TNT强：符号串G在新系统内不再是不可判定的了，因为它是定理。

TNT的弱点何在？本质上就在于它能表示自指陈述------具体说，就是陈述

``我在形式系统TNT中不可证''。

或者，说详细一点：

``没有一个自然数与本符号串的哥德尔数形成TNT证明对。''

有没有什么理由可以期待TNT+G不再受到哥德尔证明的打击呢？实在没有。我们的新系统和TNT具有同样的表示能力。而由于哥德尔的证明主要是依赖于形式系统的表示能力，所以要是看到新系统也俯首就范，那也不值得大惊小怪。技巧在于找一个表示下面这个陈述的符号串：

``我在形式系统TNT+G中不可证''

其实，这里没有多少技巧，只要你已弄清对TNT是怎么做的就行了。所用的都是同样的原理，只是上下文稍有改变而已（打个比方来说，就像我们把一支熟悉的曲子重唱一遍，但音调高了八度）。和前面一样，我们要找的符号串------就称之为``G\textquotesingle''吧------是由一个``服''号串做中介而构造出来的。不过不再是以表示TNT证明对的那个公式为基础，而是建立在与之类似但稍微复杂的TNT+G证明对概念之上。TNT+G证明对概念只是先前的TNT证明对概念的一个轻微的扩充。

可以想象，WJU系统也能做类似的扩充。我们已经知道WJU证明对的纯朴形式。现在，如果把WU加进去作第二条公理，我们就要讨论一个新系统------WJU+WU系统。扩充后的系统的一个推导是：

\begin{longtable}[]{@{}ll@{}}
\toprule\noalign{}
\endhead
\bottomrule\noalign{}
\endlastfoot
WU & 　公理 \\
WUU & 　规则2 \\
\end{longtable}

存在一个与之对应的WJU+WU证明对------即m=30300, n=300。当然，这个数对并不形成WJU证明对------它是WJU+WU证明对。加入额外的公理并没有从本质上使证明对的算术性质变复杂。有关证明对的那个意义深远的事实------``是证明对''是原始递归的------依然保持不变。

再用哥德尔方法

现在回过头来看TNT+G，我们会发现类似的情形。TNT+G证明对与其前身一样，是原始递归的，从而可以在TNT+G内用一个公式来体现，我们用一种显而易见的方式把它缩写成

(TNT+G)-PROOF-PAIR\{a, a\textquotesingle\}

此后，只需把每件事都重做一遍。我们制造新版的G时仍从一个``服''号串出发，同过去一样：

\textasciitilde{}\gebfont{∃}a:\gebfont{∃}a\textquotesingle＜(TNT+G)-PROOF-PAIR\{a, a\textquotesingle\}∧ARITHMOQUINE\{a", a\textquotesingle\}＞

称它的哥德尔数为u\textquotesingle。于是我们就来\gebfont{㧟}摁这个``服''号串，得到G\textquotesingle：

\begin{center}

\pandocbounded{\includegraphics[keepaspectratio]{./media/OEBPS/Images/Formula12.png}}

\end{center}

其解释为

``不存在一个数a能与u\textquotesingle，的算术\gebfont{㧟}摁化形成（TNT+G）证明对''。

说简洁一点，就是

``我在形式系统TNT+G中不可证''。

多重分叉现象

哎（打哈欠），这以后的细节就是很乏味的了。G\textquotesingle 对于TNT+G就像G对原来的TNT一样。可以看出，无论是G\textquotesingle 还是\textasciitilde G\textquotesingle，都可以加进TNT+G，从而导致数论的进一步分裂。而且，你别以为这种事只会对``好伙伴''发生，这种偷偷摸摸的手段也能用于TNT+\textasciitilde G------即添加\textasciitilde G而得到的TNT的非标准扩充。于是我们看到（图75）数论中有着多层次的分叉现象：

\pandocbounded{\includegraphics[keepaspectratio]{./media/OEBPS/Images/img_075.png}}

图75．TNT的``多重分叉现象''

TNT的每个扩充都有它自己的哥德尔句子，可以把这个句子或其否定添加进去，使得从TNT的每个扩充都生出一对进一步的扩充。这是一个直到无穷的过程。

当然，这个图还只是开始部分。我们来设想，沿着这棵倒长的树的最左边的枝走下去，每次都是把相应的哥德尔句子（而不是其否定）塞进去。这是摆脱超自然数的最好办法。加进G之后再加G\textquotesingle，然后加G"，然后G"\textquotesingle，等等，每做一次，我们就得到TNT的一个新的扩充。由于它会遭到乌龟方法------哦，请原谅，我指的是哥德尔方法------的打击，于是就会产生出一个新的符号串，其解释为

``我在形式系统X中不可证''。

自然，过上一会，这整个的过程就变成按步就班的例行公事了。真是，所有的``漏洞''都是用同一种技巧造出来的！这意味着，把它们视为印符对象时，它们全都是从一个模子铸出来的，而这又意味着用一条公理模式就足以把它们全体现出来！如果真是这样，何不一下子把全部的漏洞都堵死呢？何不一劳永逸地解除这种不完全性的麻烦呢？在TNT中添加一条模式，而不是每次加一条公理，就能完成这个任务。具体说，这个公理模式就是铸出G、G\textquotesingle、G"、G"\textquotesingle、\ldots\ldots 的模子。加上这个公理模式（称为G\textsubscript{ω}，我们就能击破``哥德尔化''方法。确实，看起来似乎是很清楚的，把G\textsubscript{ω}加进TNT是把所有的数论真理完全公理化时所需步骤的最后一步。

在《对位藏头诗》中就已经接近这一点：乌龟说螃蟹发明了``唱机欧米伽''。然而，这种装置的命运如何，读者仍会挂念。因为还没等把话讲完，精疲力竭的乌龟就决定还是回家睡觉为好（但他却是在抛出有关哥德尔定理的一个闪烁其词的介绍之后才去睡觉的）。现在，我们终于能有空搞清这些悬而未决的细节了\ldots\ldots 读了《生日大合唱哇哇哇乌阿乌阿乌阿\ldots\ldots》，读者也许已经受到了一点启发。

本质不完全性

你可能已经猜到了，即使是这么个有关TNT的异乎寻常的进展，也要遭到同样的厄运。而且，它之所以在劫难逃，从本质上讲，仍是由于同一原因。这个公理模式并不足够强，因而仍能实施哥德尔的构造过程。我们来详细说说（事实上应比我在这里说的严格得多）。如果有一种方法能从一个印符模子铸出各符号串G、G\textquotesingle、G"、G"\textquotesingle\ldots\ldots 那也就会有一种方法能从一个算术模子描述其相应的哥德尔数。而且，对一个无穷数类的这种算术刻划可以在TNT+G\textsubscript{ω}内用一个公式``OMEGA-AXIOM\{a\}''来体现，该公式的解释是``a是出自G\textsubscript{ω}的一条公理的哥德尔数''。当把a换成一个具体数字时，所得到的公式是TNT+G\textsubscript{ω}的定理，当且仅当此数字表示了出自该模式的一条公理的哥德尔数。

借助这个新公式，就可以在TNT+G\textsubscript{ω}内部体现诸如``TNT+G\textsubscript{ω}证明对''这类更为复杂的概念：

(TNT+G\textsubscript{ω})-PROOF-PAIR\{a, a\textquotesingle\}

有了这个公式，我们就可以构造一个新的``服''号串并用现在已彻底搞熟的方法来算术\gebfont{㧟}摁它，从而再得到一个不可判定的符号串，称之为``TNT+G\textsubscript{ω+1}''。此刻你可能会奇怪：``G\textsubscript{ω+1}为什么不会是公理模式G\textsubscript{ω+1}生成的一条公理呢？''答案是：G\textsubscript{ω}还没聪明到能预见它自己可以嵌入到数论内部去。

在《对位藏头诗》中，乌龟制造``不能播放的唱片''时所采取的一个关键步骤，是搞到他预谋破坏的那台唱机的一份设计图。这是为了搞清哪种震颤能击中它的弱点，然后再把这种震颤录进他的唱片，使音槽记有能导致这种震颤的声音。这与哥德尔的手段极为相似。按照哥德尔的手段，一个系统本身的一些性质在证明对概念之内得到了反映，然后再针对它来使用这些性质。任何一个系统，不管多复杂，多不好把握，都能进行哥德尔配数，因而就能定义证明对的概念------以子之矛，攻子之盾。只要一个系统是良定义的，即``理顺了''，它就变得脆弱了。

对于0，1之间的每个良定义的实数列，都能利用康托尔对角线手段找到一个遗漏的实数，这种对角线手段可以贴切地说明上面的原则。正是``明晰地排列''这一举动------实数们``梳理详毕''------导致了垮台。我们来看看康托尔对角线手段怎么就能一遍又一遍地反复使用。从某个数列L出发，做下面的事情，看看会怎么样：

（1a）对数列L，构造出它的对角线数d。

（1b）找个位置把d插进数列L，得到新的数列L+d。

（2a）对数列L+d，构造它的对角线数d\textquotesingle。

（2b）找个位置把d\textquotesingle 插进数列L+d，得到新的数列L+d+d\textquotesingle{}

看来，这种逐步的过程是弥补L的笨办法，因为在当初给定L之后，我们可以一下子作出全部的d、d\textquotesingle、d"、d\textquotesingle"\ldots\ldots 可是，你要是以为作出这个数列就能补全你的实数列，那就错了。当你问：``把这个对角线数列并到L内的什么地方？''这时候麻烦就出来了。无论你设计出何等聪明的方案在L内安放这些对角线数，一旦你做好了，得到的新数列就仍有弱点。如上所说，正是``明晰地排列''这一举动------实数们``梳理完毕''------导致了垮台。

对形式系统的来说，正是给出一个明确的处方这一举动------原以为它刻划了数论真理------导致了不完全性。这就是TNT+G\textsubscript{ω}的症结所在。一旦你把全部的G都以一种界说良好的方法加入TNT，那就会看到还有其它的G------某个没预见到的G------是你的公理模式捕捉不到的。而在《对位藏头诗》中的龟蟹之战里，一旦确定了唱机的``体系结构''，这台唱机可能就没指望了。

怎么办呢？这是看不到尽头的。即便是扩充无穷多次，TNT也不能完全。所以说，TNT患的是本质不完全症，因为不完全性是TNT的基本组成部分，是TNT本性的一个基本成分，不管你使用笨拙的方法还是灵巧的方法，都无法根除。而且，这个问题会纠缠住数论的任何一种形式化版本，无论是TNT的扩充，还是TNT的修正，抑或能代替TNT的其它什么系统，都是一样。事实情况是：在一个给定的系统中，是否可能利用哥德尔的自指方法构造一个不可判定的符号串，依赖于三个基本条件：

⑴．该系统要足够丰富，以便全部所需要的有关数的陈述，无论真假，都能在其中表示。（做不到这一点，就意味着形式系统从一开始就弱得不能与TNT匹配，因为它连TNT能表示的那些数论概念都表示不了。借用《对位藏头诗》来比喻，就好像我们根本没有唱机，倒有个冰箱或其他种类的什么东西。）

⑵．所有的一般递归关系都能用该系统中的公式体现。（做不到这一点，就意味着该系统不能用定理来把握一般递归的真理。这种系统要试图去产生数论的全部真理，就只能说是不自量力了。用《对位藏头诗》来比喻，就好像虽然有一台唱机，但其保真度很低。）

⑶．公理以及根据该系统的规则所确定的印符模式，都能通过某个有终止的过程来辨认。（做不到这一点，就意味着无法区分该系统中的有效推导和无效推导------因而该``形式系统''根本就不是形式化的，更谈不上良定义的了。用《对位藏头诗》来比喻，这只不过是一台仍在制图板上，只设计了一部分的唱机。）

满足了这三个条件，就保证每个一致的系统都不完全，因为哥德尔的构造能得以实施。

引人入胜的事情是，任何这样的系统都在自身挖了洞：系统本身的丰富性使它自己垮了台。由于该系统强得能有自指句子，这种垮台现象就不可避免。在物理学中，像铀这类的裂变物质有``临界质量''这样的概念。一块固态物质，如果其质量小于临界质量，它就呆在那里平安无事，可一旦超过临界质量，这块东西就会经历链式反应，发生爆炸。我们的形式系统好像也有一个类似的临界点。在临界点之下，系统``无害''，但却与形式地定义算术真理相去甚远；一旦超过临界点，它就立即获得自指能力，从而也就注定它自己不完全。那么，什么时候一个系统会具有上面三个性质？界限是模糊的。只要获得了这种自指能力，该系统就有了一个为其自身所特制的漏洞。这个漏洞注意到该系统的特征，并针对该系统来利用这些特征。

卢卡斯式的非难

哥德尔的论证的这种扰乱神智的可重复性，已经成为很多人------尤其是卢卡斯------的作战武器，想以此来说明人类智能具有某种难以捉摸、不可言状的特点。因此``机械自动机''------也就是计算机------无法达到人类智能。卢卡斯的文章《心智、机器与哥德尔》是这样开头的：

``依我看，哥德尔定理证明了机械论是错的，也就是说，心智不能解释为机器。''\footnote{}

接下来，他作了一番论证，意思是说，要认为一台计算机具有和人一样的智能，它就得能够完成人所能做的每一项智能工作。现在卢卡斯宣称没有一部计算机能做人所能做的那种``哥德尔化过程''（这是他开玩笑用的一个不正规的术语）。为什么呢？好，随便考察一个具体的形式化系统，比方说TNT或者TNT+G，甚至也可以是TNT+G\textsubscript{ω}。我们可以轻松地写出有条理地生成该系统定理的计算机程序，而且，按这种方式最终能把事先选定的任何一条定理都打印出来。也就是说，这个生成定理的程序不会跳过定理``空间''的任何一块地盘。这样的程序由两个主要部分组成：（1）给定某些公理模式的``模子''（如果有这种东西的话）的情况下，可以把公理标出来的一个子程序；（2）利用已知的定理（当然包括公理在内）和推理规则生成新定理的子程序。主程序就是让两个子程序交替运行。

我们可以用拟人化的语言，说该程序``知道''某些数论事实------也就是说，它知道它打印出来的那些事实。如果有一个真的数论事实打印不出来，那它自然就不``知道''这一事实。因此，要是能说明人确实知道程序不能知道的某些事情，那计算机程序与人相比就等而次之了。这就是卢卡斯的出发点。他说，对任何一个强如TNT的系统------因而究竟是哪个系统是不重要的------我们人类总能实施哥德尔手段，所以我们总比它知道得多。这样听起来好像只是一场有关形式系统的争辩，不过稍加修改之后，似乎就变成反驳人工智能终会重现人类水平的智能的一个战无不胜的论证了。其要点是：

\begin{itemize}
\tightlist
\item
  刻板的内部编码完全控制了计算机和机器人，于是\ldots\ldots{}
\item
  计算机同构于形式系统，那么\ldots\ldots 任何一台想要和我们一样聪明的计算机就必须能对数论做我们能做的事，所以\ldots\ldots{}
\item
  除其它事情之外，它必须能完成原始递归算术，而正因为如此\ldots\ldots{}
\item
  它就很容易上哥德尔的``钓钩''，这就意味着\ldots\ldots{}
\item
  由于我们的人类智能，我们能编造出一个数论语句，它是真的，但计算机对该语句的真实性却视而不见（也就是说，永远打印不出来），其原因恰恰在于哥德尔那种``自食恶果''的论证。
\item
  这意味着有那么一件事是必定不能给计算机编上程序来做的，而我们却能做。所以我们要更聪明一些。
\end{itemize}

让我们和卢卡斯一道来享受一会人类至上论的荣耀吧：

无论我们构造出多么复杂的机器，只要它是机器，就都对应于一个形式系统。接着就能找到一个在该系统内不可证的公式而使之受到哥德尔过程的打击。尽管这个公式是真的，机器却不能生成它，而人类心智则能看出它是真的。因而这部机器仍然不是心智的一个够格的模型。我们总是力图制造心智的一种模型，它是机械的------从本质上讲是``死''的------而心智事实上是``活''的，它总能比任何形式的、僵死的系统干得好些。幸亏哥德尔定理，永远是心智说最后一句话。\footnote{}

乍一看------而且就是仔细分析一遍大概也一样------卢卡斯论证似乎无可辩驳。这引起了某种两极分化的反应。某些人就把它当成灵魂存在性的一个近乎宗教式的证明，而另一些人则嘲笑说它不值一评。我认为它不对，但却又十分迷人，所以------而且十分------值得花点时间去驳斥它。事实上，它也是推动我思考本书内容的早期动力之一。我打算在本章中用一种方法反驳它，而在第十七章中再用另一种方法来反驳。

我们必须设法更深刻地理解卢卡斯为什么会说不可能编出程序来使计算机``知道''得和我们一样多。从根本上讲，其思想就在于我们总是处于系统之外，从外面我们总能实施``哥德尔化''运算，以得到程序无法从内部看到的某种真的东西。可是，为什么不能把``哥德尔化算子''（按卢卡斯的叫法）作为第三个主要成分编到程序中去呢？卢卡斯解释说：

用以构造哥德尔公式的那个过程是个标准过程------因而就能确保对每个形式系统都能构造出一个哥德尔公式。可既然它是个标准过程，那也就能编出程序使一台机器能实施之\ldots\ldots 而这又对应着存在一个具有一条附加推理规则的系统，使我们能把原系统的哥德尔公式作为一条定理加进去，然后又把这个新的、强化了的形式系统的哥德尔公式加进去，如此下去。这相当于在原先的形式系统中加进一个无穷的公理序列，即加进逐次得到的每个系统的哥德尔公式\ldots\ldots 我们或许会期望，有一个心智在面对着这台拥有哥德尔化算子的机器时，把这一点也考虑到了，然后对新的机器、哥德尔化算子、以及一切的一切再来个哥德尔化。事实上，已经证明的确如此。即使我们把由各个哥德尔公式所组成的那个无穷公理集加进该系统，所得到的系统仍是不完全的，仍会含有一个在该系统内不可证的公式，而一个有理性的人站在该系统之外时却能看出它是真的。我们已经料到这一点，因为即使加进一个无穷的公理集，它们也得靠某种有穷的规则或规格来指明，而这种进一步的规则或规格也会被心智在考察扩大了的形式系统时考虑进去。在某种意义上，恰是由于心智说最后一句话，所以它总能从任何一个被当作它自己工作的模型的形式系统中挖出一个洞。从某种角度看，机械模型必定是有穷的、确定的，所以人的心智总能做得更好。\footnote{}

跳高一维

\pandocbounded{\includegraphics[keepaspectratio]{./media/OEBPS/Images/img_076.jpg}}

图76．龙

艾舍尔作（木刻，1952）

艾舍尔提供的一个形象画面对于支持这里的直观极有用处：那就是他的作品《龙》（图76）。这幅画最突出的特点------自然也是它的主题------是一条咬着自己尾巴的龙。这里具有全部的哥德尔内涵。不过，这幅画还有更深一层的主题，艾舍尔本人曾写了下述很有意思的注释。第一条注释说的是他的一组素描，它们全与``平面和空间的冲突''有关，第二条注释是专门讲龙的。

Ⅰ．我们所在的三维空间是我们所知道的唯一的现实世界，二维空间就像四维空间一样完全是虚构的，因为没有什么东西是平的，哪怕是磨得最光的镜子。不过，我们仍然认为一堵墙或一张纸是平面，而且------说起来也怪------就像人们自古以来所做的那样，我们仍将继续由于这种平的表面而对空间产生错觉。画上几条线，然后宣布：``这是房子。''这的确有点荒谬。这种奇特的情形就是下面五幅画（包括《龙》在内）的主题。\footnote{}

Ⅱ．不管这条龙多么想变成空间的形象，它仍然是待在平面上。印着龙的这张纸上有两处切口，它们折起来露出了两个方孔。可这条龙是个顽强的家伙，它不顾自己的二维性，坚持认为它具有三维性，于是把头伸出其中一个孔，把尾巴伸出另外一个孔。\footnote{}

这第二条注释很能说明问题。其要点是，不管你多么聪明地想在二维中模拟三维，总要遗漏某些``三维的本质''。这条龙极力想挣脱它的二维性，无视自己画在其上的那张纸是二维的，把头伸了出来。然而自始至终，我们站在画外却可以看到这完全是可悲而徒劳的。因为，龙也好，孔也好，折起来也好，统统不过是这些概念在二维空间中的模拟，没有一个是真的。然而龙却走不出它的二维空间，因而不能像我们一样知道这一点。事实上，我们还可以把艾舍尔的画里所表述的想法无限地做下去。例如，我们可以把它从书上撕下，叠起来，在上面剪个洞，再从洞里把它掏过来，为把这团东西变成二维的，我们再给它照相。对照片，我们还可以施用同样的手法。每一次，只要是最后又变到二维的------不论我们以为已经多么巧妙地在二维空间中模拟了三维空间------它就又得遭到切开，再叠起的厄运。

现在，借用艾舍尔的这个极好的隐喻，我们回过头来看看程序与人的对垒。我们谈过打算把``哥德尔算子''装进程序自身的事情。那好，即使我们已经写出了一个执行这一运算的程序，这个程序仍然抓不住哥德尔方法的实质。因为再一次地，我们在系统外面还能用它自己做不到的办法来``将死''它。可是，这样一来，我们还是在与卢卡斯辩论吗？还是在反对他吗？

智能系统的限度

是在反对。正是由于我们无法写出执行``哥德尔化''的程序这样一个事实，才使我们多少有点怀疑我们自己是不是真能在一切情形下都能运用哥德尔手法。抽象地论证``能做''哥德尔化是一回事，具体搞清在每种情形下究竟该如何做又是一回事。事实上，随着形式系统（或程序）的复杂性逐步升级，我们自己做哥德尔化的能力最终也会开始动摇。因为，如我们前面说过的，毕竟没有任何算法手段能描述如何去施行它。如果我们没有办法明确地说出在所有的情形下使用哥德尔方法会遇到些什么，那么对我们每个人来讲，最后都会碰到复杂得简直无法搞清该如何下手的情形。

自然，人的能力的这种分界线总有些说不清楚，就像一个人到底能从地上提起多重的东西这个界限一样。某一天你可能提不起250斤的东西，可另一天可能又提得起。不过不管怎么说，总不会有一天你能提起250吨重的东西。按照这样一种意义，尽管每个人哥德尔化的界限是模糊的，但对任何一个人，都存在一个系统是大大超出他的哥德尔化能力的。

在《生日大合唱哇哇哇乌阿乌阿乌阿\ldots\ldots》中直观地说明过这种想法。一开始，看上去很明显，乌龟想要纠缠阿基里斯多久就能纠缠多久。而后来阿基里斯则试图一举概括全部答案。这是和前面做出的任何反应都截然不同的一招，从而得到一个新名字``ω''。这个名字的新颖性十分重要。它是不得不超出老的命名方案------那只包括了全体自然数的名字------的第一个例子。此后就得到某些进一步的扩充，其中一些扩充的名字颇为淸楚，而另一些则相当玄妙。不过，我们最终再一次用光了名字，这时我们把这些答案模式

ω、ω\textsuperscript{ω}、\gebfont{ꪣ}

全部集结为一个复杂得难以想象的答案模式。给它起个全新的名字，叫``ε\textsubscript{0}''。之所以需要一个新名字，是因为采取了一类全新的步骤------遇到了某种不规则性。所以必须要提出一个特定的新名字。

不存在给序数命名的递归规则

现在，你可能会冒出一个想法：在从序数到序数（我们用``序数''指称那些无穷量的名字）这一进程中，那些不规则性可以由一个计算机程序来处理。也就是说，可以有一个程序按某种规则的方法来产生新的名字，而在它力不从心的时候，则乞灵于能提供新名字的``非规则性处理器''，然后再回到那个简单程序上。但这样还是不行。结果会是：这些不规则性本身也是按一种不规则的方式发生的。所以我们还得有一个二阶程序------也就是说，要有一个生产那些制造新名字的程序的程序。即使如此，也仍然不够，最终还得要一个三阶程序。如此下去，再如此下去。

这种也许看上去颇为荒唐的复杂性，完全出自阿朗佐·丘奇和斯蒂芬·克利尼关于这些``无穷序数''的结构的一条深刻的定理：

{不存在能给每个构造性序数命名的递归相关的标号系统。}

什么是``递归相关的标号系统''，什么是``构造性序数''，我们只得留待技术性更强的资料------如哈特利·罗杰斯的著作------去解释。不过，直观的想法已经给出了。随着序数越来越大，就有不规则性和不规则性中的不规则性，以及不规则性中的不规则性中的不规则性，等等。任何一个方案，无论它多复杂，都不能给所有的这种序数命名。由此可知，没有一个算法型的方法能说清如何对所有种类的形式系统使用哥德尔方法。除非我们固执得不可理喻，我们只能承认任何一个人都必将在某一点上达到他自己作哥德尔化能力的极限。超过这一点之后，具有这种复杂性的形式系统尽管由于哥德尔的理由是不完全的，但却会和人一样强有力。

对卢卡斯的其它反驳

这不过是反驳卢卡斯立场的一种方法。还有一些方法可能更为有力，我们后面将会谈到。不过，这里的这个驳论特别会使人感兴趣，因为它提出了试图创造一个能够走出自身、完全从外部观察自己、并对自己使用哥德尔手段的计算机程序的想法。当然，这就像一台唱机能播放会导致自己崩溃的唱片一样，是不可能的。

不过，不应该因为这个理由而认为TNT有缺陷。如果什么地方出了毛病，也一定不是在TNT内，而是由于我们对它所能做的事情期望过高。此外，认清下面这一点是有好处的：我们同样很容易受到被哥德尔移植到数学形式主义中的文字圈套的打击：那就是说谎者悖论。这一点由怀特利巧妙地指了出来。他提出了一个句子：``卢卡斯不能前后一致地断定本句子。''如果你想一下这句话，那就会看出：（1）它是真的，以及（2）卢卡斯不能前后一致地断定它。所以，对于世界上的真理而言，卢卡斯也是不完全的。他大脑组织中反映世界的方法使他不能既是前后一致的，同时又能断定那个真句子。不过卢卡斯也并不比我们中的任何一个人更脆弱。他和一个老练的形式系统是一个数量级的。

搞清卢卡斯论证中的错误的一个有趣方法，是把它翻译成男人和女人之间的一场较量\ldots\ldots 从前有那么一天，思想家卢克思在游历的途中，碰见了一个不认识的东西------一个女人。这样的东西他以前从未见过。一开始，他由于她长的与自己很像而非常激动，但之后就有点害怕了。他冲着周围的男人们喊道：``瞧！我能看见她的脸，而这是她做不到的事------所以女人决不会和我一样！''这样，他就证明了男人比女人优越，变得轻松多了（而他的男同胞也都放下心来）。顺便插一句，同样的论述也能证明卢克思比其他所有的男人也都优越------可他并没有对他们指明这一点。那个女人转过身来争辩说：``对，你能看见我的脸，这我做不到------但我能看见你的脸，这件事你也做不到！我们打平手。''然而，卢克思却提出一个出人意料的反驳：``很遗憾，如果你以为你能`看见'我的脸，那就错了。你们女人所做的事情与我们男人并不一样。正如我已经指出的，你们是低能的，因而不配用同一个名称。你可以称之为`妇见'。你能`妇见'我的脸，这并不说明问题，因为情况并不对称。懂了吗？''\,``我妇懂了''，那个女人``妇答''道，然后就``妇走''了\ldots\ldots{}

是的，这是一种鸵鸟式的论点。可你要是也确信在这场智力竞赛中男人们和女人们永远跑在计算机前面，那你就不得不这样看问题。

超越自我------一个现代的神话

仔细想一想我们人类是否能跳出自身------或者是，计算机程序是否能跳出自身------仍是极为引人入胜的。当然，程序有可能自己修改自己------不过这种可修改性必须首先是该程序中固有的部分，所以不能算做``跳出系统''的例子。不论一个程序怎样迂回曲折地来到自己外面，它仍然遵循自身的固有规则。对程序来说，逃脱的可能性并不比一个活人决计不服从物理定律的可能性更大。物理定律是个强制系统，什么也逃不掉。然而一个不那么宏大的抱负是可能实现的，那就是：一个人当然可以从他的大脑的一个子系统跳到一个更宽广的子系统中去。我们偶尔也能打破常规。不过这仍然属于人的大脑中各个子系统之间的相互作用，只不过会觉得它很像是完全走出了自身。与此类似，完全可以想象把部分``走出自身''的能力嵌入一个计算机程序。

无论如何，弄清楚``感知自身''和``超越自身''之间的区别是很重要的。你可以用各种各样的方法来看到自己------用镜子、用照片或电影、用录像带、靠别人的描绘、采用精神分析法等等，但你决不可能冲破皮肤站到自己外面来（这里不谈那些现代神秘运动和流行的``超心理学''流派）。TNT能够谈论自身，但不能跳出自身。一个计算机程序可以修改自己，但不能违背自身的指令------充其量也只能通过服从自身的指令来改变自己的某些部分。这倒使人想起一个幽默的悖论问题：``上帝能不能造出一块他自己举不动的石头？''

广告和框架手法

这种跳出系统的愿望到处可见。在艺术、音乐及人类其它事业的发展过程中，都有它的踪迹。就连制作广播电视广告这类不足称道的行业里，也有它的存在。欧文·高夫曼在他的著作《框架分析》中出色地领悟和描绘了这种潜在的趋向：

例如，一位著名的专业演员做完了一个广告，摄影机还在对着他，他就明显地做出解除了任务的样子，沉浸于享用他刚刚宣传的那种产品的实际快乐之中。

当然，这只是在制作电视和广播广告时利用框架手法的一个例子。这种表面上的自然而然能克服（他们这么希望）观众中蔓延着的冷漠性。于是就常常会有儿童的声音，以期造成一种似乎是未经训练印象；利用街上的噪音和其他一些效果，以期造成被访问者并不领取酬金的印象；错了的开始、不时停顿、插科打浑、以及重叠的话语，给人造成现场谈话的假象；另外，按维勒斯的方法：打断一个企业单调重复的广告而报告该企业推出新产品的消息，偶尔也插入公众所感兴趣的简短通讯。这种作法似乎可以保持现代观众的信任。

观众越是退而把一些次要的细节当作检验标准，广告商就追得越紧。结果就是一种相互污染，一种杂乱无章。这种现象同时也在被政治家的公共关系顾问以及微观社会学传播着（后者稍好一些）。\footnote{}

这样，我们又有了另一个逐步升级的``龟蟹之战''的例子------这次是事实和广告对垒。

辛普利奇奥、萨尔维亚蒂、萨哲杜：为什么要三个？

跳出系统的问题与追求完全的客观性之间，有一个引人入胜的联系。我在读尧奇的《量子是实在的吗？》中的四段对话（基于伽利略的四段《关于两种新科学的对话》）的时候，发现自己搞不清为什么要有三个角色出场：辛普利奇奥、萨尔维亚蒂和萨哲杜。只有辛普利奇奥这位受过教育的傻瓜和萨尔维亚蒂这位学识渊博的思想家两个人为什么就不行呢？萨哲杜有什么用处？是的，他被当作中立的第三派，不偏不倚地掂量两边的份量，作出``公平的''、``无偏见的''裁决。看起来很公正，可还有一个问题：萨哲杜总是赞成萨尔维亚蒂而从不赞成辛普利奇奥。人格化的客观性怎么变成了人的好恶了呢？当然，一种答案就是：萨尔维亚蒂在阐述正确观点，所以萨哲杜没有选择的余地。可是，公平性和``机会均等''又在哪儿呢？

由于增加了萨哲杜，伽利略（以及尧奇）就把牌洗得对辛普利奇奥多一些不利，而不是少一些。也许需要增加一个更高层次的萨哲杜------他对整个环境都是客观的\ldots\ldots 你可以看到这会导致什么。我们将得到一个``客观性逐步升级''的无终点序列，它有着``决不会比第一层更客观''这样一个吸引人的性质：在第一层上，萨尔维亚蒂就是对的，而辛普利奇奥是错的。所以，疑难点仍在于：为什么增加萨哲杜？答案是：就某种直观感染力而言，这种做法造成了走出该系统的错觉。

禅宗和``走出''

在禅宗里，我们也能看到这种对超出系统的概念的神往。例如那个洞山告诉弟子们``佛性上人是非佛''的公案。也许，超越自我乃是禅宗的中心主题。禅宗门人总是企图更深刻地了解他自己是怎么回事，靠的就是逐步走出他所看到的自己，并打破那些他领悟到是在束缚自己的清规戒律------当然，也包括禅宗本身的那些清规戒律。沿着这条难以捉摸的路线下去，也许到某一处就顿悟了。在任何情形（如我见到的那样），希望都在于：通过逐步加深一个人的自我意识，逐渐扩展``该系统''的范围，他最终将会感到与整个宇宙相一致。

\begin{enumerate}
\item
  \phantomsection\label{Chapter15.xhtml_B_1}
\item
  \phantomsection\label{Chapter15.xhtml_B_2}
\item
  \phantomsection\label{Chapter15.xhtml_B_3}
\item
  \phantomsection\label{Chapter15.xhtml_B_4}
\item
  \phantomsection\label{Chapter15.xhtml_B_5}
\item
  \phantomsection\label{Chapter15.xhtml_B_6}
\end{enumerate}

\phantomsection\label{Dialog16.xhtml}{}

\chapter{一位烟民富于启发性的思想}

阿基里斯应邀来到螃蟹的家里。

\dialogname{阿基里斯：}我发现比我上次来这儿时你又添了些新东西，老蟹。你那幅新弄到的画尤其引人注目。

\dialogname{螃蟹：}谢谢。有些画家我很喜欢------特别是雷内·马格里特。我所拥有的画中大部分都是他作的。他是最对我口味的艺术家。

\dialogname{阿基里斯：}我得说它们的造形都挺有趣。马格里特的这些画中有许多地方都使我想起最对我口味的画家艾舍尔的作品。

\dialogname{螃蟹：}这一点我也能看得出来。马格里特和艾舍尔在表现悖理的和虚幻的世界时，都运用十分写实的手法，两个人在运用视觉形象以唤起观众的情感方面都具有很准确的感觉能力，而且------他们作品的爱好者们常常忽视这一点，就是他们都具有对优美线条的感受力。

\dialogname{阿基里斯：}不过，他们彼此之间还是有很大不同的。我不知道怎么描述这种区别。

\dialogname{螃蟹：}细致入微地比较这两个人一定非常引人入胜。

\dialogname{阿基里斯：}我认为，马格里特对现实的把握是惊人的。比如那儿挂着的那幅画吧，画面上有一棵树，树后面还有一支巨大的烟斗，画得太迷人了。

\pandocbounded{\includegraphics[keepaspectratio]{./media/OEBPS/Images/img_077.jpg}}

图77．影子

马格里特作（1966）

\dialogname{螃蟹：}你是想说一只大小正常的烟斗和它前面的一棵微型树吗？

\dialogname{阿基里斯：}哦，是这样，我看错了？嗯，不管怎么说，我第一眼看见它时，就觉得闻到了烟斗发出的烟味儿！你想象得出我有多荒唐吗？

\dialogname{螃蟹：}我能理解。我的客人们经常被它骗住了。

（他一边说着，一边过去伸手把画中的烟斗从树后拿出，然后翻过来，在桌上磕了磕，房间里马上充满了烟味儿。他开始往里面装新的烟丝。）

这是只很不错的旧烟斗，阿基。信不信由你，我总爱往烟丝里掺点稻草，这样抽起来味道更好。

\dialogname{阿基里斯：}掺稻草！真的吗！

\dialogname{螃蟹：}\kaitismall{（拿出一盒火柴，点上烟斗）}你想尝尝吗，阿基？

\dialogname{阿基里斯：}不，谢谢。我只是偶尔吸吸雪茄。

\dialogname{螃蟹：}没问题！雪茄我这儿也有！\kaitismall{（把手伸向马格里特的另一幅画，画面上一辆自行车放在一支点燃的雪茄上。）}

\pandocbounded{\includegraphics[keepaspectratio]{./media/OEBPS/Images/img_078.jpg}}

图78．优雅状态

马格里特作（1959）

\dialogname{阿基里斯：}噢------不，谢谢，现在还不想抽。

\dialogname{螃蟹：}你随便。我这人可是个不可救药的烟鬼。这倒叫我想起来------你无疑知道老巴赫有用烟斗吸烟的嗜好？

\dialogname{阿基里斯：}我想不大起来了。

\dialogname{螃蟹：}老巴赫喜欢作诗、玄想、吸烟，还喜欢作曲演奏（当然喜好程度不一定是这么个顺序）。他把这四者弄进一首不太高明的诗里，然后还为它谱了曲。这写在他为妻子安娜·马达莱娜所保存的那本著名的音乐笔记本里，它的题目是

一位烟民富于启发性的思想\footnote{}

我拿起烟斗装满叶烟，

吞云吐雾以消磨时间，

人在那儿坐着，却溜走了思想------

它飘向一幅图画，晦黯而忧伤：

那画告诉我这是多么相像：

我和这烟斗简直就是一样。

这喷香溢郁的烟斗同我相仿，

都不过是尘芥，不过是土壤；

我最终也要归为尘灰。

烟斗落地，没等听到声响，

就已经拦腰摔断，不幸遭殃；

相同的命运我也将承当。

洁白的烟斗没有污斑，

从未玷染，从未弄脏。

可总有一天会命归无常，

在草地下掩埋这副皮囊；

我的肉体将变黑，一副腌{ꪡ}相，

就像这烟斗，若是它使用得经常。

当烟斗被点燃，闪耀着火光，

瞧，它们立刻就冒出青烟袅袅飘荡；

轻烟散进空气，无处寻访，

烟斗里便只剩有灰烬留藏。

虚名也将消泯，正同那青烟一样，

而躯体最终只不过化作土壤。

吸烟时这种事也发生得经常：

架上的塞烟器不翼而飞，给你添忙，

这下你只好上手，虽然不很便当，

可手指伸进烟锅，难免要被烫伤。

既然烟斗里都有痛苦伏藏，

地狱里的痛苦将会何其难当！

由对烟斗的沉思，导向了别的玄想，

沉溺在格致冥思中，

也有益于身心的健康；

这样吞云吐雾，心中着实舒畅，

因此无论在国内、国外、在陆地或海洋，

我都要一边抽我的烟斗，一边坚定对上帝的信仰。

这种哲学倒是满有意思的，不是吗？

\dialogname{阿基里斯：}的确是。老巴赫是舞文弄墨的一把好手。

\dialogname{螃蟹：}你正好说出了我要说的话。其实我这辈子也曾写过一些打油诗，可我担心我的诗鄙陋不中雅听。我没有这种遣词炼句的本领。

\dialogname{阿基里斯：}哦，哪里，老蟹。你------我怎么说呢------论俏皮、戏谑样样都有。要是你能给我唱一支你的歌，老蟹，我将感到不胜荣幸。

\dialogname{螃蟹：}太抬举我了。我给你放一张我的唱片怎么样？我记不得它是什么时候灌制的了。歌的名字是``云俦梅友''。

\dialogname{阿基里斯：}太有诗味了！是首山水诗吧？

（螃蟹从他的架子上拿下一张唱片，走到一架庞大而复杂的机器面前，把它打开之后，将唱片放入一个看上去怪吓人的机械嘴里。突然，一道绿色的亮光扫过唱片的表面，过了一会儿，那唱片就悄没声地闪进那架古怪机器的肚子里的什么地方去了。又过了一会，才传出了螃蟹的歌声。）

舞文弄墨的一把好手，

论俏皮、戏谑样样都有。

诗到最后他有点疏漏，

似乎留有一点缺口，

我的意思是不知所云。

\dialogname{阿基里斯：}好极啦！不过，有一处我不太满意。在我看来，诗到最后你有点疏漏------

\dialogname{螃蟹：}似乎留有一点缺口？

\dialogname{阿基里斯：}不------我的意思是韵脚没有。

\dialogname{螃蟹：}也许你对。

\dialogname{阿基里斯：}除了这一点之外，这是首很好的诗。可我得说更叫我着迷的是这台异常复杂的怪家伙。这是台超大型号的唱机吗？

\dialogname{螃蟹：}哦，不是，远不止这么简单。这是我那架吞龟唱机。

\dialogname{阿基里斯：}天呐！

\dialogname{螃蟹：}嗯，我不是说它能吞吃乌龟，但是它能嚼啐龟兄制造的唱片。

\dialogname{阿基里斯：}哦，这还不那么吓人。这就是不久前你跟龟兄那场奇特的音乐战的一部分吗？

\dialogname{螃蟹：}也可以这么说。让我给你详细讲讲吧。龟兄鬼极了，他几乎可以摧毁我能弄到的任何一架唱机。

\dialogname{阿基里斯：}不过我最后一次听到你们俩的较量时，我记得你好像最后弄到一种打不败的唱机------一种带有内隐电视摄像机和微484型计算机这一类东西的机器，它可以自行拆卸，然后重新组装成不能被摧毁的结构。

\dialogname{螃蟹：}呜乎哀哉！我的设计失败了。龟兄利用了被我忽略的一个细节：那个控制着拆卸和组装的机构在整个过程中是要保持不变的。这就是说，很显然，它没法儿把它自己拆开或组装，所以它始终原样不动。

\dialogname{阿基里斯：}对，可这会怎么样呢？

\dialogname{螃蟹：}哦，这样一来就完蛋了！龟兄就完全把攻击点对准了那个机构。

\dialogname{阿基里斯：}怎么？

\dialogname{螃蟹：}就是说他只要制造这样一张唱片就行了：这张唱片能引起那个始终保持不动的结构------拆卸-组装机构------产生致命的颤动，这样一来\ldots\ldots{}

\dialogname{阿基里斯：}噢，我明白了------够狡猾的。

\dialogname{螃蟹：}没错儿，我也这么想。他的方法还真奏效了。可是，你知道，这回他不是一次成功的。这台唱机在遭受了他第一次进攻之后，并没有被损坏。当时我以为在这场斗智中我终于嬴了。我那时真有点乐不可支。可当他又一次带着一副冷峻的表情到我这里来的时候，我意识到这回是来者不善。我把他的新唱片放到我那架唱机转盘上。我们俩都急切地瞧着电脑控制的机构仔细检查着唱片上的纹线，然后拿下唱片，拆散唱机，又用一种完全不同的方式重新把它组装好，再次放上唱片------再缓缓地把唱针放到最靠外的纹线上。

\dialogname{阿基里斯：}乖乖！

\dialogname{螃蟹：}唱机刚一出声，就听轰隆一声，响彻全屋！整个机器都震散架啦，那个组装-拆卸器损坏得尤其厉害。这时我终于痛苦地意识到乌龟总是能把攻击目标对准系统的那个阿基里斯脚跟------噢，我是说它的致命弱点。

\dialogname{阿基里斯：}真要我的命！你一定非常沮丧吧。

\dialogname{螃蟹：}是啊，我一时绝望极了。不过，幸运的是这个故事到这里并没有结束。这件事的一个结果是我从中得出了很有益的教训，让我来说给你听。在乌龟的建议下，我浏览了一本稀奇古怪的书，里面穿插有好多奇特的对话，那些对话涉及了许多学科，有分子生物学、赋格曲、禅宗等等五花八门的玩艺儿。

\dialogname{阿基里斯：}说不定是个疯子写的呢。书名是什么？

\dialogname{螃蟹：}要是我记得不错，书名是《金、银、铜------聚宝藏之精华》。

\dialogname{阿基里斯：}哦，龟兄跟我说起过。是他的一个朋友写的，听说那人对各种材料制成的怪圈很着迷。

\dialogname{螃蟹：}我想知道这是他的哪位朋友------不管怎么说。在其中的一篇对话里，我还看到一些富于启发性的思想，涉及烟草花叶病毒、核糖体，以及其它一些我闻所未闻的东西。

\dialogname{阿基里斯：}烟草花叶病毒是什么玩艺儿？核糖体呢？

\dialogname{螃蟹：}我说不清楚，因为一讲到生物学，我就全傻了。我所知道的只是对话里讲到的那一点儿，那篇对话里说烟草花叶病毒是一种卷烟形状的、能使烟草生病的小东西。

\pandocbounded{\includegraphics[keepaspectratio]{./media/OEBPS/Images/img_079.jpg}}

图79．烟草花叶病毒

选自莱宁格{[}A, Leninger{]}，《生物化学》{[}Biochemistry{]}，纽约：Worth Publishers，1976年版

\dialogname{阿基里斯：}生癌吗？

\dialogname{螃蟹：}不，不完全一样，而是------

\dialogname{阿基里斯：}难道烟草自己吸烟，得了癌症？正好！

\dialogname{螃蟹：}我看你太着急下结论了。阿基。烟草不吸这些``卷烟''。是那些脏乎乎的小``卷烟''来进攻它们，不请自来。

\dialogname{阿基里斯：}我明白了。嗯，我现在知道烟草花叶病毒是怎么回事了。那么核糖体呢？

\dialogname{螃蟹：}核糖体显然是一种细胞里的东西，可以把某一形式的信息转换成另一种形式的信息。

\dialogname{阿基里斯：}有点像微型录音机或唱机什么的，是吗？

\dialogname{螃蟹：}作为比喻来说是像。让我感兴趣的是有个活宝在一句话里提到了这样一个事实：同烟草花叶病毒和其它一些稀奇古怪的生物结构一样，核糖体也具有``自发地自组装''这种不可思议的能力。这就是那家伙的原话。

\dialogname{阿基里斯：}这是他那些鬼话里的一句，是吧。

\dialogname{螃蟹：}这也正是那篇对话里另一个角色的看法。不过这是对那句话的蠢解释。

（螃蟹又往烟斗里添了一点稻草末，然后深吸了几口烟斗，一口口地把烟喷到空气中。）

\dialogname{阿基里斯：}嗯，``自发地自组装''是什么意思？

\dialogname{螃蟹：}意思是当某个位于细胞内的生物单位分解时，它们可以自发地把它们自己再组装起来------用不着被别的单位来控制。那些部分只要凑在一起------疾！------就粘在一块儿了。

\dialogname{阿基里斯：}像魔术似的。要是普通型号的唱机也有这种本事该多好！我是说，既然一部像核糖体这样的微型``唱机''能有这种功能，为什么大一点儿的就不行呢？这样你就能制造出一架不能被摧毁的唱机了，对吗？无论什么时候它被拆散，都能再自行组合起来。

\dialogname{螃蟹：}我正是这么想的。我急不可耐地给我的制造商写了一封信，向他解释了自组装的设想，问他能不能给我制造一架可以自行拆散，并能够自发地组装成另一种形式的唱机。

\dialogname{阿基里斯：}造这么一个玩艺可太劳神了！

\dialogname{螃蟹：}的确，主要是太劳财神了，因为几个月以后，他写信时不仅告诉我他终于成功了，还送来了一张金额巨大的帐单。在这样一个好天里，嘿！我那了不起的自组装唱机终于寄来啦，我信心十足地给龟兄打电话，邀他过来检验检验我这完备的唱机。

\dialogname{阿基里斯：}我们眼前这台大家伙一定就是你说的那台吧。

\dialogname{螃蟹：}恐怕不是，阿基。

\dialogname{阿基里斯：}是不是它又给------

\dialogname{螃蟹：}我的老朋友，你所担心的正是事实。我不知道事情是如何发生的。描述当时的情景真太叫人伤心了。眼巴巴地瞧着电线啦、碎片啦乱七八糟地摊了一地，还这儿一股那儿一股地冒着烟------啊，我------

\dialogname{阿基里斯：}啊，啊，老蟹，别太想不开。

\dialogname{螃蟹：}我倒没什么，只是这一次次的太快了。哼，龟兄开始还幸灾乐祸，最后他看我是真伤心了，于是就同情起我来了。他为了安慰我，就解释说这是理应如此的事，谁也没办法------说这跟什么人的什么``定理''有关，不过那定理我一个字也记不得了，只记得好像叫什么``疙瘩定理''。

\dialogname{阿基里斯：}我怀疑这就是他以前跟我说起过的``哥德尔定理''\ldots\ldots 是有那么点不祥的色彩。

\dialogname{螃蟹：}也许吧，我记不起了。

\dialogname{阿基里斯：}我向你保证，老蟹，我听你这个故事时是非常同情你的处境的。真太叫人伤心了。不过，你提到过什么稻草的事，请跟我说说，是怎么回事？

\dialogname{螃蟹：}噢，对------是有根救命稻草。嗯，我最后终于放弃了寻求``完备的''唱机的念头，而是决定进一步完善对乌龟的唱片的抵制。我抛开了那种能播放一切唱片的唱机的奢望，而是想要一种能避免破坏、能够保存下来的唱机：一种能避免被摧毁的唱机------即使这意味着它只能播放少数一些特殊的唱片。

\dialogname{阿基里斯：}于是你就决定以牺牲掉能重现所有可能的声响这一功能为代价，制造一种复杂的反乌龟的机器，对吗？

\dialogname{螃蟹：}嗯\ldots\ldots 说不上我``决定''。更严格地说，是我被迫如此。

\dialogname{阿基里斯：}对，我明白你的意思。

\dialogname{螃蟹：}我的新设想是不让任何``异己的''唱片在我的唱机上播放。我知道我自己的唱片是不会对我的机器有损害的，所以如果我不让别的唱片混进来，我就能保住我的唱机，用它来欣赏我灌制的唱片。

\dialogname{阿基里斯：}这一策略对你的目标很合适。我们眼前这个庞然大物就是你刚才说的那种想法的产物吧？

\dialogname{螃蟹：}是的。龟兄自然也认识到他也必须改变策略。他现在的主要目的就是要搞出一种能混过我的检查的唱片来------这是种新挑战啊！

\dialogname{阿基里斯：}作为你，你打算怎么排除他或别人的``异己的''唱片呢？

\dialogname{螃蟹：}你保证不把我的计划透露给龟兄吗？

\dialogname{阿基里斯：}以龟格担保。

\dialogname{螃蟹：}什么？

\dialogname{阿基里斯：}哦------别担心，这是我从龟兄那儿学来的话------我发誓坚守秘密。

\dialogname{螃蟹：}那好吧。我的基本方案是使用标识技术。我的每一张唱片上都有一个秘密标记。现在你面前的这台唱机跟以前那几台一样，都装有一部检验唱片的电视摄像机，这种检验器上配有一台计算机，负责处理那些由检验得到的数据，并控制相应的操作。我打算嚼碎那些没有正确标记的唱片。

\dialogname{阿基里斯：}哈，报复他一下！不过我觉得你的方案很容易被挫败。龟兄只要搞到一张你的唱片，复制下那标记，你就没咒念了。

\dialogname{螃蟹：}没那么简单，阿基。谁告诉你他能从唱片的非播放状态中知道那个标记的？事情比你想的要复杂哩。

\dialogname{阿基里斯：}你是不是说标记跟播放出的音乐是掺在一起的？

\dialogname{螃蟹：}正是。但是有办法把它们分解开。这需要从视觉上将那种数据从唱片中抽取出来，然后------

\dialogname{阿基里斯：}那种明亮的绿光就是干这个用的，是吗？

\dialogname{螃蟹：}对。那是电视摄像机在检查音纹。音纹格式被送到微型电脑里，由它来分析我放进去的唱片的音乐风格------这一切都是无声地进行的。这时不播放唱片。

\dialogname{阿基里斯：}是不是有一个甄别过程，将风格不合适的唱片排除掉？

\dialogname{螃蟹：}你说对了，阿基。能通过这第二道检查的唱片就只有同我那些唱片风格一样的唱片了------龟兄要摹仿这些简直就是不可能的。所以我坚信我一定能打赢这场新的音乐战。不过，我得说龟兄也同样相信他会混过我的检查。

\dialogname{阿基里斯：}最后把这巨大的机器震成碎片？

\dialogname{螃蟹：}哦------不，这种结局他已经用事实验证过了。他现在只是想证明不管我采取什么措施，他都能------用一张对我的唱机无害的唱片------溜过我的检查。我听他嘴里不停地嘀嘀咕咕，总是提到一首歌曲，那歌曲的名字挺古怪，叫什么``我可以在唱机X上被播放''。可他吓唬不了我！唯一叫我有些担心的是他似乎同以前一样，又有些什么晦涩的观点，那些观点\ldots\ldots{}\kaitismall{（他陷入沉默，随后，满脸沉思的吸了几口烟斗。）}

\dialogname{阿基里斯：}哼\ldots\ldots 我看龟兄那方面有无法克服的困难。他终于碰到对手了！

\dialogname{螃蟹：}原来你也这么想\ldots\ldots 我想你不知道汉肯定理的来龙去脉吧！

\dialogname{阿基里斯：}知道谁的定理的来龙去脉？我从没听到过这类事。我相信那一定很吸引人，可我更想多听一些有关``唱片设法混入唱机''的事。这小故事引人入胜呐。其实我都可以说出结局了。显然，龟兄没有找出混进去的突破口，于是只好乖乖认输，就这么回事，不是吗？

\dialogname{螃蟹：}至少我是希望这样。你想看看我这台防御式唱机的内部情况吗？

\dialogname{阿基里斯：}当然想。我一直想瞧瞧那架电视摄像机。

\dialogname{螃蟹：}说干就干，老朋友。\kaitismall{（把手伸进那台巨大的唱机张着的``大嘴''里，揿下两个按钮，拿出一台包装整齐的仪器。）}瞧，整个仪器是由一个个独立的模块组成的，可以分别卸下来独立使用。比如这部电视摄像机就能良好地独立工作。看着那边的屏幕，画有燃烧的大号的那幅画下面的那个屏幕。\kaitismall{（他将摄像机对着阿基里斯，阿基里斯的脸立刻出现在大屏幕上。）}

\dialogname{阿基里斯：}天哪！我可以试试吗？

\dialogname{螃蟹：}当然可以。

\dialogname{阿基里斯：}\kaitismall{（将摄像机对着螃蟹）}你上屏幕了，老蟹。

\dialogname{螃蟹：}是啊。

\dialogname{阿基里斯：}我要把摄像机对着画有燃烧的大号的那幅画。它现在也在屏幕上啦！

\pandocbounded{\includegraphics[keepaspectratio]{./media/OEBPS/Images/img_080.jpg}}

图80．漂亮的俘虏

马格里特作（1947）

\dialogname{螃蟹：}摄像机可以移近移远，阿基。你试试。

\dialogname{阿基里斯：}太妙了！让我先对准那些紧挨着画框的火焰\ldots\ldots 真是不可思议，它能在屏幕上瞬间``复制''出屋里的任何东西------任何我想摄进的东西。我只须把摄像机对准它，像变魔术一样，它立刻就会跳上屏幕。

\dialogname{螃蟹：}屋里的任何东西吗，阿基？

\dialogname{阿基里斯：}只要是看得见的任何东西，很显然。

\dialogname{螃蟹：}要是你把摄像机对准电视屏幕上的火焰，那会怎么样？

（阿基里斯调转摄像机，用它对准正播放有火焰的电视屏幕。）

\dialogname{阿基里斯：}太有趣了！这样一来火焰就从屏幕上消失了！它们哪儿去了？

\dialogname{螃蟹：}你无法在移动了摄像机时还让屏幕上保持原来的形象。

\dialogname{阿基里斯：}我明白了\ldots\ldots 可我一点也看不懂现在屏幕上是什么！它看上去像一条长长的奇怪的走廊。可我没有把摄像机对着什么走廊啊？我只是把它对着一个普通的电视屏幕。

\dialogname{螃蟹：}看仔细些，阿基。你真的看到一条走廊吗？

\dialogname{阿基里斯：}啊哈，我现在明白了。那是一组叠套在一起的电视屏幕，一个比一个小\ldots\ldots 没错儿！火焰的形象早没了，因为我的摄像机不再指向那幅画了。当我把摄像机对着屏幕时，屏幕上就出现了屏幕，包含着屏幕的屏幕上有什么，被包含的屏幕上也就有什么------而被包含的屏幕上只有屏幕，因此被包含的屏幕上有什么，被被包含的屏幕所包含的屏幕上也就有什么------而被被包含的屏幕所包含的屏幕上只有屏幕，因此------

\pandocbounded{\includegraphics[keepaspectratio]{./media/OEBPS/Images/img_081a.jpg}}\pandocbounded{\includegraphics[keepaspectratio]{./media/OEBPS/Images/img_081b.jpg}}\pandocbounded{\includegraphics[keepaspectratio]{./media/OEBPS/Images/img_081c.jpg}}

图81．十二个自噬的电视屏幕

如果13不是个素数的话，我会再加上一个的。

\dialogname{螃蟹：}我想我也能顺着说下去了，阿基。你干嘛不试试旋转旋转摄像机？

\dialogname{阿基里斯：}哦！我得到了一根螺旋的走廊！每个屏幕都与它的外框屏幕拧成一定角度，屏幕越小，与最外面的屏幕拧的角度就越大。这种让电视屏幕``自噬''的想法太奇特了。

\dialogname{螃蟹：}你说的``自噬''一词是什么意思，阿基？

\dialogname{阿基里斯：}我用``自噬''一词是指我把摄像机对准屏幕------或对准屏幕的一部分这一操作，这就是``自噬''，也就是自己吞掉自己。

\dialogname{螃蟹：}要是我要求你再说得详细点，你不会不同意吧。我有点叫这个概念给迷住了。

\dialogname{阿基里斯：}我也一样。

\dialogname{螃蟹：}那好吧。要是你把摄像机对准屏幕的一角，那还算是你说的``自噬''吗？

\dialogname{阿基里斯：}让我试试看。哎------屏幕的走廊好像跑到外边去了。这下就没有无穷的叠套了。它挺好看的，可我觉得它不具备``自噬''的因素。它是一种``不成功的自噬''。

\dialogname{螃蟹：}要是你把电视摄像机转回屏幕的中心部位，也许你又能把它定住了\ldots\ldots{}

\dialogname{阿基里斯}对！走廊变得越来越长啦\ldots\ldots 瞧啊，完全变回来啦。顺着它一直看下去，能看到它在远处消失了。当摄像机整个儿摄进屏幕时，走廊就又变成无限的了。嗯------这叫我想起不久前龟兄对我说，一个句子要是谈论整个自己，就会出现自指\ldots\ldots{}

\dialogname{螃蟹：}什么？请再说一遍。

\dialogname{阿基里斯：}没什么------我是在自言自语。

（当阿基里斯摆弄起摄像机的镜头和其它控制装置时，屏幕上就出现了前所未有的，花样繁多的``自噬''图像：形似星系的旋转螺线，变化万端的花状图形和其它一些各式各样的图形\ldots\ldots）

\dialogname{螃蟹：}你好像玩得挺起劲儿。

\dialogname{阿基里斯}是的！这么一个简单的办法能弄出多少图像啊！\kaitismall{（他又转回头盯着屏幕，脸上掠过一丝惊诧的神情。）}哎呀，老蟹！屏幕上出现了一种跳动的花瓣状图形！这种跳动是怎么产生的？电视机是静止的，摄像机也没动呀。

\dialogname{螃蟹：}偶尔你也可以弄出些在时间中有变化的图形。这是由于在摄像机``看到''什么东西和把它显示在屏幕上这两者之间时间上稍有先后之差，这是由线路造成的，这个时间差大约是百分之一秒。所以要是你弄出一个由五十个环像组成的套叠，其时间差大约就要有半秒。要是把某种活动的东西弄进屏幕------比方说，把你的手指放到摄像机前，那么，就得过一会儿才能在套得更深一些的屏幕上``把它找出来''。这种时间延宕反映到整个系统，就形成了某种视觉``回声''。如果创造出一种条件，使得这种回声不消逝，就会形成跳动的图形。

\dialogname{阿基里斯：}太了不起啦！你说，要是我们试试弄出一个全自噬会怎么样？

\dialogname{螃蟹：}你是什么意思？说得明白些。

\dialogname{阿基里斯：}嗯，我是觉得这种屏幕里面有屏幕的东西非常有趣，可我更想得到一幅里面有这部电视摄像机以及这个屏幕的图像，让它显现到这个屏幕上。那样我才能真正使一个系统自噬其自身。因为这个屏幕仅仅是这整个系统的一部分。

\dialogname{螃蟹：}我明白你的意思了。也许你能利用镜子获得你所要的效果。

（螃蟹递给他一面镜子，阿基里斯操纵镜子和摄像机，使得摄像机和屏幕都显现在屏幕上。）

\dialogname{阿基里斯：}瞧！我创造出了完全的自噬！

\dialogname{螃蟹：}我觉得只有镜子的前面还不够，加上它的背面怎么样？要是没镜子背面，那它就不能反射------你就没法儿使图像里也有摄像机了。

\dialogname{阿基里斯：}你说的对。不过，要既显出镜子的正面，也显出背面，我就不得不再需要一面镜子。

\dialogname{螃蟹：}可你还得显出第二面镜子的背面。除了显示摄像机的前面以外，把它的背面也弄进去怎么样？还有电源线，还有电视内部，还有------

\dialogname{阿基里斯：}哎哟哟！我的脑袋都晕了！看来我这个``全自噬计划''要出麻烦了。我现在觉得昏昏然的。

\dialogname{螃蟹：}我能体会到你这时的感觉，你干嘛不坐下来，抛开这些什么鬼自噬？放松放松！看看我的画你就会平静下来的。

（阿基里斯唉声叹气地躺下了。）

哦------也许我的烟味叫你不舒服？得，我拿开它。\kaitismall{（把烟斗从嘴上拿掉，小心地放在马格里特的另一幅画里面的一句话的上方。）}喏，好些了吧？

\dialogname{阿基里斯：}还是恶心。\kaitismall{（指着马格里特的画。）}这幅画挺有意思，我喜欢它的外框，很有质感。

\pandocbounded{\includegraphics[keepaspectratio]{./media/OEBPS/Images/img_082.jpg}}

图82．咏叹调和歌曲

马格里特作（1964）

\dialogname{螃蟹：}谢谢。那是我特地请人做的------用烟叶和稻草。

\dialogname{阿基里斯：}用烟叶和稻草？乖乖，真叫人难以相信！烟斗下面写的是什么字？不像英文，是吧？

\dialogname{螃蟹：}是法文。写的是``Ceci n\textquotesingle est pas une pipe.''，意思是``这不是一只烟斗''。千真万确。

\dialogname{阿基里斯：}可它是一只烟斗啊！你刚才还吸来着！

\dialogname{螃蟹：}哦，我想你误解了这句话的意思。``ceci''这个词指那幅画，而不是指烟斗。烟斗当然是烟斗，可一幅画却不是一只烟斗。

\dialogname{阿基里斯：}我不知道画中的``ceci''一词是指整幅画，还是指画中的烟斗。哦，我的天！这又是一个自噬！我可真不行了，老蟹，我想我快要病了\ldots\ldots{}

\begin{enumerate}
\item
  \phantomsection\label{Dialog16.xhtml_B_1}
\end{enumerate}

\phantomsection\label{Chapter16.xhtml}{}

第十六章

\chapter{自指和自复制}

在这一章里，我们要看看导致各种自指现象的机制，并把它们与一些使某几种系统能够复制自己的技巧加以比较。我们将展现这些技巧之间的一些值得注意的、漂亮的相似性。

隐性和显性的自指句子

作为开篇，我们先看几个句子，这些句子第一眼看去像是最简单的自指。下面是这一类句子中的几个例子：

⑴．本句子有七个字。

⑵．本句子无意义，因为它是自指的。

⑶．本句子无动词。

⑷．本句子是假的。（说谎者悖论）

⑸．正在写的这句话就是你正在读的那句话。

除了最后一句（这是个例外），其余的所有句子都涉及一个看上去很简单的、包含在词组``本句子''之内的技巧。其实，这种技巧一点也不简单。全部这些句子都``漂浮''在上下文之中。可以把它们与只能看见尖顶的冰山做个比较。这一串串字符就是冰山的尖顶，为了理解它们所必须做的处理就是看不见的那部分。按这种意义，它们的涵义是隐性的，而不是显性的。当然，没有一句话的意义完全是显性的，不过，这种自指性越是呈显性，就越能暴露支撑它的那种技巧。在我们这里，为了辨认出上述各句子的自指性，就不仅要熟悉像汉语这样的、可以用来讨论语言学问题的语言，而且还得能够领会``本句子''这个词组的所指。这好像很简单，但是却要依赖于我们那种非常复杂而又完全内化了的驾驭语言的能力。这里，尤其重要的一件事是领会一个带有指示代词的名词性词组的所指。这种能力建立得很慢，因而决不能等闲视之。当我们向某个不懂得悖论和语言学圈套的人（比方说一个孩子）介绍像第四句话那样的句子时，困难也许就会特别突出。他们可能会问：``什么句子是假的？''而要想使人了解这个句子是谈论它自己的，可能就得花上一些功夫。整个思想最初是有点令人吃惊的。

\pandocbounded{\includegraphics[keepaspectratio]{./media/OEBPS/Images/img_083.png}}

图83．说谎者处决自己

刘皓明绘

这里有两幅画（图83、图84）或许会有所帮助。图83是一幅能在两个层次上解释的画。在一个层次上，它是一个指着自己的句子；在另一个层次上，它是说谎者处决自己的写照。

而为了辨认一个句子的自指性所需要的处理的相对比例，就像是图84中的冰山的可见和不可见部分的相对比例。

\pandocbounded{\includegraphics[keepaspectratio]{./media/OEBPS/Images/img_084.png}}

图84．说谎者悖论的冰山

试着造一个自指句子，而又不用``本句子''等手法，是非常有意思的。可以试试在一个句子内部引用自己。下面就是一次尝试：

``这句话有七个字''这句话有七个字。

不过，这种尝试必然失败，因为一个句子引用自己后得到的新句子总是比自己长。只有当你情愿接受一个无穷长的句子时，这种做法才实际可行，例如：

\begin{longtable}[]{@{}lllll@{}}
\toprule\noalign{}
\endhead
\bottomrule\noalign{}
\endlastfoot
　``句子 & & & & \\
　　``句子 & & & & \\
　　　``句子 & & & & \\
　　　　``句子 & & & & \\
　 & • & & & \\
　 & 　 & • & & \\
　 & 　 & 　 & • & \\
　 & 　 & 　 & 　 & 等等，等等 \\
　 & 　 & 　 & • & \\
　 & 　 & • & & \\
　 & • & & & \\
　　　是无穷长的'' & & & & \\
　　是无穷长的'' & & & & \\
　是无穷长的'' & & & & \\
是无穷长的'' & & & & \\
\end{longtable}

但是对有穷的句子不能这样做。基于同样的理由，哥德尔的符号串G不能含有其哥德尔数的数字形式：它装不下。没有一个TNT符号串能含有其自身哥德尔数的TNT数字，因为这一数字总比该符号串本身含有更多的符号。不过，你可以让G含有其自身哥德尔数的一个描述，再利用``代入''和``算术\gebfont{㧟}摁化''来绕过去。

蒯恩的方法是使用描述------而不是自引用或词组``本句子''------来在一个句子中实现自指的方法之一，对话《G弦上的咏叹调》里对这种方法有形象的说明。理解蒯恩句子用不着像理解前面举的四个例子那样做精巧的智力加工。虽然乍一看它好像更为诡秘，其实倒是一种更为直接的方法。蒯恩的构造在下述意义上极像哥德尔的构造：它通过描述另外一个（已经证明）同构于蒯恩句子的字符串来创造自指。对这串新字符的描述，是由这蒯恩句子的两个部分实现的。一部分是一组指令，它告诉我们如何建立一个词组，而另一部分则含有所要的素材，也就是说，这另一部分是一块模板。这也像冰山，但更像一块漂在水面上的肥皂（见图85）。

\pandocbounded{\includegraphics[keepaspectratio]{./media/OEBPS/Images/img_085.png}}

图85．蒯恩句子的肥皂

这个句子的自指性是靠一种比说谎者悖论更为直接的方法达到的，用不着多少隐蔽的处理。此外，有意思的是，请注意``这个句子''出现在上面的句子里，但它在那里并不引起自指。你大约可以理解，它的所指是蒯恩句子，而不是它出现于其中的那个句子。这正好说明了诸如``这个句子''这类指示词组该如何根据上下文来解释。这还有助于说明对这种词组的处理必定是相当复杂的。

一个自复制的程序

``\gebfont{㧟}摁化''概念以及它在制造自指过程中的用处，对话中已经解释了，所以这里用不着再细讲。我们还是来说明一下计算机程序如何能利用完全相同的技巧复制它自己。下面的自复制程序是用一种类似BlooP的语言写出的，用于在一句话之后引用它自己（这与\gebfont{㧟}摁化次序相反，所以我把``\gebfont{㧟}摁''倒过来写成``摁\gebfont{㧟}''）：

DEFINE PROCEDURE `ENIUQ'{[}TEMPLATE{]}: PRINT {[}TEMPLATE, LEFT-BRACKET, QUOTE-MARK, TEMPLATE, QUOTE-MARK, RIGHT-BRACKET, PERIOD{]}.

定义过程``摁\gebfont{㧟}''{[}模板{]}：打印{[}模板，左括号，单引号，模板，单引号，右括号，句号{]}。

ENIUQ

{[}`DEFINE PROCEDURE `ENIUQ' {[}TEMPLATE{]}: PRINT {[}TEMPLATE, LEFT-BRACKET, QUOTE-MARK, TEMPLATE, QUOTE-MARK, RIGHT-BRACKET, PERIOD{]}. ENIUQ'{]}.

摁\gebfont{㧟}

{[}`定义过程``摁\gebfont{㧟}''{[}模板{]}：打印{[}模板，左括号，单引号，模板，单引号，右括号，句号{]}。摁\gebfont{㧟}'{]}。

``ENIUQ''是程序的前三行定义的一个过程，它的输入称作``TEMPLATE''{[}模板{]}。这个过程的意思是：调用这个过程时，TEMPLATE的值是某一串印刷字符，ENIUQ的结果是执行一项打印操作，其中要把TEMPLATE打印两次：第一次只打印它，第二次在外面加上引号和括号，最后还要缀上一个句号。这样，如果TEMPLATE的值是符号串DOUBLE-BUBBLE（意为``两个泡泡''），执行ENIUQ，就得到

DOUBLE-BUBBLE{[}\textquotesingle DOUBLE-BUBBLE\textquotesingle{]}。

上面的程序中，后一半是把ENIUQ这个过程调来处理一个具体的TEMPLATE值------即单引号内的那个长长的字串：DEFINE\ldots\ldots ENIUQ。这个值是精心挑选出来的，它由ENIUQ的定义后面再加上一个词``ENIUQ''组成。这就把程序本身------或者，如果你喜欢，也可以叫做程序的完全复本------给打印出来了。这很像说谎者句子的蒯恩式说法：

``放在其引文形式后面得到假句子''

放在其引文形式后面得到假句子

一定要清楚，上述程序的后一半中引号内所出现的那个印刷符号串------即TEMPLATE的值------决不能解释为一列指令。如果它碰巧是一组指令，在某种意义上，也只是巧合。如上面指出的，它同样也可以是``DOUBLE-BUBBLE''或任何其它的印刷符号串。这个方案的出色之处在于：当同样的符号串出现在程序的前一半中时，也会被看作是一个程序（因为它不在引号之中）。因而，在这个程序中，有一个符号串以两种方式起着作用，首先是作为程序，其次是作为数据。这就是自复制程序的秘密。而且，如我们将要看到的，这也是自复制分子的秘密。顺带说一下，把任何种类的自复制的对象或自复制的实体叫做``自复制''是有益的；类似地，我们也把任何自指的对象或自指的实体称为``自指''。以后我间或要使用这些术语。

可以使用一种并非事先设计得易于书写自复制的语言来写自复制的程序，上面那个程序就是一个极好的例子。这时候，为承担起这个任务，就不得不使用一些被看作是该语言的组成部分的概念和运算------例如``QUOTE-MARK''{[}引号{]}这个词以及``PRINT''{[}打印{]}这个命令。但是，假如把一种语言设计得便于书写自复制，那就能写出一些很短的自复制。比方说，如果``摁\gebfont{㧟}''运算是该语言的内在部分，而不必去做什么显式定义（就像我们对PRINT所做的假定一样），那么，一个极其短小的自复制程序就是

ENIUQ{[}\textquotesingle ENIUQ\textquotesingle{]}

这与乌龟对蒯恩型说谎者自指句子的说法非常相似，那里要假定``\gebfont{㧟}摁''这个动词是已知的：

``被\gebfont{㧟}摁时得到假句子''被\gebfont{㧟}摁时得到假句子。

不过，自复制还能更短，比如，在某些计算机程序中可以有这样的约定：任何以星号打头的程序在正式执行前都要复制一遍。这样，这个仅由一个星号组成的程序就是一个自复制！你可能会抱怨说这很无聊，而且要依赖一个几乎完全是任意的约定。如果真是这样，那你就在重复我先前的观点：利用词组``本句子''以达到自指，几乎是个骗局------它过多地依赖处理者，而不是对自指性的直接揭示。利用一个星号做自复制的例子，就像用``我''字做自指的例子：它们都隐瞒了各自的问题中的一切令人感兴趣的方面。

这使人想起另一种新奇的自复制典型：复印机。可以断定任何手写的文件都是一个自复制，只要把它放进复印机并按一下电钮就行了。但这有点不符合我们的自复制概念。这个过程没有征求过这张纸的意见，因此它并没有指挥它自己的复制过程。于是和前面相同，所有奥妙都在处理机内。在把一个东西叫做自复制以前，我们需要有一种感觉：它应该尽可能明显的包含复制自己的方法。

确实，这种明显性是个程度问题。不过有一个直观的界限，在界限的这一边，是真正的自指挥的自复制过程，而另一边只有那种通过一台一成不变的、自动的复印机所做的复制。

什么是副本？

在对自指和自复制的任何讨论中，我们早晚都得面临``什么是副本''这个问题。在第五章和第六章中我们曾相当严肃地讨论过它，现在再来回到这个问题。为了领略这个问题的风味，我们先来看一些高度荒诞而又似乎有理的自复制例子。

一首自复制的歌

设想在一间酒吧里有一台投币式自动唱机，只要你摁一下按钮7，它就会放出下面这首像情歌一样的曲子：

再投一枚硬币吧，放进这台投币唱机，

我将不断地歌唱，要的仅仅就是按7。

我们可以为此画一个小图（图86）：

\pandocbounded{\includegraphics[keepaspectratio]{./media/OEBPS/Images/img_086.png}}

图86．一首自复制的歌

尽管效果是歌曲复制了自己，但要把这首歌叫做自复制，总感到有些不对头。因为事实是，当它通过``7''时，全部的信息并不在那里。信息之所以得以返回，只是由于信息全部保存在自动唱机内------也就是说，是在图上的一个箭头里，并非那个椭圆内。这首歌是否包含一个有关如何使它自己被重放的完整描述，是大可怀疑的，因为``7''只是一个触发器，而不是一个副本。

一个``螃蟹''程序

下面来考虑一个把自己倒着打印出来的程序（某些读者可能会愿意考虑一下如何利用前面给出的那个自复制模型，用上述类似BlooP的语言写出这样一个程序）。能不能把这样一个滑稽的程序看作自复制呢？在某种意义上说是可以的。因为只要对它的输出做一个不足道的变换就能恢复原先的程序。说这个输出含有和投币唱机程序本身一样的信息，只是简单的重复，似乎也是公平的。然而很明显，某些人也可以眼看着这个输出而看不出它是个倒着打印出来的程序。回忆一下第六章的术语，我们可以说这个输出与这个程序的``内在消息''相同，不过却有不同的``外在消息''------也就是说，必须用不同的解码方式来读它们。现在，如果把外在消息看成信息的一部分------这好像十分合理------那么，全部的信息终归并不相同。所以这个程序不能算是自复制。

但是，这是一个使人不安的结论。因为我们习惯于认为一件东西和它的镜像含有相同的信息。不过，回忆一下在第六章中，我们曾使``固有意义''这一概念依赖于把智能假定为一个普遍观念。其思想是，在确定一个对象的固有意义时，我们可以忽视某些类型的外在消息------那些能够普遍被理解的信息。也就是说，在某种仍有缺陷的意义上，如果解码机制看上去足够基本，那么所要揭示的内在消息就是该考虑的唯一意义。在这个例子中，就似乎有充分的把握猜想一个``标准智能''会认为两个镜像含有彼此相同的信息。这也就是说，它认为两者之间的同构映射十分不足道，以至可以忽略。因而，我们可以认为在某种意义上，该程序算是个十足的自复制。这样的直观是站得住脚的。

说谎者横跨太平洋

自复制的另一个稀奇古怪的例子是一个打印自己的程序，不过是翻译成另一种计算机语言以后再打印。可以把它与下述古怪的蒯恩型说谎者自指的说法作个比较：

``is an expression which, when it is preceded by its translation, placed in quotation marks, into the language originating on the other side of the Pacific Ocean, yields a falsehood''是一段陈述，说的是把它译成太平洋另一侧的语言，然后再加上引号放在前面，就得到假句子。

你可能要试着写写这个怪诞的混和物所描述的句子（提示：它不是它自己------或者说，至少当``它自己''取朴素的意义时，它不是它自己）。如果``逆行自复制''（即倒着写出自身的程序）的概念让人想起螃蟹卡农，那么``翻译的自复制''概念大概会使人想起带有主题转调的卡农。

打印自己的哥德尔数的程序

不打印原程序的严格副本，而打印一份译文，这种想法似乎不着边际。但是，你要是想用BlooP或FlooP写一个自复制程序，你就不得不求助于某种这样的机制，因为这些语言输出的总是一个数，而不是一个印刷符号串。所以，你只好编个程序来打印它自己的哥德尔数：一个巨大的整数，其十进展开式中每三个数字作为一个代码，一个符号一个符号地为该程序编了码。在这个程序力所能及的范围内，它已是尽其所能在打印它自己了：它打印出自己在另一个``空间''中的副本，而且，在整数空间和符号串空间之间的来回转换是很方便的。因此，输出的值不像``7''那样只是起一种触发作用。原程序的全部信息在输出中都是``接近表面''的。

哥德尔式自指

这已经很接近于描述哥德尔自指句子G的机制了。毕竟，那个TNT符号串不是含有一个对它自己的描述，而是含有对一个整数（u的算术\gebfont{㧟}摁化产物）的描述。恰巧这个整数正是符号串G在自然数空间中的一个精确的``镜像''。因此，G指向它自己在另一个空间中的翻译。称G是一个自指符号串不会感到难受，因为这两个空间之间的同构是如此严密，以致我们可以把它们看成是等同的。

把TNT反映到抽象的自然数王国里的这个同构，可以比拟成利用符号把现实世界反映到我们大脑中的那个准同构。符号对于客体充当了准同构的角色，多亏了它们，我们才能思考。类似地，哥德尔数对于符号串充当了同构的角色，多亏它们，我们才能从关于自然数的语句中找到元数学的意义。有关G的一件令人吃惊的、几乎是不可思议的事情是，它使得自指得以实现，而不理会``写出G的语言------TNT------似乎根本没有希望谈论自己的结构''这一事实。这和汉语不一样，用汉语讨论汉语是世上再容易不过的事情了。

因此，G是借助翻译达到自指的一个突出例子------但难说这是最直接的情形。也可以回过头来想想我们的某些对话，其中有些也是利用翻译达到自指的。比如，拿《无伴奏阿基里斯奏鸣曲》来说，在这篇对话里，有几处是指向巴赫的《无伴奏小提琴奏鸣曲》，而乌龟的那个建议------设想一架羽管键琴的伴奏------尤其有趣。说到底，如果把这个思想用于对话本身，就能想出乌龟说的话了；但如果假定阿基里斯部分是独立成章的（像小提琴那样），那么再把什么话归结为乌龟说的，那就十分错误了。无论如何，这里再一次有了一个利用对话到巴赫作品的映射而建立起的自指，这个映射自然就留待读者加以注意了。而且，即使读者没有注意到，映射也依然存在，对话也仍然是个自指。

通过增值达到的自复制

我们已经把自复制和卡农作了对照。那么，增值卡农的类比物该是什么呢？下面是一种可能性：考虑一个含有空循环的程序，这个循环的唯一目的是使该程序变慢，有一个参数用以指明该循环重复的频度。这样就可以作出一个打印自身的副本的自复制，不过要改变一个参量，使得当副本运行时其速度只是它的母程序的速度的一半，而它的``女儿''的速度又是它的一半，如此下去。这些程序中没有一个是严格地打印自身的，不过很明显，它们都属于同一``家族''。

这使人联想起生物的自复制过程。显然任何一个生物体都不全同于父母，那为什么生儿育女这一类事情还叫``自复制过程''呢？答案是，双亲和孩子之间有一个粗略的同构。这是一个保持物种信息的同构。因此，复制的东西是类而不是例。第五章的递归图案``G图''也是如此：在那里，不同大小和形状的``磁性蝴蝶''间的那个映射是很粗糙的。任何两个蝴蝶都不尽相同，但它们都属于同一``物种''，而这个映射恰好保持了这一事实。用自复制程序的话来说，这将与一族程序相对应，这些程序都是用出自同一种计算机语言的各种``方言''写的，每一个程序都能把自己打印出来，但都稍作改变，以使得结果是原先的语言的一种方言。

凯姆式自复制

下面的例子可能是一种最为诡秘的自复制：不是用编译语言写一个合法的表达式，而是敲进一条编译程序所使用的错误信息。当编译程序读到你的``程序''时，它一下子就会陷入混乱。因为你的``程序''不合语法，于是编译程序就打印出一个错误信息。你要做的事情，就是设法安排一下，使得它打印出来的就是你键入的那条错误信息。斯科特·凯姆向我建议的这种自复制，是利用了系统中与你通常接触的层次所不同的另一个层次。尽管看起来似乎不足道，但是，在那些各种自复制为生存而相互争斗的复杂系统中，可能会有类似的情形，这我们马上就要讨论到。

什么是原件？

除了``什么是副本''这个问题，还有一个与自复制有关的基本哲学问题，那就是上述问题的另一面：``什么是原件？''弄清这是怎么回事的最好办法是看几个例子：

⑴．一个程序，如果用一个在计算机上运行的解释程序来解释它，它就把自己打印出来；

⑵．一个程序，如果用一个在计算机上运行的解释程序来解释它，它就把自己连同该解释程序的一个完整副本（毕竟也是个程序）一起打印出来；

⑶．一个程序，如果用一个在计算机上运行的解释程序来解释它，它就不仅把它自己连同该解释程序的一个完整副本都打印出来，而且还指挥着一个机械的装配过程，又装配起一部计算机，与那台运行着该程序及那个解释程序的计算机一模一样。

很明显，在（1）中，该程序是个自复制。而在（3）中，它究竟是个自复制程序，还是程序加上解释程序的混合系统，还是程序、解释程序以及处理机的联合体呢？

显然，一个自复制可以包括比单纯地打印自己更为丰富的内容。事实上，本章余下的大部分篇幅都将用来讨论那些把数据、程序、解释程序、处理机全搅在一起的自复制，而且这种自复制会一下子把它们全都复制出来。

印符遗传学

现在，我们该讨论``生命王国的分子逻辑''这个在二十世纪最吸引人而又意义深远的课题了。这里借用了阿尔伯特·莱宁格尔富有想象力的说法。这也是逻辑------不过比人的智力所能设想的任何一种逻辑都更为复杂，也更为精彩。我们要从一个比较新颖的角度来探讨它：利用一种人工的单人游戏。我把它叫做印符遗传学。在印符遗传学里，我力图在一个印符系统（乍一看很像以WJU系统为代表的那类形式系统）中，捕捉分子遗传学的某些思想。当然，印符遗传学中包含了很多简化工作，从而特别有利于教学目的。

我得赶紧解释一下，分子生物学领域是一个若干层次上的各种现象相互作用的领域，而印符遗传学只打算解释来自一两个层次的现象。具体说，纯粹化学的方面完全被排除了------因为它们低于此处所讨论的层次。与此相似，古典遗传学（非分子遗传学）的一切方面也都排除在外------它们所属的层次比这里讨论的要高。我只想利用印符遗传学给出一种直观，以期说明弗兰西斯·克里克（DNA双螺旋结构的发现者之一）所阐述的、集中在著名的``分子生物学中心法则''里的那些过程：

DNA\gebfont{⇒}RNA\gebfont{⇒}蛋白质

我的希望是，通过我所构造的这个瘦骨伶仃的模型，读者可以了解该领域的某些简单统一的原则------否则，这些原则可能会被很多不同层次上的各种现象之间庞大而复杂的相互作用所淹没。我所牺牲的东西自然是严格的精度，所获得的东西------我希望------是一点洞见。

串、基、酶

印符遗传学游戏是在字母符号的序列上进行印符操作。要用到四个字母：

A　C　G　T

由这样四个字母组成的任意序列叫做串。下面就是一些串：

GGGG

ATTACCA

CATCATCATCAT

顺便提一下，在印符遗传学中，串充当DNA片断的角色（在实际的遗传学中经常就把它们叫做``串''）。``信使RNA''的功能------它在印符遗传学中也用串来表示------是为DNA提供``邮政快递''服务。这将在后面看到。

有时我要把字母A、C、G、T称作基，并把它们所占据的位置称作单元。这样，上述例子中的中间那个串有七个单元，其第四个单元上放着基A。如果你有了一个串，你可以对它进行操作，用各种方式变化它。你还可以靠复制一个串或把一个串切成两半而多造出一些串。某些操作使串变长，某些使其变短，还有一些不涉及串的长度。

操作可以集束进行，也就是几项操作按次序接续执行。这样的一束操作有点像一台程序化的机器在串上移来移去地施加操作。这些活动着的机器就叫做``印符酶''------简称酶。酶在串上每次作用于一个单元，我们把这说成是酶``拴''在此刻所作用的那个单元上。

下面谈谈酶（任选几个）如何作用在具体的符号串上。首先该知道的事情是：每一种酶总喜欢一开始拴在某个特殊的字母上。所以就有四种酶------喜欢A的酶、喜欢C的酶、等等。给出某个酶执行的一系列操作之后，你就可以判断出它喜欢哪个字母了。不过现在我就给出它们，不加解释。下面是一个任选的酶，它由三项操作组成：

\begin{longtable}[]{@{}lll@{}}
\toprule\noalign{}
\endhead
\bottomrule\noalign{}
\endlastfoot
\multirow{3}{=}{\pandocbounded{\includegraphics[keepaspectratio]{./media/OEBPS/Images/Formula-left_bracket.png}}} & ⑴ & 删除拴着这个酶的单元（并把酶拴在右方的下一个单元上）。 \\
& ⑵ & 向右移一个单元。 \\
& ⑶ & （紧挨着这个单元，在右边）插进一个T。 \\
\end{longtable}

这个酶正巧是喜欢一开始拴在A上的。我们用下面这个串作样本，看看酶怎么工作：

ACA

如果我们的酶拴在左边的A上并开始动作，那会发生些什么呢？第一步是删除这个A，于是我们是在CA的左边------然后这个酶就拴在C上了。第二步把酶向右滑，到A上。而第三步就把T加在末端，形成CAT这样一个串，这样这个酶就完成了它的全部职责：把ACA变成了CAT。

如果酶已经拴在ACA中右边的A上，会怎么样呢？它就要删除这个A，并移出该串的端点。只要发生这种情形，酶就停止工作（这是个一般原则）。于是，全部的效果就是砍掉一个符号。

我们来再看几个例子。这又是一种酶：

\begin{longtable}[]{@{}lll@{}}
\toprule\noalign{}
\endhead
\bottomrule\noalign{}
\endlastfoot
\multirow{4}{=}{\pandocbounded{\includegraphics[keepaspectratio]{./media/OEBPS/Images/Formula-left_bracket.png}}} & ⑴ & 寻找本单元右边最近的嘧啶。 \\
& ⑵ & 进入复制状态。 \\
& ⑶ & 寻找本单元右边最近的嘌呤。 \\
& ⑷ & 从此处切断该串（从当时所在单元的右边切）。 \\
\end{longtable}

这里有两个术语``嘧啶''和``嘌呤''。它们很简单：A和G叫嘌呤，C和T叫嘧啶。所以找出一个嘧啶就是找出最近的C或T。

复制状态和双串

还有一个新术语叫复制状态。任何一个串都能``复制''成另一个串，但方式很有趣。不是把A复制成A，你得把它复制成T，反过来也一样。C也不复制成C，你得把它复制成G，反过来也一样。请注意，嘌呤总是复制成嘧啶，而嘧啶总是复制成嘌呤。这叫做补基对。这种互补关系可表示如下：

\begin{longtable}[]{@{}lllll@{}}
\toprule\noalign{}
\endhead
\bottomrule\noalign{}
\endlastfoot
\multicolumn{5}{@{}l@{}}{%
互补} \\
~ & \multirow{3}{=}{\pandocbounded{\includegraphics[keepaspectratio]{./media/OEBPS/Images/Formula-left_bracket.png}}} & A⇔T & \multirow{3}{=}{\pandocbounded{\includegraphics[keepaspectratio]{./media/OEBPS/Images/Formula-right_bracket.png}}} & ~ \\
嘌呤 & & ~ & & 嘧啶 \\
~ & & G⇔C & & ~ \\
\end{longtable}

顺便说一下，从英文看，这就像是阿基里斯{[}Achilles{]}与乌龟{[}Tortoise{]}配对，螃蟹{[}Crab{]}与其基因{[}Genes{]}配对。

``复制''一个串时，并不是实际复制它，而是产生出它的补串，并把它倒过来写在原串的上方。我们来具体看看，让上述的酶作用在下面的串上（这个酶恰好也喜欢从A出发）；

CAAAGAGAATCCTCTTTGAT

有许多地方能当作起点，比方说，我们选取左边第二个A。酶拴在它上面，然后实施步骤1：寻找右边最近的嘧啶。好，这意味着找C或T。第一个是位于该串中央附近的那个T，那就是我们到达的地方。然后，步骤2：复制状态，好，把A倒过来放在我们的T上面。不过这并没有完，复制状态会继续保持有效，直到从其中退出------或者直到酶做完为止。这意味着复制过程继续进行，于是酶所通过的每一个基上面都放上了一个补基。步骤3是说找到我们的T右边的一个嘌呤。那就是从右端数在第三个位置上的G。随着我们移到这个G上，我们一路上都必须``复制''------即创造一个补串。得到的东西是；

\begin{center}

\pandocbounded{\includegraphics[keepaspectratio]{./media/OEBPS/Images/Gene-7.png}}

\end{center}

最后一步是切断这个串，这就得到两段。

\begin{center}

\pandocbounded{\includegraphics[keepaspectratio]{./media/OEBPS/Images/Gene-8.png}}

\end{center}

和

AT

这一束指令做完了。不过，我们得到了一个双串。凡是发生这种事情，我们一律把两个互补的串彼此分开（这是一般原则），所以事实上我们的最终产品是三个串组成的集合：

AT、CAAAGAGGA、CAAAGAGAATCCTCTTTG。

应当注意到，颠倒的符号串已翻过来正面朝上了，于是左和右也颠倒。

现在，我们已经见过可以施行在串上的各种印符操作中的大部分了。还有两条命令得说一说。一条是退出复制状态，另一条是把酶从一个串转移到倒着放在它上面的串上。这样做的时候，如果你保持纸的方向不变，那就必须把全部命令中的``左''和``右''互换。更好的做法是，保持措辞不变而把纸转过来，使放在上面的符号串成为可读的。如果已经给出了``转移''指令，而当时拴着酶的地方没有补基，那么酶就简单地自行脱离该串，其工作也就结束了。

还应说说，当遇到一条``切断''命令时，这与两个串（如果有两个串的话）都有关，然而``删除''只与酶所作用的那个串有关。如果复制状态在持续，那么``插入''指令就与两个串有关------基自身插入酶所工作的串，其补基插入另一个串。而在非复制状态下，``插入''指令就只与原串相关了，所以必须在其补串上插入一个空格。

只要复制状态在持续，``移动''和``寻找''指令就要求我们给滑动着的酶所碰到的所有基都各造一个补基。顺便说一句，酶开始工作时，复制状态总是关闭的。如果本不处于复制状态，又遇见``退出复制状态''的指令，就什么也不会发生。类似地，如果复制状态已在持续，又遇到``进入复制状态''的指令，那也是什么事都不发生。

氨基酸

共有十五种类型的指令，开列如下：

\begin{longtable}[]{@{}lll@{}}
\toprule\noalign{}
\endhead
\bottomrule\noalign{}
\endlastfoot
cut & ------ & 切断串 \\
del & ------ & 从串里删除一个基 \\
swi & ------ & 把酶转移到另一个串上 \\
mvr & ------ & 右移一个单元 \\
mvl & ------ & 左移一个单元 \\
cop & ------ & 进入复制状态 \\
off & ------ & 退出复制状态 \\
ina & ------ & 在本单元右侧插入A \\
inc & ------ & 在本单元右侧插入C \\
ing & ------ & 在本单元右侧插入G \\
int & ------ & 在本单元右侧插入T \\
rpy & ------ & 寻找右边最近的嘧啶 \\
rpu & ------ & 寻找右边最近的嘌呤 \\
lpy & ------ & 寻找左边最近的嘧碇 \\
lpu & ------ & 寻找左边最近的嘌呤 \\
\end{longtable}

每一类指令都有一个用三个字母缩写的名字。我们把这些用三个字母缩写的指令叫作氨基酸。于是，每个酶都由一系列氨基酸制成。我们随便写一个酶：

rpu-inc-cop-mvr-mvl-swi-lpu-int

和一个任意的串：

TAGATCCAGTCCATCGA

看看这个酶如何对这个串起作用。碰巧，这个酶只拴在G上，我们就把它拴在中间的G上，并开始工作，向右寻找一个嘧啶（即A或G）。我们（即酶）跳过TCC，然后停在A上。插入一个C，于是有

\begin{center}

\pandocbounded{\includegraphics[keepaspectratio]{./media/OEBPS/Images/Gene-1.png}}

\end{center}

其中的箭头指着酶所拴的单元。进入复制状态。于是在C上放一个倒置的\gebfont{G}。右移，再左移，然后转移到另一个串上。迄今为止我们已有

\begin{center}

\pandocbounded{\includegraphics[keepaspectratio]{./media/OEBPS/Images/Gene-2.png}}

\end{center}

把它倒转过来，让酶拴在下面的串上：

\begin{center}

\pandocbounded{\includegraphics[keepaspectratio]{./media/OEBPS/Images/Gene-4.png}}

\end{center}

现在我们往左寻找一个嘌呤，于是找到A。复制状态仍在持续，不过其补基已在那里，用不着再做什么。最后我们（在复制状态下）插入一个T。停止：

\begin{center}

\pandocbounded{\includegraphics[keepaspectratio]{./media/OEBPS/Images/Gene-5.png}}

\end{center}

我们最终的产品是两个串：

ATG和TAGATCCAGTCCACTCGA

原先那个串当然就消失了。

翻译和印符遗传密码

现在你大概很想知道这些酶和串都是哪来的，以及怎样知道一个给定的酶最初该拴在哪里。一个可能的方案是把某些随机的串和某些随机的酶放到一起，看看这些酶作用于这些串及其后代时会发生什么事情。这有点WU谜题的味道了，在那里是几条给定的推理规则和一条作为出发点的公理。唯一的不同是，这里对一个串每作用一次，其原先的形式就消失了，而在WU谜题中，规则作用于WJ得到WJU时，并不清除WJ。

不过，印符遗传学就像实际的遗传学一样，其方案颇有些诡秘。我们从某个任意的串出发，这多少有点象形式系统中的一条公理。但是一开始时我们并没有``推理规则''------也就是说并没有酶。然而，我们可以把每个串翻译成一种或多种酶！这样，串自己就会口授将要在它们身上施行的操作，而这些操作又会依次产生出新的串，新串又会口授进一步的酶，如此等等！这简直是登峰造极的层次混合！为了对比，想想WU谜题，若能把每次生成的新定理都通过某种编码转换成一条新的推理规则，那么WU谜题会是多么的不同！

这种``翻译''是怎么做的呢？这涉及一种印符遗传密码。根据这种密码，一个串中的各对相邻的基------称作``二元组''------表示各种不同的氨基酸。总共有十六种可能的二元组：AA、AC、AG、AT、CA、CC、等等，而氨基酸共有十五种。印符遗传密码如图87所示。

\pandocbounded{\includegraphics[keepaspectratio]{./media/OEBPS/Images/img_087.png}}

图87．印符遗传密码

串中各个二元组按这种对应关系分别给十五种氨基酸（和一个标点符号）编码。

按照这个表，二元组GC的翻译是``inc''（``插入C''），AT的翻译是``swi''（``转移''），等等。这样，一个串显然可以十分直截了当地口授一种酶。例如，串

TAGATCCAGTCCACATCGA

断成二元组，就是

TA GA TC CA GT CC AC AT CG A

有一个A留在边上。翻译成酶，是

rpy - ina - rpu - mvr - int - mvl - cut - swi - cop

（注意：剩下的A没什么用处。）

酶的三级结构

每个方格中右下角的字母``s''，``l''和``r''是干什么的？在确定酶必须拴在哪里时，它们是关键，而且其行为方式也很奇特。要弄清一个酶喜欢拴在哪个字母上，就得弄清这个酶的``三级结构''。这是由其``一级结构''所决定的。所谓``一级结构''指的是其氨基酸序列，所谓``三级结构''指的是它喜欢的``折叠''方式。要紧的是，酶不喜欢像我们前面展示的那样成为一条直线。每个内部的氨基酸（两头的氨基酸除外），都可能会是一个``转弯''，这由角上的字母来决定。具体说。``l''和``r''分别代表``左''和``右''，而``s''代表``直''。下面我们拿刚才的那个作样本的酶为例，让它自行折起来以显示其三级结构。我们从酶的一级结构出发，顺着它从左向右移动。在角上有``l''的酶上，我们向左转；有``r''的，向右转；而对于``s''，则不转弯。图88显示了我们这个酶的二维造型。

\begin{longtable}[]{@{}lllllll@{}}
\toprule\noalign{}
\endhead
\bottomrule\noalign{}
\endlastfoot
cop & & & & & & \\
\gebfont{⇑} & & & & & & \\
swi & \gebfont{⇐} & cut & \gebfont{⇐} & mvl & \gebfont{⇐} & int \\
& & & & & & \gebfont{⇑} \\
& & & & & & mvr \\
& & & & & & \gebfont{⇑} \\
& & rpy & \gebfont{⇒} & ina & \gebfont{⇒} & rpu \\
\end{longtable}

图88．印符酶的三级结构

注意在``rpu''处的向左转、在``swi''处的向右转，等等。还要注意第一段（``rpy\gebfont{⇒}ina''）与最后一段（``swi\gebfont{⇒}cop''）是正交的。这就是决定拴在哪里的关键。事实上，酶的三级结构中，第一段和最后一段的相对方向决定了该酶拴于何处。我们总是可以调整酶的方位，使其第一段指向右方，如果这样做了，那么最后一段就决定了酶所拴的位置，如图89所示。

\begin{longtable}[]{@{}lll@{}}
\toprule\noalign{}
\endhead
\bottomrule\noalign{}
\endlastfoot
第一段 & 　最后一段　 & 所拴的字母 \\
\gebfont{⇒} & \gebfont{⇒} & A \\
\gebfont{⇒} & \gebfont{⇑} & C \\
\gebfont{⇒} & \gebfont{⇓} & G \\
\gebfont{⇒} & \gebfont{⇐} & T \\
\end{longtable}

图89．印符酶所拴的位置表

所以，我们现在所讨论的是一个喜欢字母C的酶。如果在折起的时候，一个酶自己居然会交叉，那也没什么不好------只要把它看成从自己下方或上方过去的就行了。应当注意，在确定酶的三级结构时，全部氨基酸都起了作用。

标点、基因与核糖体

还有一件事有待澄清：印符遗传密码的AA一格为什么是空白？答：二元组AA在串里起标点符号的作用，它标志着一个酶的密码的结束。就是说，如果一个串里有一个或几个二元组AA的话，这个串就在给两个或更多的酶编码。例如，串

CG GA TA CT AA AC CG A

就是为两个酶编码，它们是

cop - ina - rpy - off

和

cut - cop

AA把该串分成两个``基因''。基因的定义是：一个串中给一个酶编码的那个部分。要注意，在一个串内有AA并不就意味着该串是给两个酶编码。比如，CAAG给``mvr - del''编码。这里的AA从偶数单元开始，因而并不读成二元组！

读串并产生出它们编码的酶，这种装置称作核糖体（在印符遗传学中，做游戏的人起着核糖体的作用）。核糖体一点也管不着酶的三级结构，因为一旦造出了一级结构，三级结构就完全确定了。顺便说一句，翻译的过程总是从串到酶，而从不会反过来进行。

谜题：印符遗传学的自复制

印符遗传学的规则都已给出了，你如果试验一下这个游戏，就会发现它是很有趣的。尤为有趣的是设计一个自复制的串，下面我们就来做这件事。写下一个串，核糖体作用在它上面，就产生出在该串内编码的所有的酶。然后让这些酶接触原先的串，并在上面工作。这就得到一批``子串''。那些子串本身又通过那些核糖体，得到第二代酶，它们又作用在其子串上，如此循环往复。这个过程可以进行任意多阶。我们的希望是：最终，在某一阶段出现的串里，能找到最初那个串的两个副本（副本之一事实上可能就是原先的串）。

印符遗传学的中心法则

\pandocbounded{\includegraphics[keepaspectratio]{./media/OEBPS/Images/img_090.png}}

图90．印符遗传学的中心法则

``缠结的层次结构''的一个例子

印符遗传过程可以用一个图（图90）来表示。

这个图说明了印符遗传学的中心法则。它显示了串如何（借助于印符遗传密码）定义酶，以及酶反过来如何作用于那些导出它们的串，从而得到新的串。因此，左边一条线绘出了老信息怎样流上去，因为酶是串的一个翻译，从而含有和串相同的信息，只不过形式不同，特别是作用的形式不同。然而右边那条线并不表示信息流下来，而是表示如何通过在串中调动符号，创造新的信息。

印符遗传学中的酶就像形式系统中的推理规则，它们闭着眼睛调动串中的符号，而不顾有可能潜藏在符号内的任何``意义''。所以这里就有一个怪异的层次混合状态。一方面，串是被作用的，因而充当了数据的角色（如图90中右边的箭头所指出的）；另一方面，它们又支配这个施加在数据上的作用，因而又充当了程序的角色（如左边的箭头所指出的）。当然，起解释程序和处理机作用的是做游戏的人。联结``上''\,``下''两个印符遗传层次的这个双向通道表明：事实上，不管是串还是酶，哪个也不能认为比另一个高。作为对照，WJU系统的中心法则图示是：

\begin{longtable}[]{@{}ll@{}}
\toprule\noalign{}
\endhead
\bottomrule\noalign{}
\endlastfoot
& \\
推理规则 & \\
↓ & （印符操作） \\
符号串 & \\
\end{longtable}

在WJU系统中，层次有明显的区别：推理规则绝对属于比符号串高的层次。对于TNT以及所有的形式系统，情况都类似。

怪圈、TNT及实际的遗传学

但是，我们已经看到，在另一个意义上，TNT中的各层次是混合在一起的。事实上，语言和元语言之间的界限已经打破了：谈论该系统的句子在该系统内部有一个镜像。这表明，如果我们画一个表示TNT与其元语言之间的关系的图，那它就将以一种十分醒目的方式相似于表示分子生物学中心法则的图。事实上，详细做这个对比正是我们的目标。不过，要想这样做，就得指出在哪些地方印符遗传学与真正的遗传学一致，哪些地方又不一致。当然，实际的遗传学远比印符遗传学复杂------不过，读者把在了解印符遗传学的过程中所掌握的``概念骨架''，作为实际遗传学迷宫中的向导，还是十分有用的。

DNA与核苷酸

我们从讨论``串''和DNA之间的关系入手。缩写``DNA''表示``脱氧核糖核酸''。大多数细胞的DNA位于细胞核内，有一层膜保护着。冈瑟·斯坦特曾把细胞核描述成是细胞的``金銮殿''，DNA就充当着里面的君主。DNA由比较简单的、称为核苷酸的分子的长链组成，每一个核苷酸由三个部分构成：（1）脱去一个特殊氧原子的、称为``核糖''的一种醣，所谓``脱氧''就是指这个；（2）磷酸基；以及（3）一个基。区别各种核苷酸的东西只是基，因而完全可以用基来鉴别核苷酸。出现在DNA核苷酸中的四种类型的基是：

\begin{longtable}[]{@{}lllllll@{}}
\toprule\noalign{}
\endhead
\bottomrule\noalign{}
\endlastfoot
A：腺嘌呤 & \multirow{3}{=}{\pandocbounded{\includegraphics[keepaspectratio]{./media/OEBPS/Images/Formula-right_bracket.png}}} & ~ & \multirow{3}{=}{　} & G：鸟嘌呤 & \multirow{3}{=}{\pandocbounded{\includegraphics[keepaspectratio]{./media/OEBPS/Images/Formula-right_bracket.png}}} & ~ \\
~ & & 嘌呤 & & ~ & & 嘧啶 \\
T：胸腺嘧啶 & & ~ & & C：胞嘧啶 & & ~ \\
\end{longtable}

\pandocbounded{\includegraphics[keepaspectratio]{./media/OEBPS/Images/img_091.png}}

图91．DNA的四种基

腺嘌呤，鸟嘌呤，胞嘧啶，胸腺嘧啶{[}摘自哈那瓦尔特和海因斯{[}Hanawalt; Haynes{]}，《生命的化学基础》{[}The Chemical Basis of Life{]}，旧金山：W. H. Freeman，1973年版，第142页{]}

一个DNA串就这样由彼此串在一起的许多核苷酸组成，像一串珍珠一样。使一个核苷酸与其两边邻居相连的化学键很强，称作共价键，而``珍珠串''通常叫做DNA的共价脊柱。

\pandocbounded{\includegraphics[keepaspectratio]{./media/OEBPS/Images/img_092.png}}

图92．DNA结构像一架梯子

边上两条由脱氧核糖和磷酸盐的单元交替组成，横档以特殊方式配对------A配T，G配C------的基形成，基之间分别由两个或三个氢键相互连接。{[}摘自哈那瓦尔特和海因斯，《生命的化学基础》第142页{]}

DNA通常由一对串构成------就是说，由两条逐个配对的核苷酸串组成（见图92）。站在两个串之间负责这种特殊类型的配对的是基。这个串中的每个基面对着另一个串中的补基，并拴住它。互补关系就像印符遗传学中一样：A与T配对，C与G配对，总是一个嘌呤配一个嘧啶。

与顺着脊柱方向的强共价键相比，串之间的键是相当弱的。它们不是共价键，而是氢键。氢键出现时，两个分子络合物以下面的方式相互结合：原先属于其中一个分子的某个氢原子，变得``弄不清''属于哪边了。它游移于两个络合体之间，不知加入哪一边好。由于双串型DNA的两个部分只靠氢键彼此相联，所以比较容易把它们分开，也比较容易结合。这一事实对细胞的工作至关重要。

DNA形成双串时，两个串彼此缠绕，像扭在一起的藤条（图93）。每缠一圈恰好有十个核苷酸对，换句话说，每个核苷酸``扭''36度。单串型DNA就不是这种盘绕状，因为这是基对的产物。

\pandocbounded{\includegraphics[keepaspectratio]{./media/OEBPS/Images/img_093.png}}

图93．DNA的双螺旋分子模型

摘自维农·英格莱姆{[}Vernon M. Ingram{]}，《生物合成》{[}Biosynthesis{]}，加利福尼亚州，Menlo Park：W. A. Ben-jamin，1972年版，第13页

信使RNA与核糖体

如上所述，在很多细胞里，DNA这个细胞君主住在隐蔽的``金銮殿''------细胞核------之内。但细胞中的大多数``活物''都在核外，即在细胞质中------那个相对于细胞核``图形''而言的``衬底''。特别是实际主宰着每个生命进程的酶，是在细胞质中由核糖体产生的，而且酶的大多数工作都在细胞质中进行。如同印符遗传学一样，全部酶的蓝图都储存在串里------即在DNA里，而DNA又藏在小小的细胞核之中。那么有关酶的结构的信息是怎样从细胞核里传给核糖体的呢？

这就该引入信使RNA------mRNA了。早先，mRNA曾被戏称为一种DNA``邮政快递''服务机构。这并不意味着mRNA在物理意义上把DNA带到什么地方去，而是说它把储存在细胞核内的DNA中的信息或命令，传递给细胞质中的核糖体。这是怎么做的呢？很简单：细胞核内部一种特殊类型的酶把DNA基序列中的一个个长段忠实地复制到一个新串------信使RNA串------上。然后，这个mRNA就离开细胞核，漫步走入细胞质，在那里它又跑进很多核糖体内，这些核糖体就开始根据它来制造酶。

在细胞核内把DNA复制到mRNA上的过程叫转录。在这一过程中，双串型的DNA必须暂时拆成两个单串，其中的一个用作mRNA的模板。顺便说说，``RNA''表示``核糖核酸''，它很像DNA，只不过其全部核苷酸里的核糖中都有DNA核苷酸所缺的那个特殊的氧原子，所以去掉了``脱氧''这个词头。而且RNA使用尿嘧啶这个基，而没有胸腺嘧啶，所以RNA串中的信息可以用``A''、``C''、``G''、``U''四个字母的任意序列来表示。从mRNA转录出DNA时，转录过程可以通过常用的基配对法来进行（不过这时不再有T了，而是U），所以，这种DNA模板与其mRNA伴侣看上去有点像是这样：

DNA:（模板↓）

\ldots\ldots CGTAAATCAAGTCA\ldots\ldots{}

\ldots\ldots GCAUUUAGUUCAGU\ldots\ldots{}

mRNA:（副本↑）

RNA并不总与自己形成长的双串，尽管它能这样做。所以一般说来，它不是处于DNA的那种螺旋状态中，而是一个长长的、有时带点弯的串。

一个mRNA串一旦离开细胞核，就会遇到那些奇妙的亚细胞物质------``核糖体''。不过，在继续解释核糖体如何使用mRNA之前，我想对酶和蛋白质做些说明。酶属于分子生物学中所谓的蛋白质这样一个一般性范畴，而核糖体的工作是制造所有的蛋白质，并不只是制造酶。不是酶的那些蛋白质多是一些不活跃的家伙，例如，它们中有很多是结构分子，这意味着它们像是建筑物中的梁、柱等等：它们把细胞的各个部分连在一起。还有其他一些种类的蛋白质。不过就我们的目的而言，主要的蛋白质是酶，今后就不再特别区分了。

氨基酸

蛋白质由一系列氨基酸组成。氨基酸共有二十个品种，每一种都由三个字母来表示：

\begin{longtable}[]{@{}lll@{}}
\toprule\noalign{}
\endhead
\bottomrule\noalign{}
\endlastfoot
ala & ------ & 丙氨酸 \\
arg & ------ & 精氨酸 \\
asn & ------ & 天门冬氨酰胺 \\
asp & ------ & 天门冬氨酸 \\
cys & ------ & 胱氨酸 \\
gln & ------ & 谷酰胺 \\
glu & ------ & 谷氨酸 \\
gly & ------ & 甘氨酸 \\
his & ------ & 组氨酸 \\
ile & ------ & 异白氨酸 \\
leu & ------ & 白氨酸 \\
lys & ------ & 赖氨酸 \\
mef & ------ & 蛋氨酸 \\
phe & ------ & 苯基丙氨酸 \\
pro & ------ & 脯氨酸 \\
ser & ------ & 丝氨酸 \\
thr & ------ & 苏氨酸 \\
trp & ------ & 色氨酸 \\
tyr & ------ & 酪氨酸 \\
val & ------ & 缬氨酸 \\
\end{longtable}

应当注意到这里与印符遗传学在数目上的轻微差别，在那里，我们只有十五种用以组成酶的氨基酸。氨基酸是大体上与核苷酸同样复杂的一个小分子，因而，建造蛋白质的砖块与建造核酸（DNA、RNA）的砖块差不多大。但是蛋白质是由短得多的序列组成的：在典型的情形，大约三百个氨基酸就组成一个完整的蛋白质，而一个DNA串却可以含有数十万乃至上百万个核苷酸。

核糖体与录音机

当一个mRNA串游离出来进入细胞质，然后遇到一个核糖体时，会发生一个极其复杂而又精彩的过程，我们称之为翻译。可以说，这个翻译过程就是一切生命的根本所在，其上罩着许多神秘的光环。不过，从本质上说，描述它并不难。我们先做一个形象化的比喻，然后再比较严格地加以说明。把mRNA设想成一长条录音磁带，而核糖体就像一台录音机。当磁带通过录音机的放音磁头时，信号就被``读出''并转变成音乐或其它声响。就这样，磁记号``翻译''成了音符。与此类似，当mRNA``磁带''通过核糖体的``放音磁头''时，产生出的音符是氨基酸，而作出的``乐曲''是蛋白质。所谓翻译也就是这些。图96表明了这一点。

遗传密码

\pandocbounded{\includegraphics[keepaspectratio]{./media/OEBPS/Images/img_094.png}}

图94．遗传密码

信使RNA串中的每个三元组就这样为二十种氨基酸（及一个标点符号）编码

可是，一个核糖体在读到一系列核苷酸时，怎么能生产出一系列氨基酸呢？经过很多人的努力，这个秘密在六十年代初被揭开了。答案的核心是遗传密码------从核苷酸的三元组到氨基酸的映射（见图94）。这与印符遗传密码在精神上极为相似，除了这里是三个连续的基（或核苷酸）形成一个密码子，而在那里只要两个基就行了。因而，现在的表上必定要有4x4x4（即64）个不同条目，而不只是十六个。核糖体一次从RNA串上切下三个核苷酸------也就是说，一次切一个密码子------而且，它每这样做一次，就把一个单独的新氨基酸加到了正在生产的蛋白质上。于是，蛋白质就一个氨基酸接一个氨基酸地从核糖体中制造出来了。

\begin{center}

\pandocbounded{\includegraphics[keepaspectratio]{./media/OEBPS/Images/Gene-6.png}}

\end{center}

三级结构

然而，随着一个蛋白质从核糖体中走出来，它不仅会越来越长，而且会逐渐地折叠成一个奇特的三维形状，很像是那种滑稽的蛇形焰火，在点亮时一面伸长，一面卷曲。这种形状称作蛋白质的三级结构（图95），而原先的氨基酸序列叫做蛋白质的一级结构。这个三级结构暗含在一级结构之中，这正和印符遗传学的情形一样。

\pandocbounded{\includegraphics[keepaspectratio]{./media/OEBPS/Images/img_095.png}}

图95．肌红蛋白的结构

大尺寸的``扭曲管道''形状是三级结构，内部的细螺旋------``阿尔法螺旋''------是二级结构。{[}选自莱宁格尔，《生物化学》{]}

可是，如果你只知道一级结构，那么导出三级结构的方法可就远比印符遗传学要复杂了。事实上，搞清某些规则，以便只要知道了蛋白质的一级结构就能指出其三级结构，乃是当代分子生物学中悬而未决的问题之一。

蛋白质功能的简化论解释

印符遗传学和真正遗传学之间的另一个差别------这可能是所有差别中最重要的一个------是：在印符遗传学中，组成酶的每个氨基酸都负责某个具体的``作用片断''；而在实际的酶中，单个的氨基酸不能充当这种明确的角色。确定酶的作用方式的，是作为一个整体的三级结构。不能说``这个氨基酸的存在就意味着如此这般的一个操作得以执行''这类的话。换言之，在实际遗传学中，一个单独的氨基酸对酶的整个功能的贡献不是``上下文无关''的。不过这一事实无论如何都不能当作反简化论者的武器以得出``整体{[}的酶{]}不能解释为其各部分之和''的结论。这种结论还完全没有得到验证。得到验证的东西是排除了下面这个简单的断言：``每个氨基酸都用一种不依赖其它氨基酸的存在的方式对总体发生作用''。换句话说，酶的功能不能看成是由各个部分的上下文无关的功能建立起来的，而必须考虑各部分之间的相互作用。原则上，仍有可能写出一个计算机程序把蛋白质的一级结构作为输入，然后先确定其三级结构，再确定酶的功能。这将是对蛋白质工作的一个彻底的简化论解释，但要确定各部分的``总和''得需要一个高度复杂的算法。在给出酶的一级的、乃至三级的结构之后，如何阐明酶的功能是现代分子生物学的又一个大问题。

也许，整个酶的功能最终仍可以看成是按照一种上下文无关的方法，从各个部分的功能上建立起来的，不过这里得把这些部分看成是单个的粒子，诸如电子、质子之类，而不是看成像氨基酸那样的``大块头''。这正是``简化主义者的窘境''：为了用上下文无关的总和这类概念来解释每一件事情，就必须降到物理的层次上，可是此时粒子的数目会如此之大，以至于只能说出理论上``大概''如何之类的话。所以我们只好满足于一个上下文有关的总和。它又有两个缺点，第一，总和的各个部分是相当大的单元，其性状就只能在高层次上描述，因而并不确定。第二，``总和''这个词意味着：每个部分都分担一个简单的功能，而整体的功能正好是这些个别功能的一个上下文无关的总和。当我们在给定一个酶的各个氨基酸后要解释整个酶的功能时，却正好做不到这一点。但不论是吉是凶，这是解释复杂系统时总要发生的一般现象。为了对各部分之间如何相互作用有一个既直观又易于处理的理解------简言之，为了能继续下去------我们时常不得不牺牲掉依靠微观的及上下文无关的画面所能得到的精确性。这纯粹是因为它不易处理。但此时不必失去信心，以致怀疑这样一个解释在原则上是否存在。

转移RNA与核糖体

现在我们回到核糖体、RNA和蛋白质。前面已经说过，蛋白质是核糖体根据DNA的信使RNA从DNA的``金銮殿''带来的蓝图生产的。这似乎蕴涵着核糖体可以把密码子的语言翻译成氨基酸的语言，而这又等于说核糖体``懂得''密码。然而，全部信息恰恰并不在核糖体内。它是怎么做到的呢？遗传密码贮存在哪里呢？令人惊异的事实是：遗传密码就在------还能在哪儿呢？------DNA自身之内。这当然要作些解释。

我们把总体的解释丢开一会儿，先给出一个部分的解释。在任何给定的时刻，细胞质周围都四散漂浮着大量的四叶草形分子，在它们的一个叶片上，松松地拴着（即靠氢键相连）一个氨基酸，而在对面的叶片上有一个称为反密码子的核苷酸三元组。就我们的目的而言，另外两个叶片无关紧要。下面说明核糖体在制造蛋白质时如何使用这些``草''。当一个新的mRNA密码子进入核糖体的``放音磁头''的位置时，核糖体就伸手到细胞质中抓住一片其反密码子恰与该mRNA密码子互补的草，然后把这片草拽到这样一个位置上：能揭下叶子上的氨基酸，用一个共价键把它粘到正在生成的蛋白质上（顺便说一句，一个氨基酸在蛋白质中与其邻居之间的键是一种很强的共价键，称为``肽键''。因此，蛋白质有时也叫``多肽''）。当然，``草''上带有合适的氨基酸并不是偶然的，因为它们全是根据来自``金銮殿''的严格指令生产出来的。

\pandocbounded{\includegraphics[keepaspectratio]{./media/OEBPS/Images/img_096.png}}

图96．一段mRNA通过核糖体

漂浮在附近的是些tRNA分子，它们带有能被核糖体抓来并加进正在生成的蛋白质中的氨基酸。遗传密码作为一个整体包含在tRNA分子内。要注意，基对（A-U, C-G）在图中由卡在一起的字型表示。{[}斯科特·凯姆绘{]}

这种草的真正名字叫转移RNA（tRNA）。tRNA分子相当小------大约相当于一个很小的蛋白质分子的尺寸------含有大约八十个核苷酸组成的一个链。与mRNA类似，tRNA分子由主要的细胞模板DNA转录出来。可是，与含有数千个核苷酸形成的大长链的庞大mRNA分子相比，tRNA是个侏儒。此外，tRNA在这方面更像蛋白质（而不像mRNA串）：它们都有确定的、良定义的三级结构------由其一级结构确定。一个tRNA分子的三级结构恰好允许一个氨基酸拴在它的氨基酸位置上。确实，它就是对面叶片上的反密码子根据遗传密码表指出的那个氨基酸。tRNA功能的一个生动的形象化类比，是一些漂浮在一位同声翻译周围的云雾之中的单词卡片，只要他需要翻译某个词，就从空中抓来一片------百发百中！在这种情形中，这位翻译就是核糖体，词就是密码子，译文就是氨基酸。

为了让核糖体解出DNA的内在消息，这些tRNA单词卡片就必须漂浮在细胞质中。在某种意义上，这些tRNA包含了DNA的外在消息的本质，因为它们是翻译过程的关键。但是它们自身又都来自DNA。因而外在消息就要通过一种方法成为部分的内在消息，这种方法能使人联想起那个装在瓶子里的、说明它自己是用哪种语言写下来的消息。自然，这样的努力都不可能完全成功：没有一种方法能使DNA揪着鞋带把自己举起来。细胞里事先一定已经有某些遗传密码的知识了，所以才能制造那些酶，把tRNA作为DNA的主要副本给转录出来。而这种知识又储存在早先制造好的那些tRNA分子之内。这种完全排除对外在消息需要的努力就像艾舍尔的那条龙，在困住它的那个二维世界的环境中，它竭尽全力想要成为三维的，它似乎做了很多------不过当然决不会成功，尽管它给出了三维性的一个很好的模拟。

标点和阅读框架

核糖体怎么会知道什么时候一个蛋白质就制作好了呢？和印符遗传学一样，mRNA中也有一个记号，指示出一个蛋白质的结束或开始。事实上，有三个特殊的密码子------UAA、UAG、UGA------都起标点符号的作用，而不是给氨基酸编码。当这样的一个三元组一下一下地滑入核糖体的``放音磁头''时，核糖体就放走正在构造的那个蛋白质，然后开始构造一个新的蛋白质。

最近，已经揭示了已知的最小病毒φX174的全部基因组。一项极为意外的发现正在完成之中：它的9个基因中，有一些是部分重叠的------也就是说，两个不同的蛋白质由同一个DNA段编码！甚至有一个基因整个地包含在另一个里面！之所以如此，是因为两个基因的阅读框架彼此恰好错开一个单元。按这种方式压缩信息所导致的稠密性是难以置信的。当然，这正是《音程增值的卡农》里藏在阿基里斯的福气小甜饼中那奇妙的``4/17俳句''后面的灵感。

概括

于是，简而言之，我们眼前是这样一幅画面：从DNA御座出发，DNA把信使RNA的一些长串送到细胞质中的核糖体上；而核糖体利用盘旋在自己周围的那些tRNA``单词卡片''，根据mRNA中含有的蓝图，一个氨基酸接一个氨基酸地、高效率地构造蛋白质。DNA仅指示了蛋白质的一级结构。但这足够了，因为当蛋白质从核糖体中出来时，它们便``魔术般地''折叠成复杂的形态，从而有能力像一台大功率的化学机器一样行动。

蛋白质与音乐中的多层结构和意义

我们已经把核糖体的形象比作录音机，mRNA比作磁带，而蛋白质则是音乐。这看上去像是信手掂来，然而它们之间却有一些精彩的平行之处。音乐不仅仅是线性的音符序列。我们的心智是在比这高得多的层次上来理解乐曲的。我们把音符组块化成为乐句，若干乐句构成旋律，一些旋律形成乐章，几个乐章组成全曲。类似地，蛋白质只有在作为组块化单元起作用时才有意义。尽管一级结构已经携带着制造三级结构所需要的全部信息，但仍然``嫌少''，因为只有在物理上实际造出这个三级结构来，才能理解其潜力。

顺带提一句，我们只说到一级结构和三级结构，读者可能会奇怪：二级结构哪去了？其实是有的，就连四级结构也有。蛋白质的折叠现象不止存在于一个层次上。具体讲，在氨基酸链的某些地方，有形成一种螺旋的趋势，这叫阿尔法螺旋（别与DNA双螺旋搞混）。蛋白质的这种螺旋形扭曲处在比其三级结构低的层次上。这一层结构可以在图95中看到。四级结构可以比拟成由独立乐章组成的整个乐曲大厦，因为它牵涉到把已经绽开三级鲜花的若干不同的多肽组合成一个更大的结构。这些独立的链之间通常由氢键相连，而不是靠共价键。当然，这正像一个由若干乐章组成的乐曲，其各章之间的联系远不如章内的联系紧密，但仍能形成一个``有机''的整体。

一级结构、二级结构、三级结构和四级结构，这四个层次还可以比作《前奏曲，蚂蚁赋格》中的``无之图''（图60）的四个层次。总的结构------由两横一撇及一个竖弯钩组成一个``无''字------是它的四级结构；那两横及那一撇一竖弯钩本身是三级结构，分别由``整体论''和``简化论''组成；而三级结构的``整体论''与``简化论''其笔划分别由二级结构的``简化论''和``整体论''组成，最后，一级结构则是一遍又一遍地重复``无''字。

多核糖体与二排卡农

我们再来看看把磁带翻译成音乐的录音机，与把mRNA翻译成蛋白质的核糖体两者之间的另一个优美的平行关系。设想有许多磁带录音机一字排开，间隔相等。我们可以把这个阵势叫做``多录音机''。再设想同一条磁带陆续通过各个录音机的放音磁头，如果磁带上录有一个完整的长旋律，那么输出自然是一个多部轮唱，各声部之间的延迟由磁带从一台录音机起到下一台录音机所用的时间来确定。在细胞中，确实也存在这种``分子卡农''，那里，很多核糖体一字排开，形成所谓多核糖体。它们全都``使用''同一个mRNA串，错开一定时间，生产同样的蛋白质（见图97）。

\pandocbounded{\includegraphics[keepaspectratio]{./media/OEBPS/Images/img_097.png}}

图97．多核糖体

同一个mRNA串通过一个接一个的核糖体，就好像磁带依次通过若干台录音机。其结果是一批处于不同程度的正在生成的蛋白质：类似于相互错开的若干录音机所生成的音乐卡农。{[}选自莱宁格尔，《生物化学》{]}

不仅如此，大自然做得还要好些。我们还记得，mRNA是通过转录DNA而制成的，负责这个过程的酶叫RNA聚合酶。常有这样的事发生：一系列RNA聚合酶会并排工作在同一个DNA串上，结果就造出很多彼此分离（但却相同）的mRNA串，每一串都相对上一串延迟一点时间，这正是DNA从一个RNA聚合酶移到下一个聚合酶所需要的时间。同时，在并排出现的每一个mRNA上，又可以有若干不同的核糖体在工作。这样，就达到了一个双层的或二排的``分子卡农''（图98）。在音乐中的对应形象是一个十分荒唐但却很有意思的情景：若干名誊写员同时工作，每个人都把一份原稿从长笛手无法演奏的谱表抄成一个他们能演奏的谱表。每个誊写员抄完一页原稿，就把原稿传给下个人，然后自己再抄新的一页。与此同时，根据出自每个誊写员笔下的抄件，有一组长笛在吹奏这个旋律，每个长笛手相对于其他吹奏同一张曲谱的人又有所延迟。也许这个极度夸张的形象能使你对你身体中每分每秒都在进行着的过程的复杂性的某些方面有所想象。

\pandocbounded{\includegraphics[keepaspectratio]{./media/OEBPS/Images/img_098.png}}

图98．二排的分子卡农

这是一个更为复杂的图示，不只一个而是有若干个mRNA串，它们全都转录自同一个DNA串，并受到多核糖体作用，其结果是一个二排的分子卡农。{[}选自哈那瓦尔特和海因斯，《生命的化学基础》第271页{]}

谁生谁------核糖体与蛋白质

我们一直在谈论那些称为核糖体的惊人活物，可它们自己又是怎么构成的、如何制造的呢？核糖体由两种东西构成：（1）各类蛋白质，（2）另一种RNA，叫做核糖体RNA（rRNA）。于是，为了造出核糖体，得有某些种类的蛋白质，还得有rRNA。自然，要有蛋白质就得有核糖体来制造它。那么，你怎么应付这个恶性循环呢？先有谁------核糖体还是蛋白质？谁制造谁？当然不会有答案，因为往前追溯总是会遇到当前对象的同类------就像先有鸡还是先有蛋的问题一样------直到一切都消失在时间的地平线上为止。无论何时，核糖体总是由两块东西制成，一大一小，每一块中都含有一个rRNA和几个蛋白质。核糖体差不多和大的蛋白质一样大小。它们远远小于作为它们的输入并且它们要顺着移动的mRNA串。

蛋白质的功能

我们已经谈过一些蛋白质的结构，尤其是酶的结构，但我们还没有实际讲述它们在细胞内执行什么任务，以及如何去做。全部的酶都是催化剂，这意味着，在某种意义上，它不过是有选择地加速细胞中的一些化学过程，而不是使那些没有它们就不能发生的事情发生。每种酶都使千万条潜在通道中的一条特定通道得以实现。因此，在选定要引入哪一种酶的同时，你也就选定了将要发生的事情和将不发生的事情------尽管从理论上说，任何细胞过程不靠催化剂帮助而自动发生的概率都不是零。

酶怎样在细胞里的分子上起作用呢？如已经讲过的那样，酶是折起的多肽链，每个酶中都有一个裂缝或口袋或某种别的清晰定义的外观特性，酶在那里拴住某个其它种类的分子。这个地点叫活性部位，拴在那里的分子叫基质。酶可以有不止一个活性部位和不只一个基质。像印符遗传学一样，酶对于它们所要发挥作用的部位实在是很挑剔的。活性部位通常都很``挑食''：只许一种分子拴在上面。尽管如此，有时也会遇上``李鬼''------一些其它分子，它们也能适应这个作用点并赘在上面，哄骗这个酶，从而事实上使酶起不了作用。

一旦把酶与其基质拴在一起，就会发生某种电荷不平衡状态，随后，电荷------以电子或质子的形式------就围绕着拴住的分子流动并重新排布自己。达到平衡时，基质可能已经发生了意义深刻的化学变化。例如：可能已经发生了一个``焊接''过程，把一些标准的小分子连成核苷酸、氨基酸或其它的普通细胞分子；一个DNA串可能已在一个特殊的部位``豁开''；一个分子的某一部分可能已被割掉；等等，等等。事实上，生物酶对分子所起的作用十分类似于印符酶所执行的印符操作。只不过大多数酶基本上都只执行单项任务，而不是执行一系列任务。印符酶和生物酶之间还有一个显著的区别，那就是：印符酶只作用在串上，而生物酶则可以作用在DNA、RNA、其它蛋白质、核糖体以及细胞膜之上------简言之，作用于细胞里的每一样东西。换句话说，酶是在细胞内使各种事件得以发生的通用机器。有各种各样的酶，有的能把东西粘在一起，有的又把它们分开，有的修改它们，有的激活它们，有的安定它们，有的复制它们，有的恢复它们，有的毁坏它们\ldots\ldots{}

细胞内某些极为复杂的过程是``多级瀑布''形的，在那里面，某种类型的单个分子导致某种酶的生产过程；生产过程开始后，从``组装线''中制出的酶又开辟一条新的化学途经使第二种酶开始生产。这种事情可以持续三四层，每一种新生产出的酶又引起另一类型的生产过程。最后，生产出一大批最终型号的酶的副本，然后所有的副本就脱离出来，去执行它们各自的任务，也许是去除掉某个``异端''DNA，或协助制作细胞所``渴望''的某种氨基酸，或是别的什么事情。

需要有一个足够强有力的支撑系统

我们来说说印符遗传学所揭示的大自然对``什么样的DNA串能指挥自己的复制？''这个谜题的解答。当然，并非每个DNA串生来都是自复制。关键在于：任何一个想要指挥自己的复制工作的串，都必须含有一些命令，以便能恰好调集那些能够执行这一任务的酶。指望一个DNA串孤家寡人就能成为自复制，这是枉费心机。因为，要把那些潜在的蛋白质从DNA中拽出来，不仅得有核糖体，而且还得有RNA催化酶。后者要用来制造输送给核糖体的mRNA。所以我们必须从假定有一种刚好强到允许转录过程和翻译过程得以进行的``极小支撑系统''开始。于是，这个极小支撑系统就得包括：（1）某些保证mRNA能从DNA制造出来的蛋白质------例如RNA聚合酶，（2）某些核糖体。

DNA如何自复制

``足够强有力的支撑系统''和``足够强有力的形式系统''两种说法听起来很有些相似，这无论如何不能算是巧合。一个是产生自复制的前提，一个是产生自指的前提。其实，从本质上讲，这只是一种现象在两个不同的外观下出现，而我们马上就要把它们明显地对应起来。不过，做这件事情以前，我们先得说清楚一个DNA串怎么能成为自复制。

DNA必定含有一组蛋白质的密码，而这组蛋白质就将复制这个DNA。有一种十分有效而精致的办法能复制一个由互补的两个串组成的双串DNA。这包括两步：

⑴．把两个串彼此分开，

⑵．给刚刚得到的两个新的单串各``配''上一个新串。

这个过程能造出两个新的DNA串，每一个都与原先的双串等同。如果我们的解决办法是建立在这个想法上的，那就必定要涉及一组在DNA中有自己的密码、并将要执行这两个步骤的蛋白质。

据信，在细胞内这两个步骤是通过一种协调的方式一起进行的，而且需要三种主要的酶：DNA核酸内切酶，DNA催化酶和DNA连接酶。第一种是个``拉开酶''：它把原有的两个串剥开一小段距离，然后停止。这之后，另外两种酶登场。催化酶基本上是一种完成``复制并移动''操作的酶：它切下DNA的一个短短的单串，以一种会使人联想起印符遗传学中的复制状态的方式用补基来复制它们。为了复制，它要利用在细胞质中四处漂浮的一些原料------尤其是核苷酸。由于这个过程是阵发性的，每次都剥开一些东西，复制一些东西，所以就产生了一些小缝隙，而DNA连接酶就是堵缝的东西。这个过程一遍又一遍地重复。这部精密的三酶机就以一种小心翼翼的方式一路沿着DNA分子运行着，直到整个东西都被剥开，同时也被复制，最后有了两个副本为止。

DNA的自复制方法与{㧟}摁相比较

应该注意，在酶对DNA串的作用过程中，全然没有用到信息储存在DNA之中这一事实。酶只是在执行调拨符号的职能，就像WJU系统中的推理规则一样。这三种酶在某一时刻就在实际复制作为他们自身的密码的那些基因，但它们对此却无动于衷。对它们来说，DNA只是一块既无意义也无趣味的模板。

把这与蒯恩句子那种描述如何构造自己的副本的方法作个比较是颇为有趣的。在那里，也有一种``双串''------同一信息的两个副本，其中一个起指令作用，另一个起模板作用。在DNA中，过程有一点类似，因为三种酶（DNA核酸内切酶、DNA催化酶、DNA连接酶）只在其中的一个串中被编码，这个串就起着程序的作用，而另一个串只是块模板。这个平行还不完美，因为在执行一个复制动作时，两个串都作模板使用，而不只是一个串。即便如此，这个类比也是具有高度启发性的。对这种使用与谈论的二重性，有一个生物化学的类比：当DNA被看作只是一系列有待复制的化学物质时，就相当于谈论印刷符号；而当DNA指示该执行哪些操作时，就相当于使用印刷符号。

DNA的意义层次

从一个DNA串里可以读出好几个不同层次的意义，这要看你把DNA划成多大的块，以及你使用的解码器有多强。在最低的层次上，每一个DNA串给一个等效的RNA串编码------作为转录的编码过程。如果将DNA组块而形成三元组，那么，利用``遗传解码器''就能把DNA读成一系列氨基酸。这是（在转录之上的）一个翻译。在下一个自然的层次上，DNA又可读成一组蛋白质的代码。这个把蛋白质从基因里拽出来的物理过程叫做基因表示。目前，这是我们所了解的DNA意义的最高层次了。

然而，肯定有更高层次的意义，它们更难觉察。例如，有种种理由使人们相信（比如说人的）DNA含有个人特征（如鼻子的形状、音乐天份、反应速度等等）方面的密码。我们能不能在原则上学会直接从一个DNA串读出这些信息，而不必经历实际的物理渐成过程（从遗传型到表现型的物理显示）呢？大概是可以的。因为在理论上我们可以拥有一个强得难以置信的计算机程序来模拟整个的过程，包括每个细胞、细胞中的每种蛋白质、以及DNA复制时所涉及的每个细微末节，等等，一直做到底。这样一个伪渐成程序就是表现型的一个高层次描述。

另外还有一种（极其微弱的）可能：我们毋需去同构模拟渐成过程，而是通过找到某种较为简单的解码器，直接从遗传型读出表现型来。这可以叫做``捷径伪渐成过程''。捷径也好，非捷径也好，伪渐成过程目前当然还远未达到------尽管有个惹人注目的例外：对于猫科的一个种------家猫{[}Felis-catus{]}，深入的探索已经揭示出，确有可能直接从遗传型读出表现型来。直接观察一下家猫的DNA中下述典型节段，可能就会更好地体会这个值得注意的事实：\footnote{}

\ldots\ldots CATCATCATCATCATCATCATCATCATCAT\ldots\ldots{}

下面列出了各个层次上DNA读出方式的一览表，以及不同的译码层次的名称。DNA能被读成这样的序列：

\begin{longtable}[]{@{}lllll@{}}
\toprule\noalign{}
\endhead
\bottomrule\noalign{}
\endlastfoot
\multicolumn{3}{@{}l}{%
⑴基（核苷酸）} & \ldots{} & 转录 \\
\multicolumn{3}{@{}l}{%
⑵氨基酸} & \ldots{} & 翻译 \\
⑶蛋白质 & （一级结构） & \multirow{3}{=}{\pandocbounded{\includegraphics[keepaspectratio]{./media/OEBPS/Images/Formula-right_bracket.png}}} & ~ & ~ \\
⑷蛋白质 & （二级结构） & & \ldots{} & 基因表示 \\
⑸蛋白质 & （三级结构） & & ~ & ~ \\
\multicolumn{3}{@{}l}{%
⑹蛋白质束} & \ldots{} & 高层基因表示 \\
· & · & \multirow{3}{=}{\pandocbounded{\includegraphics[keepaspectratio]{./media/OEBPS/Images/Formula-right_bracket.png}}} & ~ & ~ \\
· & · & & \ldots{} & \multirow{2}{=}{DNA意义的未知层次} \\
· & · & & ~ \\
(N-1)　 & ？ & ~ & \ldots{} & ？ \\
\multicolumn{3}{@{}l}{%
(N)物理、精神和心理的品质} & \ldots{} & 伪渐成过程 \\
\end{longtable}

中心法则映射

有了这个背景，我们就可以在克里克的``分子生物学中心法则''（称做法则Ⅰ，这是一切细胞过程的基础）和我富有诗意的发明------``数理逻辑中心法则''（称作法则Ⅱ，是哥德尔定理的基础）------之间做一个精细的比较。图99和下面的表显示了其间的映射，这个图和表一起构成了``中心法则映射''。

\begin{longtable}[]{@{}lll@{}}
\toprule\noalign{}
\endhead
\bottomrule\noalign{}
\endlastfoot
法则Ⅰ（分子生物学） & & 法则Ⅱ（数理逻辑） \\
DNA串 & \gebfont{⇔} & TNT符号串 \\
mRNA串 & \gebfont{⇔} & N的陈述 \\
蛋白质 & \gebfont{⇔} & 元TNT的陈述 \\
作用在蛋白质上的蛋白质 & \gebfont{⇔} & 关于元TNT的陈述的陈述 \\
作用在作用在蛋白质上的蛋白质上的蛋白质 & \gebfont{⇔} & 关于关于元TNT陈述的陈述的陈述 \\
转录（DNA\gebfont{⇒}RNA） & \gebfont{⇔} & 解释（TNT\gebfont{⇒}N） \\
翻译（RNA\gebfont{⇒}蛋白质） & \gebfont{⇔} & 算术化（N\gebfont{⇒}元TNT） \\
克里克 & \gebfont{⇔} & 哥德尔 \\
遗传密码（任意约定） & \gebfont{⇔} & 哥德尔编码（任意约定） \\
密码子（基因的三元组） & \gebfont{⇔} & 密码子（数码的三元组） \\
氨基酸 & \gebfont{⇔} & 元TNT中引用的TNT符号 \\
自复制 & \gebfont{⇔} & 自指 \\
使自复制得以发生的充分强的细胞系统 & \gebfont{⇔} & 使自指得以发生的充分强的算术形式系统 \\
\end{longtable}

\pandocbounded{\includegraphics[keepaspectratio]{./media/OEBPS/Images/img_099.png}}

图99．中心法则映射

在两个基本的缠结的层次系统------分子生物学层次系统和数理逻辑层次系统------之间的一个类比

请注意A和T组成的基对（算术化{[}Arithmetization{]}和翻译{[}Translation{]}），以及G和C的基对（哥德尔{[}Gödel{]}和克里克{[}Crick{]}）。数理逻辑居嘌呤一方，分子生物学居嘧啶一方。

为了使这个映射从美学角度看更为完善，我情愿让哥德尔配数法向遗传密码绝对看齐。事实上，只要按下面的对应，遗传密码表也就变成了哥德尔编码表：

\begin{longtable}[]{@{}lll@{}}
\toprule\noalign{}
\endhead
\bottomrule\noalign{}
\endlastfoot
（奇）1 & \gebfont{⇔} & A（嘌呤） \\
（偶）2 & \gebfont{⇔} & C（嘧啶） \\
（奇）3 & \gebfont{⇔} & G（嘌呤） \\
（偶）6 & \gebfont{⇔} & U（嘧啶） \\
\end{longtable}

\pandocbounded{\includegraphics[keepaspectratio]{./media/OEBPS/Images/img_100.png}}

图100．哥德尔编码

按照这个哥德尔配数方案，每个TNT符号有一个或几个密码子。小椭圆圈表明这个表怎样包含着第九章中的哥德尔配数表

每种氨基酸（共二十种）恰好对应一个TNT符号（共二十个）。这样，我当初编造``简朴TNT''的动机也就终于明朗了------要它恰好有二十个符号！哥德尔编码如图100所示。试将它与遗传密码（图94）比较一下。

看到二十世纪所获得的知识中，两个极其深奥而又十分重要的进展如此深刻地分享了这样一种抽象的结构，简直会让人感到有什么近乎神秘的东西存在。中心法则映射无论如何不是这两种理论等价的一个严格证明，不过它清楚地显示了一个意义深远的亲缘关系，这是值得深究的。

中心法则映射中的怪圈

这个映射的两边有一个更为有趣的相似，即那种具有任意复杂度的``圈子''在两边渐渐达到顶层的方式：左边，是作用在作用在蛋白质上的蛋白质上的蛋白质，等等，直至无穷；而右边，是关于关于元TNT陈述的陈述的陈述，等等，直至无穷。这些都类似于我们在第五章中讨论过的异层结构，在那里，一个足够复杂的基质就可以使高层怪圈出现，并不断地兜圈子，以至完全隔离于较低的层次。我们将在第二十章中详细探究这一思想。

说到这里，你可能奇怪：``按照中心法则映射，\heititext{哥德尔不完全性定理}该对应于什么？''这倒是个应该在读下去之前好好想想的问题。

中心法则映射与《对位藏头诗》

看得出，中心法则映射十分类似于第九章中建立的《对位藏头诗》与哥德尔定理之间的映射。因此，我们可以找出全部三个系统之间的平行性：

⑴．形式系统和符号串，

⑵．细胞和DNA串，

⑶．唱机和唱片。

下面的图表详细解释了系统2和系统3之间的映射：

\begin{longtable}[]{@{}lll@{}}
\toprule\noalign{}
\endhead
\bottomrule\noalign{}
\endlastfoot
《对位藏头诗》 & & 分子生物学唱机 \\
唱机 & \gebfont{⇔} & 细胞 \\
``完备的''唱机 & \gebfont{⇔} & ``完备的''细胞 \\
唱片 & \gebfont{⇔} & DNA串 \\
可以在一个给定的唱机上播放的唱片 & \gebfont{⇔} & 可以由一个给定的细胞复制的DNA串 \\
不可在这台唱机上播放的唱片 & \gebfont{⇔} & 不可由这个细胞复制的DNA串 \\
把唱片的音槽转变为声音的过程 & \gebfont{⇔} & 把DNA转录到mRNA上的过程 \\
唱机产生的声音 & \gebfont{⇔} & 信使RNA串 \\
声音翻译为唱机的震颤 & \gebfont{⇔} & mRNA翻译成蛋白质 \\
从外部音响到唱机震颤的映射 & \gebfont{⇔} & 遗传密码（mRNA三元组到蛋白质的氨基酸的映射） \\
唱机的毁坏 & \gebfont{⇔} & 细胞的毁灭 \\
为唱机X特制的歌，标题是``我不能在唱机X上播放'' & \gebfont{⇔} & 为细胞X特制的DNA串高层解释：``我不能由细胞X复制'' \\
``不完备''的唱机 & \gebfont{⇔} & 至少有一个DNA串复制不出来的细胞 \\
``疙瘩定理''：``给定一台具体的唱机，总有一张它不能播放的唱片。'' & \gebfont{⇔} & 免疫定理：``给定一个具体的细胞，总有一个它复制不了的DNA串。'' \\
\end{longtable}

哥德尔定理的对应物看上去是一个奇怪的东西，也许对分子生物学家来说没什么用处（在他们看来，这似乎十分明显）：

总可以设计一个DNA串，如果把它注入细胞中，在被转录时它将促使一些将会毁掉这个细胞（或DNA）的蛋白质生产出来，而这就造成该DNA的``绝收''。

这会促成一个滑稽有趣的构想，至少从进化论的角度看是如此：一种入侵型的病毒偷偷摸摸进入了一个细胞，然后精心制作一些具有破坏这个病毒本身的效力的蛋白质！这是一种细胞水平上的自杀------如果你愿意，也可说是细胞水平上的说谎者句子。显然，这并没有从物种生存的角度证明有什么好处。不过，如果我们不拘泥于字眼，可以认为这说明了细胞与其入侵者各自发展起来的两种机制------保护和颠覆------的精神实质。

大肠杆菌与T4之战

\pandocbounded{\includegraphics[keepaspectratio]{./media/OEBPS/Images/img_101.png}}

图101．T4细菌病毒

它是一个蛋白质聚合体(a)。它的``头''是形如一种具有三十面的扁长的正十二面体的蛋白质膜，里面充满了DNA，头通过一个颈与尾部相连。尾部是包在一个可收缩的鞘内的空心芯子，安在四周有六条须根、下面有若干钉子的底盘上。钉子和须根使病毒牢牢地粘在细菌的细胞壁上(b)。鞘的收缩使芯子穿透细胞壁，病毒的DNA进入细胞。{[}选自哈那瓦尔特和海因斯，《生命的化学基础》第230页{]}

我们来考虑生物学家特别宠爱的一种细胞：大肠杆菌细胞，以及生物学家特别宠爱的入侵这种细胞的家伙：凶恶可怕的T4噬菌体，其形状如图101。这种神奇的小东西看上去有点像登月舱和蚊子的混血儿，可它却比蚊子凶得多。它有一个``头''，里面储存着它的全部``知识''------即它的DNA；它有六条``腿''，以便抓牢它选定入侵的细胞；它有一个像蚊子那样的``刺针''（叫``剌尾''可能更合适些）。主要的区别是，蚊子利用它的刺针来吸血，而T4噬菌体则利用它的剌针违背它的牺牲品的意志，把自己的遗传物质强行注入细胞。所以噬菌体在那个小小的天地里犯有``强奸''罪。

分子特洛伊木马

病毒DNA进入细胞后会发生什么事呢？按拟人的说法，病毒``希望''它的DNA能得到与宿主细胞的DNA完全相同的待遇。这就意味着要能被转录和被翻译，于是就使它可以指挥合成它自己的特殊的、与宿主细胞相异的蛋白质，然后这些蛋白质就开始各司其职。这相当于利用``密码''（即遗传密码）秘密地把外来蛋白质输送到宿主细胞中，然后再来``解码''（制造它们）。在某种意义上，这很像特洛伊木马的故事：数百名战士藏进一个看上去无害的巨大木马，当他们进城之后，就钻出来占领城市。这些异己的蛋白质一旦从它们的载体DNA中被``翻译''出来（合成出来），就投入行动。

\pandocbounded{\includegraphics[keepaspectratio]{./media/OEBPS/Images/img_102.png}}

图102．细菌受到病毒的感染

当病毒的DNA进入细菌之后，病毒感染就开始了。细菌的DNA被瓦解，而病毒的DNA得到复制。病毒结构蛋白的合成以及它们装配成病毒的过程一直在继续，直到撑破细胞，释放出微粒为止。{[}选自哈那瓦尔特和海因斯，《生命的化学基础》第230页{]}

由噬菌体T4所指挥的一系列行为已得到仔细研究，差不多是如下所述（也见图102和图103）：

\begin{longtable}[]{@{}ll@{}}
\toprule\noalign{}
\endhead
\bottomrule\noalign{}
\endlastfoot
发生时刻 & 发生的行为 \\
第0分 & 注入病毒DNA。 \\
第1分 & 破坏了主体的DNA。细胞本身的蛋白质停止生产，异己的（T4）蛋白质开始生产。最先生产出的一批蛋白质是负责指挥复制异己（T4）DNA的。 \\
第5分 & 开始复制病毒DNA \\
第8分 & 开始生产将要组成新噬菌体``躯体''的结构蛋白。 \\
第13分 & T4入侵物的第一个完整的复制品被生产出来。 \\
第25分 & 溶菌酶（一种蛋白质）攻击寄主的细胞壁，打破细菌，``二百元组''出现。 \\
\end{longtable}

这样，当噬菌体T4侵入大肠杆菌细胞时，经过大约二十四五分钟这样一段不长的时间，这个细胞就完全被颠覆、撕裂。啪的一下，大约二百个与原先的病毒完全相同的复制品------``二百元组''------跑了出来，准备去攻击更多的细菌细胞，而原先的细胞在这个过程中被大幅度地毁坏。

\pandocbounded{\includegraphics[keepaspectratio]{./media/OEBPS/Images/img_103.png}}

图103．T4病毒的形态

形态发生通道有三个主分支，分别形成头、尾和尾须，然后结合起来形成完整的病毒微粒。{[}选自哈那瓦尔特和海因斯，《生命的化学基础》第237页{]}

尽管从一个细菌的角度来看，这种事情是一种极为严重的威胁，但从我们的大尺度的优劣观点来看，可以把它看成两名对手之间的一种娱乐性博弈：入侵者，或T方（依照病毒的T偶数族命名，包括T2，T4，等等），和``C''方（代表``细胞''{[}Cell{]}）。T方的目标是：为了复制自己而侵入C的细胞，并从内部接管C的细胞。C方的目标则是：保护自己并消灭入侵者。用这种方式描述的话，分子TC-博弈可以被认为是十分类似于前面对话中描述的宏观龟{[}Tortoise{]}蟹{[}Crab{]}之战。

识别、伪装、标识

这种``博弈''强调了一个事实：识别过程是细胞和亚细胞生物学的中心主题。分子（或一些高层结构）如何互相识别？为了做到这一点，酶必须能锁在其基质的特殊``拴缚部位''上，细菌必须能区别它自己的DNA和噬菌体的DNA，两个细胞必须能通过一种有控制的方式彼此识别和相互作用。这种识别问题可能会使你想起当初关于形式系统的那个关键问题：你怎么才能说清一个符号串有没有某种性质（比方说``是定理''）？有没有一个判定过程？这类问题并不限于数理逻辑：它还渗入计算机科学乃至（我们已经看到了）分子生物学。

对话中描述的那种标识技术，实际上就是大肠杆菌智斗噬菌体入侵者的计策。其思想是：DNA串能通过一种化学手段来标识，即把一个小分子------甲基------排在各种核苷酸上。这种标识操作并不改变DNA的正常生物性质，换句话说，甲基化（标识了的）的DNA能够转录得和未甲基化（未标识的）的DNA一样，所以它仍然能指挥蛋白质的合成。但是，如果主体细胞有某些特殊机制能检验DNA是否有标识，那么这种标识就会是非常重要的了。具体地说，主体细胞可能有一个酶系统来搜寻未加标识的DNA，而且只要找到这样的DNA就毫不留情地把它砍碎破坏掉。此刻，呜呼哀哉的是所有未加标识的入侵者。

核苷酸上甲基标识可以比作一个书法流派的特征。利用这种比喻，我们可以说大肠杆菌细胞在寻找以其``本派字体''------即与它自己特征一致的书法------写下的DNA，并砍掉以``异端''书法写下的一切DNA。当然，噬菌体的反策略则是学会标识自己，从而哄骗那些它们打算侵入的细胞来复制自己。

这种TC之战一类的博奕可以继续到任意复杂的层次，但我们不再继续了。基本的事实是：这是一场力图排斥全部入侵者的主体细胞与力图把自己的DNA渗入宿主细胞、并使它把这个DNA转录成mRNA（从而保证它的复制）的噬菌体之间的战斗。任何一种成功地获得这种自复制方法的病毒DNA，都可以说是具有这样一种高层次的解释：``我可以在X型细胞内被复制。''一定要把这与早先谈到的、在进化论上没有什么意义的那种噬菌体区别开来。早先谈的那种噬菌体是给破坏自己的那些蛋白质编码，其高层解释则是一个自戕句子：``我不能在X型细胞内被复制。''

汉肯句子与病毒

分子生物学中这两种互为对照的自指类型，在数理逻辑中都有对应的东西。我们讨论过自戕噬菌体的类比物------哥德尔型的符号串，它断言自己在特定的形式系统中不可制造。此外，我们还可做出一个与实际噬菌体相应的句子：该噬菌体断言了它自己在特定的细胞内可以生成。这样的句子就是断言它自己在特定形式系统中可以生成。这类句子被冠以数理逻辑学家列昂·汉肯的名字，称作汉肯句子。可以严格顺着哥德尔句子的线路去构造它们，唯一的不同是丢掉一个否定号。当然，还得从一个``服''号串开始：

\gebfont{∃}a:\gebfont{∃}a\textquotesingle:＜TNT-PROOF-PAIR\{a, a\textquotesingle\}∧ARITHMOQUINE\{a", a\textquotesingle\}＞

然后利用标准的手段继续下去，令上述``服''号串的哥德尔数是h，于是就对这个``服''号串作算术\gebfont{㧟}摁，从而得到一个汉肯句子：

\begin{center}

\pandocbounded{\includegraphics[keepaspectratio]{./media/OEBPS/Images/Formula14.png}}

\end{center}

（顺便说一句，你能看出这个句子与\textasciitilde G的不同点吗？）我明白地把它写出来，是为了指出：汉肯句子并没有给它自己的推导开出完整的处方，它只是断言存在一个推导。你可能对它的断定能不能得到验证深感怀疑。汉肯句子真有推导吗？它们的确像它们自己宣称的那样是个定理吗？想一下下面的情形是会有所帮助的：人们怎么会相信一个自称``一贯正确''的``死亡研究专家''呢------他可能正确，但也可能不正确。汉肯句子就比那种专家更靠得住吗？如果那种专家对你说：``其实，人都会死于饭前洗手''，你相信吗？

已经证明，这些汉肯句子一律都讲的是真话。理由并不显然，不过我们将不加证明地接受这一使人奇怪的事实。

隐式汉肯句子之别于显式汉肯句子

刚才说了，一个汉肯句子对于它自己的推导过程并没说出什么，只是断言有那么一个推导存在。现在，就有可能发明出汉肯句子这一主题的一个变奏------明确描述了自己的推导过程的句子。这种句子的高层解释就不是``存在某个符号串序列是我的推导''，而是``此处描述的符号串序列\ldots\ldots 是我的推导。''我们把前一种句型叫隐性汉肯句子，而把这种新型的句子称作显性汉肯句子，因为它们显式地描述了自己的推导。应当注意的是，显性汉肯句子与其隐性的弟兄们不同，它们不一定是定理。事实上，很容易写出一个符号串，它断言自己的推导由单个符号串0=0组成，但这是个假句子，因为0=0不是任何东西的推导。不过，也可以写出一个是定理的显性汉肯句子------即事实上给出它自己的推导方案的句子。

汉肯句子和自组装

我之所以要区分显性和隐性汉肯句子，是因为这很好地对应着两类病毒之间的一个重要的区别。有一类病毒（比如所谓的``烟草花叶病毒''），叫做自组装病毒；同时也还有另一类病毒，比如我们最喜欢用的T偶数族病毒，叫做非自组装病毒。区别何在？这可以直接类比于隐性和显性的汉肯句子。

一个自组装病毒的DNA只给新病毒的各个部分编了码，而不给任何酶编码。一旦这些部分都造了出来，那个诡秘的病毒就依靠它们自行连接起来，而毋需什么酶的帮助。这样的过程依赖于各个部分在细胞的醇厚饮料中游泳时彼此之间所具有的化学亲和力。不光是病毒，某些细胞器------例如核糖体------也是自组装的。有些时候可能也需要酶------不过在这种情形，酶是从宿主细胞中硬拉来为其所奴役的，这就是自组装的含义所在。

反之，诸如T偶数号那样的比较复杂的病毒的DNA，就不仅给各个部分编码，而且也给在各部分组装成整体的过程中起特殊作用的各种酶编码。由于这个组装过程不是自发的而是需要``一些装置''，所以这些病毒不能看作是自组装。自组装元件与非自组装元件之间的区别，从本质上讲就是：前者不必告诉细胞任何有关它们结构的事情，就安然完成自复制，而后者则需要给出一些有关如何组装它们自己的指令。

这样一来，与隐性汉肯句子和显性汉肯句子相平行的东西也就非常清楚了。隐性汉肯句子是自证明的，但关于它们的证明却什么也没说出来------它们类似于自组装病毒；显性汉肯句子指示了自己证明的构造------这类似于在复制自己的过程中指挥主体细胞的那些比较复杂的病毒。

像病毒这样复杂的自组装生物结构的概念，会使我们想到建造复杂的自组装机器的可能性。设想有这样一系列的零部件，当把它们放入一个适当的背景环境中时，就能通过某种方式自动地集合起来形成一台复杂的机器。这似乎不着边际，实际上却是描述烟草花叶病毒依靠自组装来进行自复制的过程的一种精确方式。有机体（或机器）的整个构造的信息，散布在它的各个部分之中，而不是集中在某一个地方。

这个概念还能引导我们走向一个奇怪的方向，如在《一位烟民富于启发性的思想》中所看到的那样。在那里，我们看到螃蟹如何利用这样一种思想：自组装信息可以四下散开，而不集中在一处。他希望这能防止他的唱机毁于乌龟的破坏手段。不幸的是，正像那些最最精致的公理模式一样，一旦这个系统建立起来并``梳理''完毕，它的良定义性就使它要受到一个足够聪明的``哥德尔化算子''的攻击，而这就是螃蟹所讲的那些令人伤心的故事。尽管看起来荒唐，对话中异想天开情节的与现实的距离并不那么遥远，它就存在于陌生的、梦幻般的细胞世界之中。

两个突出的问题：分化与形态发生

这样，自组装可以是构造细胞的某些子单元以及某些病毒的手段------可是构造那些极为复杂的宏观结构------比如大象或蜘蛛的躯体，或者维氏捕蝇草------的手段是什么呢？归巢的本能是怎样在鸟的脑子中建立起来的？追猎的本能又是如何在狗的脑子里建立的？简言之，DNA如果仅仅能指挥在细胞中合成一些蛋白质，它又怎么能如此惊人准确地控制宏观生物的严格结构和功能呢？这里主要有两个不同的问题。一个是细胞分化的问题：如何区分具有同样的DNA却又扮演不同角色的细胞------例如肾细胞、骨髓细胞、脑细胞？另一个是形态发生的问题（``形体的出生''）：局部水平上的细胞间的相互作用，如何导致大规模的、总体的结构和组织------诸如身体的各个器官、面貌、脑的各个部位等等？尽管我们现在对细胞的分化和形态发生都还知之甚少，但看起来这种手段是存在于细胞内部和细胞之间非常协调的反馈和``前馈''机制之中的，这种机制通知细胞何时``启动''何时``停止''各种蛋白质的生产过程。

反馈和前馈

在细胞中，如果所需要的某种物质太多了或太少了，就会发生反馈，那时细胞就得用某种方式调整组装这种物质的生产线。前馈也涉及对组装线的调整，不过不是根据当时最终产物的总量，而是根据该组装线的某些``已出厂的产品''的总量。有两种主要方式实现否定性的前馈或反馈。一种方法是阻止有关的酶起作用------也就是说``堵住''它们的活性部位。这叫做抑制。另一种方法是干脆不让有关的酶产生！这叫做阻遏。从概念上讲，抑制比较简单：只须堵住组装线上第一个酶的活性部位，整个合成过程就要停下来。

阻遏物和诱导物

阻遏则比较诡秘。一个细胞如何阻止一个基因，不让它被表示出来呢？回答是：细胞使这个基因转录不出来。这意味着它必须阻止RNA聚合酶开展工作。顺着DNA的长线，恰好在细胞不想转录的那个基因面前的通道中设置一个巨大的障碍物，就可以达到这个目的。这种障碍物是有的，叫做阻遏物。阻遏物本身也是蛋白质，而且它们能拴在DNA上特定的放置障碍物的部位，这种部位（我弄不清为什么）叫做操纵基因。于是，一个操纵基因就是对紧挨在它后面的基因进行控制的部位，它后面的基因则称为操纵子。由于在实现一个长的化学变化的过程中，是一系列的酶在一齐发生作用，所以它们经常是一个挨一个地编码的。这就说明了为什么一个操纵子常常含有几个基因而不只是一个。成功地阻遏一个操纵子，产生的效果是这一系列基因都得不到转录，这意味着会有一大批酶都处于未合成的状态。

肯定性的前馈和反馈又怎么样呢？仍有两种可能：（1）释放受阻的酶，（2）停止对有关操纵子的阻遏。（应当注意到大自然多么喜欢双重否定！这可能有着十分深刻的道理。）使阻遏受到阻遏的机制要用到一类分子，称为诱导物。诱导物的作用很简单：它在一个阻遏蛋白有机会拴在DNA分子的某个操纵基因上之前，与这个阻遏蛋白结合起来。这样得到的``阻遏物与诱导物复合体''就不可能拴在一个操纵基因上，于是就为把联着的操纵子转录到mRNA上，然后再翻译成蛋白质打开了方便之门。通常，最终的产品或某些``已出厂产品''可以起诱导物的作用。

反馈与怪圈的对比

话说到这里，让我们来区别一下抑制和阻遏这类过程中的简单型反馈，和中心法则映射中所展示的不同信息层次之间的怪圈。在某种意义上讲，两者都是``反馈''，但后者要比前者深刻得多。当诸如色氨酸或异亮氨酸这样的一个氨基酸（以诱导物的形式）拴住它的阻遏物起着反馈作用，从而使其更多的副本得以制造的时候，它并没有说出如何构造它自己。它只是通知酶去制造那些副本。这可以与收音机的音量作个比较。当声音注入听者的耳朵时，有可能导致人去减小或增大音量本身。但这与广播本身命令你打开或关掉收音机------或者要你调到另一个波长上，或者，甚至是教你如何组装另一台收音机------相比，是完全不同的两码事！后面的这些事更像是各种不同的信息层次之间的那种兜圈子。因为此时收音机信号中的信息被``解码''并翻译成心智上的结构。收音机的信号是一些符号构成物，是这些符号的意义在起作用------这时是使用，而非谈论。另一方面，当声音太响时，这些符号就不是在传达意义。它们只被看作很响的声音，因而等于没有意义------这时是谈论，而非使用。这种情形就更像蛋白质调整其合成速率所使用的反馈圈子了。

两个相邻的、具有完全相同的基因类型、但有不同功能的细胞之间的区别，在于它们的基因组中的不同节段受到阻遏，因而它们的蛋白质工作点不同，这已经形成理论了。这样的假说可以解释人体的不同器官中细胞何以会有显著的差别。

分化的两个简单例子

从初始的细胞开始，一遍又一遍地复制，造出大量各种各样的、具有特定功能的细胞的过程，可以比做一种在人们之间传送连锁信的过程。其中每个新的参加者都得忠实地传达消息，但都要加一点润色。最后就会得到很多彼此相差极大的书信。

对这种分化思想给出另一种说明的，是下面这种与分化型自复制相似的极为简单的计算机。设想一个很短的、被一个双向开关控制的程序，程序中有一个内部参数N------一个自然数。这个程序能按两种方式运行------向上方式和向下方式。当按照向上方式运行时，它把自己复制到该机器的存储器的邻近部位，不过它的``女儿''的内部参数是原先的N加上1。当按照向下方式运行时，它不做自复制，而是计算数值

(-1)\textsuperscript{N}/(2N+1)

然后把它加到前面运行过程中所得的总数上。

好了，假定一开始存储器中有一份这个程序，其中的N=0，状态为``向上''。那么该程序就要在存储器里相邻部位中复制自己，并使N=1。重复这个过程，新程序又要在与自己相邻处复制自己，得到N=2的副本。就这样一遍一遍地重复下去\ldots\ldots 所发生的事是一个很大的程序在存储器中膨胀。存储器装满后，过程就停止了。可以看成是整个存储器中装满了一个大程序，它含有很多彼此相似但却是分化了的模块------或叫``细胞''。如果此刻把程式开关扳到``向下''，并运行这个大程序，那会发生什么呢？首先第一个``细胞''运行起来，算出1/1。接着第二个``细胞''运行，算出-1/3，并和前面的结果相加，然后第三个``细胞''运行，算出1/5，并与前面的和相加\ldots\ldots 最终的结果是整个``组织''------那个大程序------算出大量（与存储器内能装下的``细胞''一样多）的项的和：

1-1/3+1/5-1/7+1/9-1/11+1/13-1/15+\ldots\ldots{}

由于这个级数收敛（尽管很慢）于π/4，所以我们就有了一个``表现型''的东西，它的功能是计算一个著名的数学常数的值。

细胞中的层次混合

我希望，对诸如标识、自组装、分化、形态发生、以及转录和翻译等过程的描述，能有助于描述一个十分复杂的系统，那就是细胞------一种面貌全新的信息加工系统。在中心法则映射中，我们已经看到，尽管可以试着在程序和数据之间划一条分界线，但这种区别多少有些随意。按这种思路再进一步，我们就会发现，不仅程序和数据犬牙交错，就是程序的解释程序以及实际的处理机、乃至语言，都是如此紧密相关的。尽管（在一定程度上）有可能给各层次划定边界，把它们分开，但是认识到各层次之间的交叉、混合也同样重要------而且同样吸引人。在生物系统中一个令人惊讶的事实说明了这一点：为做到自复制而必需的所有因素（即：语言、程序、数据、解释程序及处理机）高度协同动作，以致于到了能把它们同时复制出来的程度------这表明生物的自复制比人类沿这些方向所设计的任何东西都深奥得多。例如本章开头所列出的自复制程序，就得假定事先存在三个客观的方面：语言，解释程序和处理机，而且并不复制它们。

我们来试着用计算机科学的术语总结一下给细胞的各个子单元分类的各种方法。先看DNA。由于DNA包含着构造作为该细胞的活性物质的各种蛋白质的全部信息，所以可以把DNA看成是用一种高层语言写出、随后又被翻译（或解释）成细胞``机器语言''的一个程序。另一方面，DNA本身又是受各种酶操纵的被动的分子，在这个意义上，DNA分子又恰像一长段数据。第三，DNA包含着能生成tRNA``单词卡片''的模板，这意味着DNA也含有它自己的高层语言的定义。

我们转向蛋白质。蛋白质是活性分子，并执行细胞的全部功能，所以，把它们看成用``细胞的机器语言''写成的程序（细胞本身做为处理机）就十分合适。另一方面，由于蛋白质是硬件，而大多数程序是软件，所以把蛋白质看成处理机也许更好一些。第三，蛋白质经常受到其它蛋白质的作用，这意味着蛋白质经常是数据。最后，还可以把蛋白质看作解释程序，这时是把DNA看成一组高级语言程序，在这种情形中，酶只是在执行DNA密码写下的程序，也就是说蛋白质起解释程序的作用。

还有核糖体及tRNA分子。它们是从DNA到蛋白质的翻译媒介，这种翻译可以比拟成一个程序从高级语言到机器语言的翻译过程。换句话说，核糖体有解释程序的功用，而tRNA分子规定了高级语言的定义。不过，对这个翻译过程还有一种看法：核糖体是处理机而tRNA是解释程序。

分析全部这些生物分子之间的相互关系时，我们还只是蜻蜓点水而已。所看到的是，我们倾向于认为有明显区别的那些层次，被大自然十分惬意的混合在了一起。实际上，在计算机科学中，已经有着把一个信息处理系统中看上去不同的所有方面混为一谈的明显趋势。人工智能的研究尤其如此，而人工智能的研究通常是站在计算机语言设计的最前沿。

生命的起源

了解了这些难以置信的、复杂的、相互联结的软件和硬件玩意儿，会使人提出一个自然而又基本的问题：``它们怎么开始的呢？''这倒真是一件令人困惑的事情。我们不得不设想有一个``揪鞋带举自己''的过程，这有点像开发新计算机语言时使用的那个过程------可是从简单分子到整个细胞的``揪鞋带''简直超出了我们的想象力。有许多关于生命起源的理论，它们遇到下述这个最最关键的问题时都不灵了：``遗传密码及其翻译机能（核糖体和tRNA分子）都是如何发生的？''目前，我们不得不满足于一种惊奇和敬畏的感觉，这不是一个答案带来的满足。也许体味这种惊奇和敬畏比有一个答案更令人满意------至少暂时是这样。

\begin{enumerate}
\item
  \phantomsection\label{Chapter16.xhtml_B_1}
\end{enumerate}

\phantomsection\label{Dialog17_split_000.xhtml}{}

\pandocbounded{\includegraphics[keepaspectratio]{./media/OEBPS/Images/img_104.jpg}}

图104．卡斯特罗瓦尔瓦

艾舍尔作（石板画，1930）

\phantomsection\label{Dialog17_split_001.xhtml}{}


\end{document}
